\RequirePackage{snapshot}
\documentclass[a4paper,oneside]{alpenspecs/alpenspecs}
% <<< Preamble
\hextrue
\papertrue
%https://tex.stackexchange.com/a/210456/131649
\renewcommand\partnumberlinebox[2]{#2\hspace{2em}}
\usetikzlibrary{positioning}
\usetikzlibrary{fit}
\usetikzlibrary{backgrounds}
\tcbuselibrary{breakable}
\tcbset{shield externalize}
\usepackage{alpenspecs/a3pages}
\makeindex
% >>>
\begin{document}
\begin{titlingpage}
    \fullhexpage{q1}{q0}
    \flushright\sffamily
    \enlargethispage{10ex}

    \vspace*{5em}
    
    \Huge\bfseries{Front-End Signal Processing for Red Pitaya Spectrum Analyzer}\\[1ex]
    \Large\mdseries{Bachelor Thesis: Specifications and Planning}\\[3ex]

    \normalsize\mdseries

    \tikzexternaldisable
%\tikzsetnextfilename{titlepage}
\begin{tikzpicture}[
    signalPath/.style={
        draw=q4,,
        fill=q4!50!white,
        circle,
        inner sep=0.3mm,
    },
    remember picture,
    overlay,
]
     \pgfplotsset{
         every axis/.style={
            anchor=center,
            at={(current page.center)},
            width=165mm,
            height=100mm,
            grid=none,
            axis line style={draw=none},
            tick style={draw=none},
            ticks = none,
            each nth point=4,
            filter discard warning=false,
            unbounded coords=discard,
            yshift=-6.75mm,
            xshift=16mm,
        }
    }
    \begin{axis}
        \addplot[ycomb,mark=*,mark options={scale=0.75},very thin,q1] 
            table[x=t,y=h,col sep=comma] {images/titlepage/impzChain25.csv};
    \end{axis}
\end{tikzpicture}
\tikzexternalenable


    \vfill
    \begin{tabular}{>{\bfseries}rl}
        Degree Program: & Electrical Engineering and Information Technology \\[2mm]
        Course:         & Project \num{6}\\[2mm]
        Coaches:        & Prof. Dr. Richard Gut, Michael Pichler \\[2mm]
        External Expert & Dr. J\"urg Stettbacher \\[2mm]
        Team:           & Raphael Frey, Noah H\"usser \\[2mm]
        Date:           & August 18, 2017 \\[2mm]
        Revision:       & 1.0.0 \\[2mm]
    \end{tabular}

    \tikzexternaldisable
        \begin{tikzpicture}[remember picture,overlay]
            \node[anchor=north east,yshift=-2.5mm,xshift=-2.5mm,] 
                at (current page.north east) {\includegraphics[width=10cm]{images/titlepage/logo-fhnw.pdf}};
            \node[anchor=south east,yshift=2.5mm,xshift=-2.5mm,] 
                at (current page.south east) {\includegraphics[width=4cm]{images/titlepage/logo-ime.pdf}};
        \end{tikzpicture}
    \tikzexternalenable
\end{titlingpage} % >>>

\frontmatter % <<<
\tableofcontents*
% >>>

\mainmatter

\chapter{Initial Situation} % <<< ------------------------------------------------ %
\label{ch:initial}
% ---------------------------------------------------------------------------- %

To replace  common lab tools which  are usually expensive and  big, Red Pitaya
produces a board  called STEMlab which features  a capable (125$\frac{MS}{s}$,
14bit) ADC and a ZYNQ SoC from Xilinx. The idea is to record data on the board
and transmit it  using the network infrastructure  to a PC which  does all the
number crunching.

To  reduce the  data rate  and processing  load, downsampling  filters can  be
enabled to decimate the signal. The filters are necessary in order to keep the
Nyquist theorem is fulfilled and to prevent aliasing. Sadly, the stock filters
are not  quite as capable  as desired, as they  do not attenuate  the stopband
well enough and have a frequency response not as flat as desired.

A  previous project  implemented  new filter  chains which  are  based on  IIR
filters. While those  filter chains worked  well in theory (i.e.  their design
was  fundamentally  sound),  they  did not  produce  satisfactory  results  in
practice.

This project  aims to  analyze the  previous work with  respect to  design and
implementation of those  filter chains, and to reproduce and  fix the problems
that exist in the earlier solution.

The   structure  of   the   idea   behind  the   STEMlab   can   be  seen   in
Figure~\ref{fig:system}.

\begin{figure}[h]
    \centering
    %\includegraphics[width=0.5\textwidth]{images/stemlab125.png}
    \includegraphics[width=0.60\textwidth]{images/system-overview.pdf}
    \caption{Block diagram of the STEMlab 125}
    \label{fig:system}
\end{figure}

\chapter{Objectives} % <<< ------------------------------------------------ %
\label{ch:objectives}
% ---------------------------------------------------------------------------- %

\emph{Note:} \emph{Revised}  goals are  the result  of decisions  and findings
during the project, which have lead to amendments in its goals.

\section{Must Haves}
    \begin{enumerate}
        \item Analysis of the previous project
        \begin{enumerate}
            \item What was the exact strategy?
            \item Why did it fail?
            \item How can potential mistakes be fixed?
            \item Is the project salvageable?
        \end{enumerate}
            \item Comparing and finding a solution to implement the desired filter chains
        \begin{enumerate}    
            \item Start over from the beginning, or
            \item Continue working on the existing solution.
        \end{enumerate}
            \item A complete specification for all the filter chains with respect to
        \begin{enumerate}
            \item Decimation Rates
            \item Attenuation in stopband
            \item Maximum ripple in passband
        \end{enumerate}
        \item A complete simulation of the filter chains either in Matlab or another tool.
        \item Filter implementation on the FPGA
        \item Verification of the performance of the implemented filters
    \end{enumerate}

\clearpage
\section{Must Haves (revised)}

After  analysis of  the previous  project and  the overall  STEMlab ecosystem,
these are the modified/amended objectives. For the full analysis, see the main
project report.

\begin{enumerate}
    \item Develop our own STEMlab toolchain including
    \begin{enumerate}
        \item Build system
        \item Good documentation
        \item Vivado TcL scripts
    \end{enumerate}
    \item Develop our own oscilloscope application
    \begin{enumerate}
        \item A server on the STEMlab
        \item A scope application in JavaScript
    \end{enumerate}
\end{enumerate}

\section{Nice to Haves}
\begin{enumerate}
    \item Additional scoping application features
    \begin{enumerate}
        \item Automatic SNR calculation
        \item Automatic THD calculation
        \item Multiple trace sources
        \item Different types of triggers
    \end{enumerate}
\end{enumerate}

\section{Do Not Implement (revised)}

\begin{enumerate}
    \item IIR filter chains
\end{enumerate}


\chapter{Objects of Delivery} % <<< ------------------------------------------------ %
\label{ch:objectives}
% ---------------------------------------------------------------------------- %

\begin{itemize}
    \item Specifications (periodically, 18.08.2017)
    \item Working Implementation of the Filter Chains (18.08.2017)
    \item Working Implementation of the JavaScript Scoping Application (18.08.2017)
    \item Factsheet (31.08.2017)
    \item Technical Report (18.08.2017), including Developer and User Guides
    \item Working Repository (containing  READMEs, Toolchain, Virtual Machine,
            Calculations, Simulations, Measurements, etc.) (18.08.2017)
    \item Presentation (15.09.2017)
\end{itemize}

\chapter{Concept (revised)} % <<< ------------------------------------------------ %
\label{ch:concept}
% ---------------------------------------------------------------------------- %

\section{Toolbox (revised)}

Since it has become  apparent that the STEMlab project base  led by Red Pitaya
is going through a major overhaul and it is uncertain whether this will change
in the future, we will develop our own toolbox.

We will implement a  toolchain that can build an ARM Linux  image with all the
necessary  components. This will  be done  based  on the  original Red  Pitaya
patches to  make the Linux  run on  the ZYNQ and  some scripts of  Open Source
projects which we will adapt to our needs and a more general toolchain.

We will use  GNU/Linux as a Platform  and Matlab and Vivado  for filter design
and implementation, respectively.

For  the  JavaScript  Application  we  will  use  node.js,  yarn,  rollup  and
mithril.js for the implementation.

\section{Realization (revised)}

The  preceding project  implemented filter  systems by  cascading several  IIR
filters.  While IIR filters tend to use fewer resources for a given transition
band width than FIR  filters, it still turned out that  resource usage was too
high to implement filter cascades for both channels present on the STEMlab.

Consequently, the  decision was made to  exploit CIC filters in  this project.
They are very computationally efficient and can be used as initial stages in a
filter cascade,  followed by FIR  filters to  shape the passband  response and
clean  up any  undesired  effects  which may  have  been  introduced into  the
frequency response of the system. Two types  of CIC filters are used: One with
a  decimation ratio  of  \num{25},  and another  with  a  decimation ratio  of
\num{125}.  Both are then followed by their respective compensation FIR filter
and more FIR filters for the final rate reduction.

As it has  been decided to develop an entirely  self-dependent solution of the
project, general components  on the FPGA such  as an ADC controller  or a data
logger are required. For the ADC part an  IP from Pavel Demin will be utilized
(see the project  report for further remarks and source). For  the logger part
an IP core that Mr. Hüsser developed in a previous project will be used.

For the ARM Core, a Linux will be  set up that can communicate with the logger
core via  a GNU/Linux kernel  module. A server application will  be developed
which will run on the Linux to send  the data written to the RAM by the logger
to the network. This will be done using WebSockets.

The    oscilloscope    web    application   will    be    implemented    using
JavaScript. WebSockets are used  to transmit and receive data to  and from the
network.  As build systems, yarn and rollup will be used. To implement the GUI
side of the application, mithril.js will be deployed.

\section{Testing}

To test the performance of the filters, a series of tests will be performed.

The following metrics will be calculated from the measured data.

\begin{itemize}
    \item The RMS frequency response
    \item The SNR frequency response
    \item The attenuation in the stopband
\end{itemize}

Optional testing would involve a comparison to the reference implementation and to a commercial spectrum analyzer.

\chapter{Project Timeline (revised)} % <<< ------------------------------------------------ %
\label{ch:initial}
% ---------------------------------------------------------------------------- %

\begin{a3pages}
    \begin{figure}
        \centering
        %% Creator: Matplotlib, PGF backend
%%
%% To include the figure in your LaTeX document, write
%%   \input{<filename>.pgf}
%%
%% Make sure the required packages are loaded in your preamble
%%   \usepackage{pgf}
%%
%% Figures using additional raster images can only be included by \input if
%% they are in the same directory as the main LaTeX file. For loading figures
%% from other directories you can use the `import` package
%%   \usepackage{import}
%% and then include the figures with
%%   \import{<path to file>}{<filename>.pgf}
%%
%% Matplotlib used the following preamble
%%   \usepackage{fontspec}
%%   \setmainfont{Times New Roman}
%%   \setsansfont{Lucida Grande}
%%   \setmonofont{Andale Mono}
%%
\begingroup%
\makeatletter%
\begin{pgfpicture}%
\pgfpathrectangle{\pgfpointorigin}{\pgfqpoint{15.000000in}{8.000000in}}%
\pgfusepath{use as bounding box, clip}%
\begin{pgfscope}%
\pgfsetbuttcap%
\pgfsetmiterjoin%
\definecolor{currentfill}{rgb}{1.000000,1.000000,1.000000}%
\pgfsetfillcolor{currentfill}%
\pgfsetlinewidth{0.000000pt}%
\definecolor{currentstroke}{rgb}{1.000000,1.000000,1.000000}%
\pgfsetstrokecolor{currentstroke}%
\pgfsetdash{}{0pt}%
\pgfpathmoveto{\pgfqpoint{0.000000in}{0.000000in}}%
\pgfpathlineto{\pgfqpoint{15.000000in}{0.000000in}}%
\pgfpathlineto{\pgfqpoint{15.000000in}{8.000000in}}%
\pgfpathlineto{\pgfqpoint{0.000000in}{8.000000in}}%
\pgfpathclose%
\pgfusepath{fill}%
\end{pgfscope}%
\begin{pgfscope}%
\pgfsetbuttcap%
\pgfsetmiterjoin%
\definecolor{currentfill}{rgb}{0.992157,0.964706,0.890196}%
\pgfsetfillcolor{currentfill}%
\pgfsetlinewidth{0.000000pt}%
\definecolor{currentstroke}{rgb}{0.000000,0.000000,0.000000}%
\pgfsetstrokecolor{currentstroke}%
\pgfsetstrokeopacity{0.000000}%
\pgfsetdash{}{0pt}%
\pgfpathmoveto{\pgfqpoint{4.009158in}{0.624448in}}%
\pgfpathlineto{\pgfqpoint{14.801389in}{0.624448in}}%
\pgfpathlineto{\pgfqpoint{14.801389in}{7.801389in}}%
\pgfpathlineto{\pgfqpoint{4.009158in}{7.801389in}}%
\pgfpathclose%
\pgfusepath{fill}%
\end{pgfscope}%
\begin{pgfscope}%
\pgfpathrectangle{\pgfqpoint{4.009158in}{0.624448in}}{\pgfqpoint{10.792231in}{7.176941in}} %
\pgfusepath{clip}%
\pgfsetrectcap%
\pgfsetroundjoin%
\pgfsetlinewidth{0.803000pt}%
\definecolor{currentstroke}{rgb}{1.000000,1.000000,1.000000}%
\pgfsetstrokecolor{currentstroke}%
\pgfsetdash{}{0pt}%
\pgfpathmoveto{\pgfqpoint{4.009158in}{0.624448in}}%
\pgfpathlineto{\pgfqpoint{4.009158in}{7.801389in}}%
\pgfusepath{stroke}%
\end{pgfscope}%
\begin{pgfscope}%
\pgfsetbuttcap%
\pgfsetroundjoin%
\definecolor{currentfill}{rgb}{0.333333,0.333333,0.333333}%
\pgfsetfillcolor{currentfill}%
\pgfsetlinewidth{0.803000pt}%
\definecolor{currentstroke}{rgb}{0.333333,0.333333,0.333333}%
\pgfsetstrokecolor{currentstroke}%
\pgfsetdash{}{0pt}%
\pgfsys@defobject{currentmarker}{\pgfqpoint{0.000000in}{-0.048611in}}{\pgfqpoint{0.000000in}{0.000000in}}{%
\pgfpathmoveto{\pgfqpoint{0.000000in}{0.000000in}}%
\pgfpathlineto{\pgfqpoint{0.000000in}{-0.048611in}}%
\pgfusepath{stroke,fill}%
}%
\begin{pgfscope}%
\pgfsys@transformshift{4.009158in}{0.624448in}%
\pgfsys@useobject{currentmarker}{}%
\end{pgfscope}%
\end{pgfscope}%
\begin{pgfscope}%
\definecolor{textcolor}{rgb}{0.333333,0.333333,0.333333}%
\pgfsetstrokecolor{textcolor}%
\pgfsetfillcolor{textcolor}%
\pgftext[x=3.554811in,y=0.179905in,left,base,rotate=30.000000]{\color{textcolor}\sffamily\fontsize{10.000000}{12.000000}\selectfont 01-Mar}%
\end{pgfscope}%
\begin{pgfscope}%
\pgfpathrectangle{\pgfqpoint{4.009158in}{0.624448in}}{\pgfqpoint{10.792231in}{7.176941in}} %
\pgfusepath{clip}%
\pgfsetrectcap%
\pgfsetroundjoin%
\pgfsetlinewidth{0.803000pt}%
\definecolor{currentstroke}{rgb}{1.000000,1.000000,1.000000}%
\pgfsetstrokecolor{currentstroke}%
\pgfsetdash{}{0pt}%
\pgfpathmoveto{\pgfqpoint{4.453544in}{0.624448in}}%
\pgfpathlineto{\pgfqpoint{4.453544in}{7.801389in}}%
\pgfusepath{stroke}%
\end{pgfscope}%
\begin{pgfscope}%
\pgfsetbuttcap%
\pgfsetroundjoin%
\definecolor{currentfill}{rgb}{0.333333,0.333333,0.333333}%
\pgfsetfillcolor{currentfill}%
\pgfsetlinewidth{0.803000pt}%
\definecolor{currentstroke}{rgb}{0.333333,0.333333,0.333333}%
\pgfsetstrokecolor{currentstroke}%
\pgfsetdash{}{0pt}%
\pgfsys@defobject{currentmarker}{\pgfqpoint{0.000000in}{-0.048611in}}{\pgfqpoint{0.000000in}{0.000000in}}{%
\pgfpathmoveto{\pgfqpoint{0.000000in}{0.000000in}}%
\pgfpathlineto{\pgfqpoint{0.000000in}{-0.048611in}}%
\pgfusepath{stroke,fill}%
}%
\begin{pgfscope}%
\pgfsys@transformshift{4.453544in}{0.624448in}%
\pgfsys@useobject{currentmarker}{}%
\end{pgfscope}%
\end{pgfscope}%
\begin{pgfscope}%
\definecolor{textcolor}{rgb}{0.333333,0.333333,0.333333}%
\pgfsetstrokecolor{textcolor}%
\pgfsetfillcolor{textcolor}%
\pgftext[x=3.999197in,y=0.179905in,left,base,rotate=30.000000]{\color{textcolor}\sffamily\fontsize{10.000000}{12.000000}\selectfont 08-Mar}%
\end{pgfscope}%
\begin{pgfscope}%
\pgfpathrectangle{\pgfqpoint{4.009158in}{0.624448in}}{\pgfqpoint{10.792231in}{7.176941in}} %
\pgfusepath{clip}%
\pgfsetrectcap%
\pgfsetroundjoin%
\pgfsetlinewidth{0.803000pt}%
\definecolor{currentstroke}{rgb}{1.000000,1.000000,1.000000}%
\pgfsetstrokecolor{currentstroke}%
\pgfsetdash{}{0pt}%
\pgfpathmoveto{\pgfqpoint{4.897930in}{0.624448in}}%
\pgfpathlineto{\pgfqpoint{4.897930in}{7.801389in}}%
\pgfusepath{stroke}%
\end{pgfscope}%
\begin{pgfscope}%
\pgfsetbuttcap%
\pgfsetroundjoin%
\definecolor{currentfill}{rgb}{0.333333,0.333333,0.333333}%
\pgfsetfillcolor{currentfill}%
\pgfsetlinewidth{0.803000pt}%
\definecolor{currentstroke}{rgb}{0.333333,0.333333,0.333333}%
\pgfsetstrokecolor{currentstroke}%
\pgfsetdash{}{0pt}%
\pgfsys@defobject{currentmarker}{\pgfqpoint{0.000000in}{-0.048611in}}{\pgfqpoint{0.000000in}{0.000000in}}{%
\pgfpathmoveto{\pgfqpoint{0.000000in}{0.000000in}}%
\pgfpathlineto{\pgfqpoint{0.000000in}{-0.048611in}}%
\pgfusepath{stroke,fill}%
}%
\begin{pgfscope}%
\pgfsys@transformshift{4.897930in}{0.624448in}%
\pgfsys@useobject{currentmarker}{}%
\end{pgfscope}%
\end{pgfscope}%
\begin{pgfscope}%
\definecolor{textcolor}{rgb}{0.333333,0.333333,0.333333}%
\pgfsetstrokecolor{textcolor}%
\pgfsetfillcolor{textcolor}%
\pgftext[x=4.443583in,y=0.179905in,left,base,rotate=30.000000]{\color{textcolor}\sffamily\fontsize{10.000000}{12.000000}\selectfont 15-Mar}%
\end{pgfscope}%
\begin{pgfscope}%
\pgfpathrectangle{\pgfqpoint{4.009158in}{0.624448in}}{\pgfqpoint{10.792231in}{7.176941in}} %
\pgfusepath{clip}%
\pgfsetrectcap%
\pgfsetroundjoin%
\pgfsetlinewidth{0.803000pt}%
\definecolor{currentstroke}{rgb}{1.000000,1.000000,1.000000}%
\pgfsetstrokecolor{currentstroke}%
\pgfsetdash{}{0pt}%
\pgfpathmoveto{\pgfqpoint{5.342316in}{0.624448in}}%
\pgfpathlineto{\pgfqpoint{5.342316in}{7.801389in}}%
\pgfusepath{stroke}%
\end{pgfscope}%
\begin{pgfscope}%
\pgfsetbuttcap%
\pgfsetroundjoin%
\definecolor{currentfill}{rgb}{0.333333,0.333333,0.333333}%
\pgfsetfillcolor{currentfill}%
\pgfsetlinewidth{0.803000pt}%
\definecolor{currentstroke}{rgb}{0.333333,0.333333,0.333333}%
\pgfsetstrokecolor{currentstroke}%
\pgfsetdash{}{0pt}%
\pgfsys@defobject{currentmarker}{\pgfqpoint{0.000000in}{-0.048611in}}{\pgfqpoint{0.000000in}{0.000000in}}{%
\pgfpathmoveto{\pgfqpoint{0.000000in}{0.000000in}}%
\pgfpathlineto{\pgfqpoint{0.000000in}{-0.048611in}}%
\pgfusepath{stroke,fill}%
}%
\begin{pgfscope}%
\pgfsys@transformshift{5.342316in}{0.624448in}%
\pgfsys@useobject{currentmarker}{}%
\end{pgfscope}%
\end{pgfscope}%
\begin{pgfscope}%
\definecolor{textcolor}{rgb}{0.333333,0.333333,0.333333}%
\pgfsetstrokecolor{textcolor}%
\pgfsetfillcolor{textcolor}%
\pgftext[x=4.887969in,y=0.179905in,left,base,rotate=30.000000]{\color{textcolor}\sffamily\fontsize{10.000000}{12.000000}\selectfont 22-Mar}%
\end{pgfscope}%
\begin{pgfscope}%
\pgfpathrectangle{\pgfqpoint{4.009158in}{0.624448in}}{\pgfqpoint{10.792231in}{7.176941in}} %
\pgfusepath{clip}%
\pgfsetrectcap%
\pgfsetroundjoin%
\pgfsetlinewidth{0.803000pt}%
\definecolor{currentstroke}{rgb}{1.000000,1.000000,1.000000}%
\pgfsetstrokecolor{currentstroke}%
\pgfsetdash{}{0pt}%
\pgfpathmoveto{\pgfqpoint{5.786702in}{0.624448in}}%
\pgfpathlineto{\pgfqpoint{5.786702in}{7.801389in}}%
\pgfusepath{stroke}%
\end{pgfscope}%
\begin{pgfscope}%
\pgfsetbuttcap%
\pgfsetroundjoin%
\definecolor{currentfill}{rgb}{0.333333,0.333333,0.333333}%
\pgfsetfillcolor{currentfill}%
\pgfsetlinewidth{0.803000pt}%
\definecolor{currentstroke}{rgb}{0.333333,0.333333,0.333333}%
\pgfsetstrokecolor{currentstroke}%
\pgfsetdash{}{0pt}%
\pgfsys@defobject{currentmarker}{\pgfqpoint{0.000000in}{-0.048611in}}{\pgfqpoint{0.000000in}{0.000000in}}{%
\pgfpathmoveto{\pgfqpoint{0.000000in}{0.000000in}}%
\pgfpathlineto{\pgfqpoint{0.000000in}{-0.048611in}}%
\pgfusepath{stroke,fill}%
}%
\begin{pgfscope}%
\pgfsys@transformshift{5.786702in}{0.624448in}%
\pgfsys@useobject{currentmarker}{}%
\end{pgfscope}%
\end{pgfscope}%
\begin{pgfscope}%
\definecolor{textcolor}{rgb}{0.333333,0.333333,0.333333}%
\pgfsetstrokecolor{textcolor}%
\pgfsetfillcolor{textcolor}%
\pgftext[x=5.332355in,y=0.179905in,left,base,rotate=30.000000]{\color{textcolor}\sffamily\fontsize{10.000000}{12.000000}\selectfont 29-Mar}%
\end{pgfscope}%
\begin{pgfscope}%
\pgfpathrectangle{\pgfqpoint{4.009158in}{0.624448in}}{\pgfqpoint{10.792231in}{7.176941in}} %
\pgfusepath{clip}%
\pgfsetrectcap%
\pgfsetroundjoin%
\pgfsetlinewidth{0.803000pt}%
\definecolor{currentstroke}{rgb}{1.000000,1.000000,1.000000}%
\pgfsetstrokecolor{currentstroke}%
\pgfsetdash{}{0pt}%
\pgfpathmoveto{\pgfqpoint{6.231088in}{0.624448in}}%
\pgfpathlineto{\pgfqpoint{6.231088in}{7.801389in}}%
\pgfusepath{stroke}%
\end{pgfscope}%
\begin{pgfscope}%
\pgfsetbuttcap%
\pgfsetroundjoin%
\definecolor{currentfill}{rgb}{0.333333,0.333333,0.333333}%
\pgfsetfillcolor{currentfill}%
\pgfsetlinewidth{0.803000pt}%
\definecolor{currentstroke}{rgb}{0.333333,0.333333,0.333333}%
\pgfsetstrokecolor{currentstroke}%
\pgfsetdash{}{0pt}%
\pgfsys@defobject{currentmarker}{\pgfqpoint{0.000000in}{-0.048611in}}{\pgfqpoint{0.000000in}{0.000000in}}{%
\pgfpathmoveto{\pgfqpoint{0.000000in}{0.000000in}}%
\pgfpathlineto{\pgfqpoint{0.000000in}{-0.048611in}}%
\pgfusepath{stroke,fill}%
}%
\begin{pgfscope}%
\pgfsys@transformshift{6.231088in}{0.624448in}%
\pgfsys@useobject{currentmarker}{}%
\end{pgfscope}%
\end{pgfscope}%
\begin{pgfscope}%
\definecolor{textcolor}{rgb}{0.333333,0.333333,0.333333}%
\pgfsetstrokecolor{textcolor}%
\pgfsetfillcolor{textcolor}%
\pgftext[x=5.788076in,y=0.186449in,left,base,rotate=30.000000]{\color{textcolor}\sffamily\fontsize{10.000000}{12.000000}\selectfont 05-Apr}%
\end{pgfscope}%
\begin{pgfscope}%
\pgfpathrectangle{\pgfqpoint{4.009158in}{0.624448in}}{\pgfqpoint{10.792231in}{7.176941in}} %
\pgfusepath{clip}%
\pgfsetrectcap%
\pgfsetroundjoin%
\pgfsetlinewidth{0.803000pt}%
\definecolor{currentstroke}{rgb}{1.000000,1.000000,1.000000}%
\pgfsetstrokecolor{currentstroke}%
\pgfsetdash{}{0pt}%
\pgfpathmoveto{\pgfqpoint{6.675474in}{0.624448in}}%
\pgfpathlineto{\pgfqpoint{6.675474in}{7.801389in}}%
\pgfusepath{stroke}%
\end{pgfscope}%
\begin{pgfscope}%
\pgfsetbuttcap%
\pgfsetroundjoin%
\definecolor{currentfill}{rgb}{0.333333,0.333333,0.333333}%
\pgfsetfillcolor{currentfill}%
\pgfsetlinewidth{0.803000pt}%
\definecolor{currentstroke}{rgb}{0.333333,0.333333,0.333333}%
\pgfsetstrokecolor{currentstroke}%
\pgfsetdash{}{0pt}%
\pgfsys@defobject{currentmarker}{\pgfqpoint{0.000000in}{-0.048611in}}{\pgfqpoint{0.000000in}{0.000000in}}{%
\pgfpathmoveto{\pgfqpoint{0.000000in}{0.000000in}}%
\pgfpathlineto{\pgfqpoint{0.000000in}{-0.048611in}}%
\pgfusepath{stroke,fill}%
}%
\begin{pgfscope}%
\pgfsys@transformshift{6.675474in}{0.624448in}%
\pgfsys@useobject{currentmarker}{}%
\end{pgfscope}%
\end{pgfscope}%
\begin{pgfscope}%
\definecolor{textcolor}{rgb}{0.333333,0.333333,0.333333}%
\pgfsetstrokecolor{textcolor}%
\pgfsetfillcolor{textcolor}%
\pgftext[x=6.232462in,y=0.186449in,left,base,rotate=30.000000]{\color{textcolor}\sffamily\fontsize{10.000000}{12.000000}\selectfont 12-Apr}%
\end{pgfscope}%
\begin{pgfscope}%
\pgfpathrectangle{\pgfqpoint{4.009158in}{0.624448in}}{\pgfqpoint{10.792231in}{7.176941in}} %
\pgfusepath{clip}%
\pgfsetrectcap%
\pgfsetroundjoin%
\pgfsetlinewidth{0.803000pt}%
\definecolor{currentstroke}{rgb}{1.000000,1.000000,1.000000}%
\pgfsetstrokecolor{currentstroke}%
\pgfsetdash{}{0pt}%
\pgfpathmoveto{\pgfqpoint{7.119860in}{0.624448in}}%
\pgfpathlineto{\pgfqpoint{7.119860in}{7.801389in}}%
\pgfusepath{stroke}%
\end{pgfscope}%
\begin{pgfscope}%
\pgfsetbuttcap%
\pgfsetroundjoin%
\definecolor{currentfill}{rgb}{0.333333,0.333333,0.333333}%
\pgfsetfillcolor{currentfill}%
\pgfsetlinewidth{0.803000pt}%
\definecolor{currentstroke}{rgb}{0.333333,0.333333,0.333333}%
\pgfsetstrokecolor{currentstroke}%
\pgfsetdash{}{0pt}%
\pgfsys@defobject{currentmarker}{\pgfqpoint{0.000000in}{-0.048611in}}{\pgfqpoint{0.000000in}{0.000000in}}{%
\pgfpathmoveto{\pgfqpoint{0.000000in}{0.000000in}}%
\pgfpathlineto{\pgfqpoint{0.000000in}{-0.048611in}}%
\pgfusepath{stroke,fill}%
}%
\begin{pgfscope}%
\pgfsys@transformshift{7.119860in}{0.624448in}%
\pgfsys@useobject{currentmarker}{}%
\end{pgfscope}%
\end{pgfscope}%
\begin{pgfscope}%
\definecolor{textcolor}{rgb}{0.333333,0.333333,0.333333}%
\pgfsetstrokecolor{textcolor}%
\pgfsetfillcolor{textcolor}%
\pgftext[x=6.676848in,y=0.186449in,left,base,rotate=30.000000]{\color{textcolor}\sffamily\fontsize{10.000000}{12.000000}\selectfont 19-Apr}%
\end{pgfscope}%
\begin{pgfscope}%
\pgfpathrectangle{\pgfqpoint{4.009158in}{0.624448in}}{\pgfqpoint{10.792231in}{7.176941in}} %
\pgfusepath{clip}%
\pgfsetrectcap%
\pgfsetroundjoin%
\pgfsetlinewidth{0.803000pt}%
\definecolor{currentstroke}{rgb}{1.000000,1.000000,1.000000}%
\pgfsetstrokecolor{currentstroke}%
\pgfsetdash{}{0pt}%
\pgfpathmoveto{\pgfqpoint{7.564246in}{0.624448in}}%
\pgfpathlineto{\pgfqpoint{7.564246in}{7.801389in}}%
\pgfusepath{stroke}%
\end{pgfscope}%
\begin{pgfscope}%
\pgfsetbuttcap%
\pgfsetroundjoin%
\definecolor{currentfill}{rgb}{0.333333,0.333333,0.333333}%
\pgfsetfillcolor{currentfill}%
\pgfsetlinewidth{0.803000pt}%
\definecolor{currentstroke}{rgb}{0.333333,0.333333,0.333333}%
\pgfsetstrokecolor{currentstroke}%
\pgfsetdash{}{0pt}%
\pgfsys@defobject{currentmarker}{\pgfqpoint{0.000000in}{-0.048611in}}{\pgfqpoint{0.000000in}{0.000000in}}{%
\pgfpathmoveto{\pgfqpoint{0.000000in}{0.000000in}}%
\pgfpathlineto{\pgfqpoint{0.000000in}{-0.048611in}}%
\pgfusepath{stroke,fill}%
}%
\begin{pgfscope}%
\pgfsys@transformshift{7.564246in}{0.624448in}%
\pgfsys@useobject{currentmarker}{}%
\end{pgfscope}%
\end{pgfscope}%
\begin{pgfscope}%
\definecolor{textcolor}{rgb}{0.333333,0.333333,0.333333}%
\pgfsetstrokecolor{textcolor}%
\pgfsetfillcolor{textcolor}%
\pgftext[x=7.121234in,y=0.186449in,left,base,rotate=30.000000]{\color{textcolor}\sffamily\fontsize{10.000000}{12.000000}\selectfont 26-Apr}%
\end{pgfscope}%
\begin{pgfscope}%
\pgfpathrectangle{\pgfqpoint{4.009158in}{0.624448in}}{\pgfqpoint{10.792231in}{7.176941in}} %
\pgfusepath{clip}%
\pgfsetrectcap%
\pgfsetroundjoin%
\pgfsetlinewidth{0.803000pt}%
\definecolor{currentstroke}{rgb}{1.000000,1.000000,1.000000}%
\pgfsetstrokecolor{currentstroke}%
\pgfsetdash{}{0pt}%
\pgfpathmoveto{\pgfqpoint{8.008632in}{0.624448in}}%
\pgfpathlineto{\pgfqpoint{8.008632in}{7.801389in}}%
\pgfusepath{stroke}%
\end{pgfscope}%
\begin{pgfscope}%
\pgfsetbuttcap%
\pgfsetroundjoin%
\definecolor{currentfill}{rgb}{0.333333,0.333333,0.333333}%
\pgfsetfillcolor{currentfill}%
\pgfsetlinewidth{0.803000pt}%
\definecolor{currentstroke}{rgb}{0.333333,0.333333,0.333333}%
\pgfsetstrokecolor{currentstroke}%
\pgfsetdash{}{0pt}%
\pgfsys@defobject{currentmarker}{\pgfqpoint{0.000000in}{-0.048611in}}{\pgfqpoint{0.000000in}{0.000000in}}{%
\pgfpathmoveto{\pgfqpoint{0.000000in}{0.000000in}}%
\pgfpathlineto{\pgfqpoint{0.000000in}{-0.048611in}}%
\pgfusepath{stroke,fill}%
}%
\begin{pgfscope}%
\pgfsys@transformshift{8.008632in}{0.624448in}%
\pgfsys@useobject{currentmarker}{}%
\end{pgfscope}%
\end{pgfscope}%
\begin{pgfscope}%
\definecolor{textcolor}{rgb}{0.333333,0.333333,0.333333}%
\pgfsetstrokecolor{textcolor}%
\pgfsetfillcolor{textcolor}%
\pgftext[x=7.540659in,y=0.172038in,left,base,rotate=30.000000]{\color{textcolor}\sffamily\fontsize{10.000000}{12.000000}\selectfont 03-May}%
\end{pgfscope}%
\begin{pgfscope}%
\pgfpathrectangle{\pgfqpoint{4.009158in}{0.624448in}}{\pgfqpoint{10.792231in}{7.176941in}} %
\pgfusepath{clip}%
\pgfsetrectcap%
\pgfsetroundjoin%
\pgfsetlinewidth{0.803000pt}%
\definecolor{currentstroke}{rgb}{1.000000,1.000000,1.000000}%
\pgfsetstrokecolor{currentstroke}%
\pgfsetdash{}{0pt}%
\pgfpathmoveto{\pgfqpoint{8.453018in}{0.624448in}}%
\pgfpathlineto{\pgfqpoint{8.453018in}{7.801389in}}%
\pgfusepath{stroke}%
\end{pgfscope}%
\begin{pgfscope}%
\pgfsetbuttcap%
\pgfsetroundjoin%
\definecolor{currentfill}{rgb}{0.333333,0.333333,0.333333}%
\pgfsetfillcolor{currentfill}%
\pgfsetlinewidth{0.803000pt}%
\definecolor{currentstroke}{rgb}{0.333333,0.333333,0.333333}%
\pgfsetstrokecolor{currentstroke}%
\pgfsetdash{}{0pt}%
\pgfsys@defobject{currentmarker}{\pgfqpoint{0.000000in}{-0.048611in}}{\pgfqpoint{0.000000in}{0.000000in}}{%
\pgfpathmoveto{\pgfqpoint{0.000000in}{0.000000in}}%
\pgfpathlineto{\pgfqpoint{0.000000in}{-0.048611in}}%
\pgfusepath{stroke,fill}%
}%
\begin{pgfscope}%
\pgfsys@transformshift{8.453018in}{0.624448in}%
\pgfsys@useobject{currentmarker}{}%
\end{pgfscope}%
\end{pgfscope}%
\begin{pgfscope}%
\definecolor{textcolor}{rgb}{0.333333,0.333333,0.333333}%
\pgfsetstrokecolor{textcolor}%
\pgfsetfillcolor{textcolor}%
\pgftext[x=7.985045in,y=0.172038in,left,base,rotate=30.000000]{\color{textcolor}\sffamily\fontsize{10.000000}{12.000000}\selectfont 10-May}%
\end{pgfscope}%
\begin{pgfscope}%
\pgfpathrectangle{\pgfqpoint{4.009158in}{0.624448in}}{\pgfqpoint{10.792231in}{7.176941in}} %
\pgfusepath{clip}%
\pgfsetrectcap%
\pgfsetroundjoin%
\pgfsetlinewidth{0.803000pt}%
\definecolor{currentstroke}{rgb}{1.000000,1.000000,1.000000}%
\pgfsetstrokecolor{currentstroke}%
\pgfsetdash{}{0pt}%
\pgfpathmoveto{\pgfqpoint{8.897404in}{0.624448in}}%
\pgfpathlineto{\pgfqpoint{8.897404in}{7.801389in}}%
\pgfusepath{stroke}%
\end{pgfscope}%
\begin{pgfscope}%
\pgfsetbuttcap%
\pgfsetroundjoin%
\definecolor{currentfill}{rgb}{0.333333,0.333333,0.333333}%
\pgfsetfillcolor{currentfill}%
\pgfsetlinewidth{0.803000pt}%
\definecolor{currentstroke}{rgb}{0.333333,0.333333,0.333333}%
\pgfsetstrokecolor{currentstroke}%
\pgfsetdash{}{0pt}%
\pgfsys@defobject{currentmarker}{\pgfqpoint{0.000000in}{-0.048611in}}{\pgfqpoint{0.000000in}{0.000000in}}{%
\pgfpathmoveto{\pgfqpoint{0.000000in}{0.000000in}}%
\pgfpathlineto{\pgfqpoint{0.000000in}{-0.048611in}}%
\pgfusepath{stroke,fill}%
}%
\begin{pgfscope}%
\pgfsys@transformshift{8.897404in}{0.624448in}%
\pgfsys@useobject{currentmarker}{}%
\end{pgfscope}%
\end{pgfscope}%
\begin{pgfscope}%
\definecolor{textcolor}{rgb}{0.333333,0.333333,0.333333}%
\pgfsetstrokecolor{textcolor}%
\pgfsetfillcolor{textcolor}%
\pgftext[x=8.429431in,y=0.172038in,left,base,rotate=30.000000]{\color{textcolor}\sffamily\fontsize{10.000000}{12.000000}\selectfont 17-May}%
\end{pgfscope}%
\begin{pgfscope}%
\pgfpathrectangle{\pgfqpoint{4.009158in}{0.624448in}}{\pgfqpoint{10.792231in}{7.176941in}} %
\pgfusepath{clip}%
\pgfsetrectcap%
\pgfsetroundjoin%
\pgfsetlinewidth{0.803000pt}%
\definecolor{currentstroke}{rgb}{1.000000,1.000000,1.000000}%
\pgfsetstrokecolor{currentstroke}%
\pgfsetdash{}{0pt}%
\pgfpathmoveto{\pgfqpoint{9.341790in}{0.624448in}}%
\pgfpathlineto{\pgfqpoint{9.341790in}{7.801389in}}%
\pgfusepath{stroke}%
\end{pgfscope}%
\begin{pgfscope}%
\pgfsetbuttcap%
\pgfsetroundjoin%
\definecolor{currentfill}{rgb}{0.333333,0.333333,0.333333}%
\pgfsetfillcolor{currentfill}%
\pgfsetlinewidth{0.803000pt}%
\definecolor{currentstroke}{rgb}{0.333333,0.333333,0.333333}%
\pgfsetstrokecolor{currentstroke}%
\pgfsetdash{}{0pt}%
\pgfsys@defobject{currentmarker}{\pgfqpoint{0.000000in}{-0.048611in}}{\pgfqpoint{0.000000in}{0.000000in}}{%
\pgfpathmoveto{\pgfqpoint{0.000000in}{0.000000in}}%
\pgfpathlineto{\pgfqpoint{0.000000in}{-0.048611in}}%
\pgfusepath{stroke,fill}%
}%
\begin{pgfscope}%
\pgfsys@transformshift{9.341790in}{0.624448in}%
\pgfsys@useobject{currentmarker}{}%
\end{pgfscope}%
\end{pgfscope}%
\begin{pgfscope}%
\definecolor{textcolor}{rgb}{0.333333,0.333333,0.333333}%
\pgfsetstrokecolor{textcolor}%
\pgfsetfillcolor{textcolor}%
\pgftext[x=8.873817in,y=0.172038in,left,base,rotate=30.000000]{\color{textcolor}\sffamily\fontsize{10.000000}{12.000000}\selectfont 24-May}%
\end{pgfscope}%
\begin{pgfscope}%
\pgfpathrectangle{\pgfqpoint{4.009158in}{0.624448in}}{\pgfqpoint{10.792231in}{7.176941in}} %
\pgfusepath{clip}%
\pgfsetrectcap%
\pgfsetroundjoin%
\pgfsetlinewidth{0.803000pt}%
\definecolor{currentstroke}{rgb}{1.000000,1.000000,1.000000}%
\pgfsetstrokecolor{currentstroke}%
\pgfsetdash{}{0pt}%
\pgfpathmoveto{\pgfqpoint{9.786176in}{0.624448in}}%
\pgfpathlineto{\pgfqpoint{9.786176in}{7.801389in}}%
\pgfusepath{stroke}%
\end{pgfscope}%
\begin{pgfscope}%
\pgfsetbuttcap%
\pgfsetroundjoin%
\definecolor{currentfill}{rgb}{0.333333,0.333333,0.333333}%
\pgfsetfillcolor{currentfill}%
\pgfsetlinewidth{0.803000pt}%
\definecolor{currentstroke}{rgb}{0.333333,0.333333,0.333333}%
\pgfsetstrokecolor{currentstroke}%
\pgfsetdash{}{0pt}%
\pgfsys@defobject{currentmarker}{\pgfqpoint{0.000000in}{-0.048611in}}{\pgfqpoint{0.000000in}{0.000000in}}{%
\pgfpathmoveto{\pgfqpoint{0.000000in}{0.000000in}}%
\pgfpathlineto{\pgfqpoint{0.000000in}{-0.048611in}}%
\pgfusepath{stroke,fill}%
}%
\begin{pgfscope}%
\pgfsys@transformshift{9.786176in}{0.624448in}%
\pgfsys@useobject{currentmarker}{}%
\end{pgfscope}%
\end{pgfscope}%
\begin{pgfscope}%
\definecolor{textcolor}{rgb}{0.333333,0.333333,0.333333}%
\pgfsetstrokecolor{textcolor}%
\pgfsetfillcolor{textcolor}%
\pgftext[x=9.318203in,y=0.172038in,left,base,rotate=30.000000]{\color{textcolor}\sffamily\fontsize{10.000000}{12.000000}\selectfont 31-May}%
\end{pgfscope}%
\begin{pgfscope}%
\pgfpathrectangle{\pgfqpoint{4.009158in}{0.624448in}}{\pgfqpoint{10.792231in}{7.176941in}} %
\pgfusepath{clip}%
\pgfsetrectcap%
\pgfsetroundjoin%
\pgfsetlinewidth{0.803000pt}%
\definecolor{currentstroke}{rgb}{1.000000,1.000000,1.000000}%
\pgfsetstrokecolor{currentstroke}%
\pgfsetdash{}{0pt}%
\pgfpathmoveto{\pgfqpoint{10.230562in}{0.624448in}}%
\pgfpathlineto{\pgfqpoint{10.230562in}{7.801389in}}%
\pgfusepath{stroke}%
\end{pgfscope}%
\begin{pgfscope}%
\pgfsetbuttcap%
\pgfsetroundjoin%
\definecolor{currentfill}{rgb}{0.333333,0.333333,0.333333}%
\pgfsetfillcolor{currentfill}%
\pgfsetlinewidth{0.803000pt}%
\definecolor{currentstroke}{rgb}{0.333333,0.333333,0.333333}%
\pgfsetstrokecolor{currentstroke}%
\pgfsetdash{}{0pt}%
\pgfsys@defobject{currentmarker}{\pgfqpoint{0.000000in}{-0.048611in}}{\pgfqpoint{0.000000in}{0.000000in}}{%
\pgfpathmoveto{\pgfqpoint{0.000000in}{0.000000in}}%
\pgfpathlineto{\pgfqpoint{0.000000in}{-0.048611in}}%
\pgfusepath{stroke,fill}%
}%
\begin{pgfscope}%
\pgfsys@transformshift{10.230562in}{0.624448in}%
\pgfsys@useobject{currentmarker}{}%
\end{pgfscope}%
\end{pgfscope}%
\begin{pgfscope}%
\definecolor{textcolor}{rgb}{0.333333,0.333333,0.333333}%
\pgfsetstrokecolor{textcolor}%
\pgfsetfillcolor{textcolor}%
\pgftext[x=9.808752in,y=0.198690in,left,base,rotate=30.000000]{\color{textcolor}\sffamily\fontsize{10.000000}{12.000000}\selectfont 07-Jun}%
\end{pgfscope}%
\begin{pgfscope}%
\pgfpathrectangle{\pgfqpoint{4.009158in}{0.624448in}}{\pgfqpoint{10.792231in}{7.176941in}} %
\pgfusepath{clip}%
\pgfsetrectcap%
\pgfsetroundjoin%
\pgfsetlinewidth{0.803000pt}%
\definecolor{currentstroke}{rgb}{1.000000,1.000000,1.000000}%
\pgfsetstrokecolor{currentstroke}%
\pgfsetdash{}{0pt}%
\pgfpathmoveto{\pgfqpoint{10.674948in}{0.624448in}}%
\pgfpathlineto{\pgfqpoint{10.674948in}{7.801389in}}%
\pgfusepath{stroke}%
\end{pgfscope}%
\begin{pgfscope}%
\pgfsetbuttcap%
\pgfsetroundjoin%
\definecolor{currentfill}{rgb}{0.333333,0.333333,0.333333}%
\pgfsetfillcolor{currentfill}%
\pgfsetlinewidth{0.803000pt}%
\definecolor{currentstroke}{rgb}{0.333333,0.333333,0.333333}%
\pgfsetstrokecolor{currentstroke}%
\pgfsetdash{}{0pt}%
\pgfsys@defobject{currentmarker}{\pgfqpoint{0.000000in}{-0.048611in}}{\pgfqpoint{0.000000in}{0.000000in}}{%
\pgfpathmoveto{\pgfqpoint{0.000000in}{0.000000in}}%
\pgfpathlineto{\pgfqpoint{0.000000in}{-0.048611in}}%
\pgfusepath{stroke,fill}%
}%
\begin{pgfscope}%
\pgfsys@transformshift{10.674948in}{0.624448in}%
\pgfsys@useobject{currentmarker}{}%
\end{pgfscope}%
\end{pgfscope}%
\begin{pgfscope}%
\definecolor{textcolor}{rgb}{0.333333,0.333333,0.333333}%
\pgfsetstrokecolor{textcolor}%
\pgfsetfillcolor{textcolor}%
\pgftext[x=10.253138in,y=0.198690in,left,base,rotate=30.000000]{\color{textcolor}\sffamily\fontsize{10.000000}{12.000000}\selectfont 14-Jun}%
\end{pgfscope}%
\begin{pgfscope}%
\pgfpathrectangle{\pgfqpoint{4.009158in}{0.624448in}}{\pgfqpoint{10.792231in}{7.176941in}} %
\pgfusepath{clip}%
\pgfsetrectcap%
\pgfsetroundjoin%
\pgfsetlinewidth{0.803000pt}%
\definecolor{currentstroke}{rgb}{1.000000,1.000000,1.000000}%
\pgfsetstrokecolor{currentstroke}%
\pgfsetdash{}{0pt}%
\pgfpathmoveto{\pgfqpoint{11.119334in}{0.624448in}}%
\pgfpathlineto{\pgfqpoint{11.119334in}{7.801389in}}%
\pgfusepath{stroke}%
\end{pgfscope}%
\begin{pgfscope}%
\pgfsetbuttcap%
\pgfsetroundjoin%
\definecolor{currentfill}{rgb}{0.333333,0.333333,0.333333}%
\pgfsetfillcolor{currentfill}%
\pgfsetlinewidth{0.803000pt}%
\definecolor{currentstroke}{rgb}{0.333333,0.333333,0.333333}%
\pgfsetstrokecolor{currentstroke}%
\pgfsetdash{}{0pt}%
\pgfsys@defobject{currentmarker}{\pgfqpoint{0.000000in}{-0.048611in}}{\pgfqpoint{0.000000in}{0.000000in}}{%
\pgfpathmoveto{\pgfqpoint{0.000000in}{0.000000in}}%
\pgfpathlineto{\pgfqpoint{0.000000in}{-0.048611in}}%
\pgfusepath{stroke,fill}%
}%
\begin{pgfscope}%
\pgfsys@transformshift{11.119334in}{0.624448in}%
\pgfsys@useobject{currentmarker}{}%
\end{pgfscope}%
\end{pgfscope}%
\begin{pgfscope}%
\definecolor{textcolor}{rgb}{0.333333,0.333333,0.333333}%
\pgfsetstrokecolor{textcolor}%
\pgfsetfillcolor{textcolor}%
\pgftext[x=10.697524in,y=0.198690in,left,base,rotate=30.000000]{\color{textcolor}\sffamily\fontsize{10.000000}{12.000000}\selectfont 21-Jun}%
\end{pgfscope}%
\begin{pgfscope}%
\pgfpathrectangle{\pgfqpoint{4.009158in}{0.624448in}}{\pgfqpoint{10.792231in}{7.176941in}} %
\pgfusepath{clip}%
\pgfsetrectcap%
\pgfsetroundjoin%
\pgfsetlinewidth{0.803000pt}%
\definecolor{currentstroke}{rgb}{1.000000,1.000000,1.000000}%
\pgfsetstrokecolor{currentstroke}%
\pgfsetdash{}{0pt}%
\pgfpathmoveto{\pgfqpoint{11.563720in}{0.624448in}}%
\pgfpathlineto{\pgfqpoint{11.563720in}{7.801389in}}%
\pgfusepath{stroke}%
\end{pgfscope}%
\begin{pgfscope}%
\pgfsetbuttcap%
\pgfsetroundjoin%
\definecolor{currentfill}{rgb}{0.333333,0.333333,0.333333}%
\pgfsetfillcolor{currentfill}%
\pgfsetlinewidth{0.803000pt}%
\definecolor{currentstroke}{rgb}{0.333333,0.333333,0.333333}%
\pgfsetstrokecolor{currentstroke}%
\pgfsetdash{}{0pt}%
\pgfsys@defobject{currentmarker}{\pgfqpoint{0.000000in}{-0.048611in}}{\pgfqpoint{0.000000in}{0.000000in}}{%
\pgfpathmoveto{\pgfqpoint{0.000000in}{0.000000in}}%
\pgfpathlineto{\pgfqpoint{0.000000in}{-0.048611in}}%
\pgfusepath{stroke,fill}%
}%
\begin{pgfscope}%
\pgfsys@transformshift{11.563720in}{0.624448in}%
\pgfsys@useobject{currentmarker}{}%
\end{pgfscope}%
\end{pgfscope}%
\begin{pgfscope}%
\definecolor{textcolor}{rgb}{0.333333,0.333333,0.333333}%
\pgfsetstrokecolor{textcolor}%
\pgfsetfillcolor{textcolor}%
\pgftext[x=11.141910in,y=0.198690in,left,base,rotate=30.000000]{\color{textcolor}\sffamily\fontsize{10.000000}{12.000000}\selectfont 28-Jun}%
\end{pgfscope}%
\begin{pgfscope}%
\pgfpathrectangle{\pgfqpoint{4.009158in}{0.624448in}}{\pgfqpoint{10.792231in}{7.176941in}} %
\pgfusepath{clip}%
\pgfsetrectcap%
\pgfsetroundjoin%
\pgfsetlinewidth{0.803000pt}%
\definecolor{currentstroke}{rgb}{1.000000,1.000000,1.000000}%
\pgfsetstrokecolor{currentstroke}%
\pgfsetdash{}{0pt}%
\pgfpathmoveto{\pgfqpoint{12.008106in}{0.624448in}}%
\pgfpathlineto{\pgfqpoint{12.008106in}{7.801389in}}%
\pgfusepath{stroke}%
\end{pgfscope}%
\begin{pgfscope}%
\pgfsetbuttcap%
\pgfsetroundjoin%
\definecolor{currentfill}{rgb}{0.333333,0.333333,0.333333}%
\pgfsetfillcolor{currentfill}%
\pgfsetlinewidth{0.803000pt}%
\definecolor{currentstroke}{rgb}{0.333333,0.333333,0.333333}%
\pgfsetstrokecolor{currentstroke}%
\pgfsetdash{}{0pt}%
\pgfsys@defobject{currentmarker}{\pgfqpoint{0.000000in}{-0.048611in}}{\pgfqpoint{0.000000in}{0.000000in}}{%
\pgfpathmoveto{\pgfqpoint{0.000000in}{0.000000in}}%
\pgfpathlineto{\pgfqpoint{0.000000in}{-0.048611in}}%
\pgfusepath{stroke,fill}%
}%
\begin{pgfscope}%
\pgfsys@transformshift{12.008106in}{0.624448in}%
\pgfsys@useobject{currentmarker}{}%
\end{pgfscope}%
\end{pgfscope}%
\begin{pgfscope}%
\definecolor{textcolor}{rgb}{0.333333,0.333333,0.333333}%
\pgfsetstrokecolor{textcolor}%
\pgfsetfillcolor{textcolor}%
\pgftext[x=11.626174in,y=0.221714in,left,base,rotate=30.000000]{\color{textcolor}\sffamily\fontsize{10.000000}{12.000000}\selectfont 05-Jul}%
\end{pgfscope}%
\begin{pgfscope}%
\pgfpathrectangle{\pgfqpoint{4.009158in}{0.624448in}}{\pgfqpoint{10.792231in}{7.176941in}} %
\pgfusepath{clip}%
\pgfsetrectcap%
\pgfsetroundjoin%
\pgfsetlinewidth{0.803000pt}%
\definecolor{currentstroke}{rgb}{1.000000,1.000000,1.000000}%
\pgfsetstrokecolor{currentstroke}%
\pgfsetdash{}{0pt}%
\pgfpathmoveto{\pgfqpoint{12.452492in}{0.624448in}}%
\pgfpathlineto{\pgfqpoint{12.452492in}{7.801389in}}%
\pgfusepath{stroke}%
\end{pgfscope}%
\begin{pgfscope}%
\pgfsetbuttcap%
\pgfsetroundjoin%
\definecolor{currentfill}{rgb}{0.333333,0.333333,0.333333}%
\pgfsetfillcolor{currentfill}%
\pgfsetlinewidth{0.803000pt}%
\definecolor{currentstroke}{rgb}{0.333333,0.333333,0.333333}%
\pgfsetstrokecolor{currentstroke}%
\pgfsetdash{}{0pt}%
\pgfsys@defobject{currentmarker}{\pgfqpoint{0.000000in}{-0.048611in}}{\pgfqpoint{0.000000in}{0.000000in}}{%
\pgfpathmoveto{\pgfqpoint{0.000000in}{0.000000in}}%
\pgfpathlineto{\pgfqpoint{0.000000in}{-0.048611in}}%
\pgfusepath{stroke,fill}%
}%
\begin{pgfscope}%
\pgfsys@transformshift{12.452492in}{0.624448in}%
\pgfsys@useobject{currentmarker}{}%
\end{pgfscope}%
\end{pgfscope}%
\begin{pgfscope}%
\definecolor{textcolor}{rgb}{0.333333,0.333333,0.333333}%
\pgfsetstrokecolor{textcolor}%
\pgfsetfillcolor{textcolor}%
\pgftext[x=12.070560in,y=0.221714in,left,base,rotate=30.000000]{\color{textcolor}\sffamily\fontsize{10.000000}{12.000000}\selectfont 12-Jul}%
\end{pgfscope}%
\begin{pgfscope}%
\pgfpathrectangle{\pgfqpoint{4.009158in}{0.624448in}}{\pgfqpoint{10.792231in}{7.176941in}} %
\pgfusepath{clip}%
\pgfsetrectcap%
\pgfsetroundjoin%
\pgfsetlinewidth{0.803000pt}%
\definecolor{currentstroke}{rgb}{1.000000,1.000000,1.000000}%
\pgfsetstrokecolor{currentstroke}%
\pgfsetdash{}{0pt}%
\pgfpathmoveto{\pgfqpoint{12.896878in}{0.624448in}}%
\pgfpathlineto{\pgfqpoint{12.896878in}{7.801389in}}%
\pgfusepath{stroke}%
\end{pgfscope}%
\begin{pgfscope}%
\pgfsetbuttcap%
\pgfsetroundjoin%
\definecolor{currentfill}{rgb}{0.333333,0.333333,0.333333}%
\pgfsetfillcolor{currentfill}%
\pgfsetlinewidth{0.803000pt}%
\definecolor{currentstroke}{rgb}{0.333333,0.333333,0.333333}%
\pgfsetstrokecolor{currentstroke}%
\pgfsetdash{}{0pt}%
\pgfsys@defobject{currentmarker}{\pgfqpoint{0.000000in}{-0.048611in}}{\pgfqpoint{0.000000in}{0.000000in}}{%
\pgfpathmoveto{\pgfqpoint{0.000000in}{0.000000in}}%
\pgfpathlineto{\pgfqpoint{0.000000in}{-0.048611in}}%
\pgfusepath{stroke,fill}%
}%
\begin{pgfscope}%
\pgfsys@transformshift{12.896878in}{0.624448in}%
\pgfsys@useobject{currentmarker}{}%
\end{pgfscope}%
\end{pgfscope}%
\begin{pgfscope}%
\definecolor{textcolor}{rgb}{0.333333,0.333333,0.333333}%
\pgfsetstrokecolor{textcolor}%
\pgfsetfillcolor{textcolor}%
\pgftext[x=12.514946in,y=0.221714in,left,base,rotate=30.000000]{\color{textcolor}\sffamily\fontsize{10.000000}{12.000000}\selectfont 19-Jul}%
\end{pgfscope}%
\begin{pgfscope}%
\pgfpathrectangle{\pgfqpoint{4.009158in}{0.624448in}}{\pgfqpoint{10.792231in}{7.176941in}} %
\pgfusepath{clip}%
\pgfsetrectcap%
\pgfsetroundjoin%
\pgfsetlinewidth{0.803000pt}%
\definecolor{currentstroke}{rgb}{1.000000,1.000000,1.000000}%
\pgfsetstrokecolor{currentstroke}%
\pgfsetdash{}{0pt}%
\pgfpathmoveto{\pgfqpoint{13.341264in}{0.624448in}}%
\pgfpathlineto{\pgfqpoint{13.341264in}{7.801389in}}%
\pgfusepath{stroke}%
\end{pgfscope}%
\begin{pgfscope}%
\pgfsetbuttcap%
\pgfsetroundjoin%
\definecolor{currentfill}{rgb}{0.333333,0.333333,0.333333}%
\pgfsetfillcolor{currentfill}%
\pgfsetlinewidth{0.803000pt}%
\definecolor{currentstroke}{rgb}{0.333333,0.333333,0.333333}%
\pgfsetstrokecolor{currentstroke}%
\pgfsetdash{}{0pt}%
\pgfsys@defobject{currentmarker}{\pgfqpoint{0.000000in}{-0.048611in}}{\pgfqpoint{0.000000in}{0.000000in}}{%
\pgfpathmoveto{\pgfqpoint{0.000000in}{0.000000in}}%
\pgfpathlineto{\pgfqpoint{0.000000in}{-0.048611in}}%
\pgfusepath{stroke,fill}%
}%
\begin{pgfscope}%
\pgfsys@transformshift{13.341264in}{0.624448in}%
\pgfsys@useobject{currentmarker}{}%
\end{pgfscope}%
\end{pgfscope}%
\begin{pgfscope}%
\definecolor{textcolor}{rgb}{0.333333,0.333333,0.333333}%
\pgfsetstrokecolor{textcolor}%
\pgfsetfillcolor{textcolor}%
\pgftext[x=12.959332in,y=0.221714in,left,base,rotate=30.000000]{\color{textcolor}\sffamily\fontsize{10.000000}{12.000000}\selectfont 26-Jul}%
\end{pgfscope}%
\begin{pgfscope}%
\pgfpathrectangle{\pgfqpoint{4.009158in}{0.624448in}}{\pgfqpoint{10.792231in}{7.176941in}} %
\pgfusepath{clip}%
\pgfsetrectcap%
\pgfsetroundjoin%
\pgfsetlinewidth{0.803000pt}%
\definecolor{currentstroke}{rgb}{1.000000,1.000000,1.000000}%
\pgfsetstrokecolor{currentstroke}%
\pgfsetdash{}{0pt}%
\pgfpathmoveto{\pgfqpoint{13.785650in}{0.624448in}}%
\pgfpathlineto{\pgfqpoint{13.785650in}{7.801389in}}%
\pgfusepath{stroke}%
\end{pgfscope}%
\begin{pgfscope}%
\pgfsetbuttcap%
\pgfsetroundjoin%
\definecolor{currentfill}{rgb}{0.333333,0.333333,0.333333}%
\pgfsetfillcolor{currentfill}%
\pgfsetlinewidth{0.803000pt}%
\definecolor{currentstroke}{rgb}{0.333333,0.333333,0.333333}%
\pgfsetstrokecolor{currentstroke}%
\pgfsetdash{}{0pt}%
\pgfsys@defobject{currentmarker}{\pgfqpoint{0.000000in}{-0.048611in}}{\pgfqpoint{0.000000in}{0.000000in}}{%
\pgfpathmoveto{\pgfqpoint{0.000000in}{0.000000in}}%
\pgfpathlineto{\pgfqpoint{0.000000in}{-0.048611in}}%
\pgfusepath{stroke,fill}%
}%
\begin{pgfscope}%
\pgfsys@transformshift{13.785650in}{0.624448in}%
\pgfsys@useobject{currentmarker}{}%
\end{pgfscope}%
\end{pgfscope}%
\begin{pgfscope}%
\definecolor{textcolor}{rgb}{0.333333,0.333333,0.333333}%
\pgfsetstrokecolor{textcolor}%
\pgfsetfillcolor{textcolor}%
\pgftext[x=13.317098in,y=0.173583in,left,base,rotate=30.000000]{\color{textcolor}\sffamily\fontsize{10.000000}{12.000000}\selectfont 02-Aug}%
\end{pgfscope}%
\begin{pgfscope}%
\pgfpathrectangle{\pgfqpoint{4.009158in}{0.624448in}}{\pgfqpoint{10.792231in}{7.176941in}} %
\pgfusepath{clip}%
\pgfsetrectcap%
\pgfsetroundjoin%
\pgfsetlinewidth{0.803000pt}%
\definecolor{currentstroke}{rgb}{1.000000,1.000000,1.000000}%
\pgfsetstrokecolor{currentstroke}%
\pgfsetdash{}{0pt}%
\pgfpathmoveto{\pgfqpoint{14.230035in}{0.624448in}}%
\pgfpathlineto{\pgfqpoint{14.230035in}{7.801389in}}%
\pgfusepath{stroke}%
\end{pgfscope}%
\begin{pgfscope}%
\pgfsetbuttcap%
\pgfsetroundjoin%
\definecolor{currentfill}{rgb}{0.333333,0.333333,0.333333}%
\pgfsetfillcolor{currentfill}%
\pgfsetlinewidth{0.803000pt}%
\definecolor{currentstroke}{rgb}{0.333333,0.333333,0.333333}%
\pgfsetstrokecolor{currentstroke}%
\pgfsetdash{}{0pt}%
\pgfsys@defobject{currentmarker}{\pgfqpoint{0.000000in}{-0.048611in}}{\pgfqpoint{0.000000in}{0.000000in}}{%
\pgfpathmoveto{\pgfqpoint{0.000000in}{0.000000in}}%
\pgfpathlineto{\pgfqpoint{0.000000in}{-0.048611in}}%
\pgfusepath{stroke,fill}%
}%
\begin{pgfscope}%
\pgfsys@transformshift{14.230035in}{0.624448in}%
\pgfsys@useobject{currentmarker}{}%
\end{pgfscope}%
\end{pgfscope}%
\begin{pgfscope}%
\definecolor{textcolor}{rgb}{0.333333,0.333333,0.333333}%
\pgfsetstrokecolor{textcolor}%
\pgfsetfillcolor{textcolor}%
\pgftext[x=13.761484in,y=0.173583in,left,base,rotate=30.000000]{\color{textcolor}\sffamily\fontsize{10.000000}{12.000000}\selectfont 09-Aug}%
\end{pgfscope}%
\begin{pgfscope}%
\pgfpathrectangle{\pgfqpoint{4.009158in}{0.624448in}}{\pgfqpoint{10.792231in}{7.176941in}} %
\pgfusepath{clip}%
\pgfsetrectcap%
\pgfsetroundjoin%
\pgfsetlinewidth{0.803000pt}%
\definecolor{currentstroke}{rgb}{1.000000,1.000000,1.000000}%
\pgfsetstrokecolor{currentstroke}%
\pgfsetdash{}{0pt}%
\pgfpathmoveto{\pgfqpoint{14.674421in}{0.624448in}}%
\pgfpathlineto{\pgfqpoint{14.674421in}{7.801389in}}%
\pgfusepath{stroke}%
\end{pgfscope}%
\begin{pgfscope}%
\pgfsetbuttcap%
\pgfsetroundjoin%
\definecolor{currentfill}{rgb}{0.333333,0.333333,0.333333}%
\pgfsetfillcolor{currentfill}%
\pgfsetlinewidth{0.803000pt}%
\definecolor{currentstroke}{rgb}{0.333333,0.333333,0.333333}%
\pgfsetstrokecolor{currentstroke}%
\pgfsetdash{}{0pt}%
\pgfsys@defobject{currentmarker}{\pgfqpoint{0.000000in}{-0.048611in}}{\pgfqpoint{0.000000in}{0.000000in}}{%
\pgfpathmoveto{\pgfqpoint{0.000000in}{0.000000in}}%
\pgfpathlineto{\pgfqpoint{0.000000in}{-0.048611in}}%
\pgfusepath{stroke,fill}%
}%
\begin{pgfscope}%
\pgfsys@transformshift{14.674421in}{0.624448in}%
\pgfsys@useobject{currentmarker}{}%
\end{pgfscope}%
\end{pgfscope}%
\begin{pgfscope}%
\definecolor{textcolor}{rgb}{0.333333,0.333333,0.333333}%
\pgfsetstrokecolor{textcolor}%
\pgfsetfillcolor{textcolor}%
\pgftext[x=14.205870in,y=0.173583in,left,base,rotate=30.000000]{\color{textcolor}\sffamily\fontsize{10.000000}{12.000000}\selectfont 16-Aug}%
\end{pgfscope}%
\begin{pgfscope}%
\pgfpathrectangle{\pgfqpoint{4.009158in}{0.624448in}}{\pgfqpoint{10.792231in}{7.176941in}} %
\pgfusepath{clip}%
\pgfsetrectcap%
\pgfsetroundjoin%
\pgfsetlinewidth{0.803000pt}%
\definecolor{currentstroke}{rgb}{1.000000,1.000000,1.000000}%
\pgfsetstrokecolor{currentstroke}%
\pgfsetdash{}{0pt}%
\pgfpathmoveto{\pgfqpoint{4.009158in}{7.420183in}}%
\pgfpathlineto{\pgfqpoint{14.801389in}{7.420183in}}%
\pgfusepath{stroke}%
\end{pgfscope}%
\begin{pgfscope}%
\definecolor{textcolor}{rgb}{0.333333,0.333333,0.333333}%
\pgfsetstrokecolor{textcolor}%
\pgfsetfillcolor{textcolor}%
\pgftext[x=0.717152in,y=7.366641in,left,base]{\color{textcolor}\sffamily\fontsize{10.000000}{12.000000}\selectfont 1.1.1 Read Literature about the Existing Project}%
\end{pgfscope}%
\begin{pgfscope}%
\pgfpathrectangle{\pgfqpoint{4.009158in}{0.624448in}}{\pgfqpoint{10.792231in}{7.176941in}} %
\pgfusepath{clip}%
\pgfsetrectcap%
\pgfsetroundjoin%
\pgfsetlinewidth{0.803000pt}%
\definecolor{currentstroke}{rgb}{1.000000,1.000000,1.000000}%
\pgfsetstrokecolor{currentstroke}%
\pgfsetdash{}{0pt}%
\pgfpathmoveto{\pgfqpoint{4.009158in}{7.236910in}}%
\pgfpathlineto{\pgfqpoint{14.801389in}{7.236910in}}%
\pgfusepath{stroke}%
\end{pgfscope}%
\begin{pgfscope}%
\definecolor{textcolor}{rgb}{0.333333,0.333333,0.333333}%
\pgfsetstrokecolor{textcolor}%
\pgfsetfillcolor{textcolor}%
\pgftext[x=2.072065in,y=7.183369in,left,base]{\color{textcolor}\sffamily\fontsize{10.000000}{12.000000}\selectfont 1.1.2 Tutorials to IIR Filters}%
\end{pgfscope}%
\begin{pgfscope}%
\pgfpathrectangle{\pgfqpoint{4.009158in}{0.624448in}}{\pgfqpoint{10.792231in}{7.176941in}} %
\pgfusepath{clip}%
\pgfsetrectcap%
\pgfsetroundjoin%
\pgfsetlinewidth{0.803000pt}%
\definecolor{currentstroke}{rgb}{1.000000,1.000000,1.000000}%
\pgfsetstrokecolor{currentstroke}%
\pgfsetdash{}{0pt}%
\pgfpathmoveto{\pgfqpoint{4.009158in}{7.053638in}}%
\pgfpathlineto{\pgfqpoint{14.801389in}{7.053638in}}%
\pgfusepath{stroke}%
\end{pgfscope}%
\begin{pgfscope}%
\definecolor{textcolor}{rgb}{0.333333,0.333333,0.333333}%
\pgfsetstrokecolor{textcolor}%
\pgfsetfillcolor{textcolor}%
\pgftext[x=0.733835in,y=6.999249in,left,base]{\color{textcolor}\sffamily\fontsize{10.000000}{12.000000}\selectfont 1.1.3 Assess Results of the Predecessor Project}%
\end{pgfscope}%
\begin{pgfscope}%
\pgfpathrectangle{\pgfqpoint{4.009158in}{0.624448in}}{\pgfqpoint{10.792231in}{7.176941in}} %
\pgfusepath{clip}%
\pgfsetrectcap%
\pgfsetroundjoin%
\pgfsetlinewidth{0.803000pt}%
\definecolor{currentstroke}{rgb}{1.000000,1.000000,1.000000}%
\pgfsetstrokecolor{currentstroke}%
\pgfsetdash{}{0pt}%
\pgfpathmoveto{\pgfqpoint{4.009158in}{6.870366in}}%
\pgfpathlineto{\pgfqpoint{14.801389in}{6.870366in}}%
\pgfusepath{stroke}%
\end{pgfscope}%
\begin{pgfscope}%
\definecolor{textcolor}{rgb}{0.333333,0.333333,0.333333}%
\pgfsetstrokecolor{textcolor}%
\pgfsetfillcolor{textcolor}%
\pgftext[x=2.261681in,y=6.816825in,left,base]{\color{textcolor}\sffamily\fontsize{10.000000}{12.000000}\selectfont 1.1.4 Outline the Project}%
\end{pgfscope}%
\begin{pgfscope}%
\pgfpathrectangle{\pgfqpoint{4.009158in}{0.624448in}}{\pgfqpoint{10.792231in}{7.176941in}} %
\pgfusepath{clip}%
\pgfsetrectcap%
\pgfsetroundjoin%
\pgfsetlinewidth{0.803000pt}%
\definecolor{currentstroke}{rgb}{1.000000,1.000000,1.000000}%
\pgfsetstrokecolor{currentstroke}%
\pgfsetdash{}{0pt}%
\pgfpathmoveto{\pgfqpoint{4.009158in}{6.687094in}}%
\pgfpathlineto{\pgfqpoint{14.801389in}{6.687094in}}%
\pgfusepath{stroke}%
\end{pgfscope}%
\begin{pgfscope}%
\definecolor{textcolor}{rgb}{0.333333,0.333333,0.333333}%
\pgfsetstrokecolor{textcolor}%
\pgfsetfillcolor{textcolor}%
\pgftext[x=2.527523in,y=6.633552in,left,base]{\color{textcolor}\sffamily\fontsize{10.000000}{12.000000}\selectfont 1.2.1 Filter Research}%
\end{pgfscope}%
\begin{pgfscope}%
\pgfpathrectangle{\pgfqpoint{4.009158in}{0.624448in}}{\pgfqpoint{10.792231in}{7.176941in}} %
\pgfusepath{clip}%
\pgfsetrectcap%
\pgfsetroundjoin%
\pgfsetlinewidth{0.803000pt}%
\definecolor{currentstroke}{rgb}{1.000000,1.000000,1.000000}%
\pgfsetstrokecolor{currentstroke}%
\pgfsetdash{}{0pt}%
\pgfpathmoveto{\pgfqpoint{4.009158in}{6.503821in}}%
\pgfpathlineto{\pgfqpoint{14.801389in}{6.503821in}}%
\pgfusepath{stroke}%
\end{pgfscope}%
\begin{pgfscope}%
\definecolor{textcolor}{rgb}{0.333333,0.333333,0.333333}%
\pgfsetstrokecolor{textcolor}%
\pgfsetfillcolor{textcolor}%
\pgftext[x=1.464223in,y=6.449432in,left,base]{\color{textcolor}\sffamily\fontsize{10.000000}{12.000000}\selectfont 1.2.2 Determine Filter Specifications}%
\end{pgfscope}%
\begin{pgfscope}%
\pgfpathrectangle{\pgfqpoint{4.009158in}{0.624448in}}{\pgfqpoint{10.792231in}{7.176941in}} %
\pgfusepath{clip}%
\pgfsetrectcap%
\pgfsetroundjoin%
\pgfsetlinewidth{0.803000pt}%
\definecolor{currentstroke}{rgb}{1.000000,1.000000,1.000000}%
\pgfsetstrokecolor{currentstroke}%
\pgfsetdash{}{0pt}%
\pgfpathmoveto{\pgfqpoint{4.009158in}{6.320549in}}%
\pgfpathlineto{\pgfqpoint{14.801389in}{6.320549in}}%
\pgfusepath{stroke}%
\end{pgfscope}%
\begin{pgfscope}%
\definecolor{textcolor}{rgb}{0.333333,0.333333,0.333333}%
\pgfsetstrokecolor{textcolor}%
\pgfsetfillcolor{textcolor}%
\pgftext[x=1.435943in,y=6.266160in,left,base]{\color{textcolor}\sffamily\fontsize{10.000000}{12.000000}\selectfont 1.2.3 Matlab Scripts for Filter Design}%
\end{pgfscope}%
\begin{pgfscope}%
\pgfpathrectangle{\pgfqpoint{4.009158in}{0.624448in}}{\pgfqpoint{10.792231in}{7.176941in}} %
\pgfusepath{clip}%
\pgfsetrectcap%
\pgfsetroundjoin%
\pgfsetlinewidth{0.803000pt}%
\definecolor{currentstroke}{rgb}{1.000000,1.000000,1.000000}%
\pgfsetstrokecolor{currentstroke}%
\pgfsetdash{}{0pt}%
\pgfpathmoveto{\pgfqpoint{4.009158in}{6.137277in}}%
\pgfpathlineto{\pgfqpoint{14.801389in}{6.137277in}}%
\pgfusepath{stroke}%
\end{pgfscope}%
\begin{pgfscope}%
\definecolor{textcolor}{rgb}{0.333333,0.333333,0.333333}%
\pgfsetstrokecolor{textcolor}%
\pgfsetfillcolor{textcolor}%
\pgftext[x=1.405154in,y=6.082888in,left,base]{\color{textcolor}\sffamily\fontsize{10.000000}{12.000000}\selectfont 1.1.4 TCL Scripts for Filter Evaluation}%
\end{pgfscope}%
\begin{pgfscope}%
\pgfpathrectangle{\pgfqpoint{4.009158in}{0.624448in}}{\pgfqpoint{10.792231in}{7.176941in}} %
\pgfusepath{clip}%
\pgfsetrectcap%
\pgfsetroundjoin%
\pgfsetlinewidth{0.803000pt}%
\definecolor{currentstroke}{rgb}{1.000000,1.000000,1.000000}%
\pgfsetstrokecolor{currentstroke}%
\pgfsetdash{}{0pt}%
\pgfpathmoveto{\pgfqpoint{4.009158in}{5.954005in}}%
\pgfpathlineto{\pgfqpoint{14.801389in}{5.954005in}}%
\pgfusepath{stroke}%
\end{pgfscope}%
\begin{pgfscope}%
\definecolor{textcolor}{rgb}{0.333333,0.333333,0.333333}%
\pgfsetstrokecolor{textcolor}%
\pgfsetfillcolor{textcolor}%
\pgftext[x=0.532487in,y=5.900463in,left,base]{\color{textcolor}\sffamily\fontsize{10.000000}{12.000000}\selectfont 1.1.5 Document Filter Design and Implementation}%
\end{pgfscope}%
\begin{pgfscope}%
\pgfpathrectangle{\pgfqpoint{4.009158in}{0.624448in}}{\pgfqpoint{10.792231in}{7.176941in}} %
\pgfusepath{clip}%
\pgfsetrectcap%
\pgfsetroundjoin%
\pgfsetlinewidth{0.803000pt}%
\definecolor{currentstroke}{rgb}{1.000000,1.000000,1.000000}%
\pgfsetstrokecolor{currentstroke}%
\pgfsetdash{}{0pt}%
\pgfpathmoveto{\pgfqpoint{4.009158in}{5.770732in}}%
\pgfpathlineto{\pgfqpoint{14.801389in}{5.770732in}}%
\pgfusepath{stroke}%
\end{pgfscope}%
\begin{pgfscope}%
\definecolor{textcolor}{rgb}{0.333333,0.333333,0.333333}%
\pgfsetstrokecolor{textcolor}%
\pgfsetfillcolor{textcolor}%
\pgftext[x=2.754370in,y=5.717191in,left,base]{\color{textcolor}\sffamily\fontsize{10.000000}{12.000000}\selectfont 1.3.1 Disposition}%
\end{pgfscope}%
\begin{pgfscope}%
\pgfpathrectangle{\pgfqpoint{4.009158in}{0.624448in}}{\pgfqpoint{10.792231in}{7.176941in}} %
\pgfusepath{clip}%
\pgfsetrectcap%
\pgfsetroundjoin%
\pgfsetlinewidth{0.803000pt}%
\definecolor{currentstroke}{rgb}{1.000000,1.000000,1.000000}%
\pgfsetstrokecolor{currentstroke}%
\pgfsetdash{}{0pt}%
\pgfpathmoveto{\pgfqpoint{4.009158in}{5.587460in}}%
\pgfpathlineto{\pgfqpoint{14.801389in}{5.587460in}}%
\pgfusepath{stroke}%
\end{pgfscope}%
\begin{pgfscope}%
\definecolor{textcolor}{rgb}{0.333333,0.333333,0.333333}%
\pgfsetstrokecolor{textcolor}%
\pgfsetfillcolor{textcolor}%
\pgftext[x=3.069922in,y=5.533919in,left,base]{\color{textcolor}\sffamily\fontsize{10.000000}{12.000000}\selectfont 1.3.2 Report}%
\end{pgfscope}%
\begin{pgfscope}%
\pgfpathrectangle{\pgfqpoint{4.009158in}{0.624448in}}{\pgfqpoint{10.792231in}{7.176941in}} %
\pgfusepath{clip}%
\pgfsetrectcap%
\pgfsetroundjoin%
\pgfsetlinewidth{0.803000pt}%
\definecolor{currentstroke}{rgb}{1.000000,1.000000,1.000000}%
\pgfsetstrokecolor{currentstroke}%
\pgfsetdash{}{0pt}%
\pgfpathmoveto{\pgfqpoint{4.009158in}{5.404188in}}%
\pgfpathlineto{\pgfqpoint{14.801389in}{5.404188in}}%
\pgfusepath{stroke}%
\end{pgfscope}%
\begin{pgfscope}%
\definecolor{textcolor}{rgb}{0.333333,0.333333,0.333333}%
\pgfsetstrokecolor{textcolor}%
\pgfsetfillcolor{textcolor}%
\pgftext[x=3.144860in,y=5.350647in,left,base]{\color{textcolor}\sffamily\fontsize{10.000000}{12.000000}\selectfont 1.4.1 Linux}%
\end{pgfscope}%
\begin{pgfscope}%
\pgfpathrectangle{\pgfqpoint{4.009158in}{0.624448in}}{\pgfqpoint{10.792231in}{7.176941in}} %
\pgfusepath{clip}%
\pgfsetrectcap%
\pgfsetroundjoin%
\pgfsetlinewidth{0.803000pt}%
\definecolor{currentstroke}{rgb}{1.000000,1.000000,1.000000}%
\pgfsetstrokecolor{currentstroke}%
\pgfsetdash{}{0pt}%
\pgfpathmoveto{\pgfqpoint{4.009158in}{5.220916in}}%
\pgfpathlineto{\pgfqpoint{14.801389in}{5.220916in}}%
\pgfusepath{stroke}%
\end{pgfscope}%
\begin{pgfscope}%
\definecolor{textcolor}{rgb}{0.333333,0.333333,0.333333}%
\pgfsetstrokecolor{textcolor}%
\pgfsetfillcolor{textcolor}%
\pgftext[x=1.889638in,y=5.167374in,left,base]{\color{textcolor}\sffamily\fontsize{10.000000}{12.000000}\selectfont 1.4.2.1 Server Design Choices}%
\end{pgfscope}%
\begin{pgfscope}%
\pgfpathrectangle{\pgfqpoint{4.009158in}{0.624448in}}{\pgfqpoint{10.792231in}{7.176941in}} %
\pgfusepath{clip}%
\pgfsetrectcap%
\pgfsetroundjoin%
\pgfsetlinewidth{0.803000pt}%
\definecolor{currentstroke}{rgb}{1.000000,1.000000,1.000000}%
\pgfsetstrokecolor{currentstroke}%
\pgfsetdash{}{0pt}%
\pgfpathmoveto{\pgfqpoint{4.009158in}{5.037644in}}%
\pgfpathlineto{\pgfqpoint{14.801389in}{5.037644in}}%
\pgfusepath{stroke}%
\end{pgfscope}%
\begin{pgfscope}%
\definecolor{textcolor}{rgb}{0.333333,0.333333,0.333333}%
\pgfsetstrokecolor{textcolor}%
\pgfsetfillcolor{textcolor}%
\pgftext[x=0.817928in,y=4.984102in,left,base]{\color{textcolor}\sffamily\fontsize{10.000000}{12.000000}\selectfont 1.4.2.2 Build External Libraries and Test them}%
\end{pgfscope}%
\begin{pgfscope}%
\pgfpathrectangle{\pgfqpoint{4.009158in}{0.624448in}}{\pgfqpoint{10.792231in}{7.176941in}} %
\pgfusepath{clip}%
\pgfsetrectcap%
\pgfsetroundjoin%
\pgfsetlinewidth{0.803000pt}%
\definecolor{currentstroke}{rgb}{1.000000,1.000000,1.000000}%
\pgfsetstrokecolor{currentstroke}%
\pgfsetdash{}{0pt}%
\pgfpathmoveto{\pgfqpoint{4.009158in}{4.854371in}}%
\pgfpathlineto{\pgfqpoint{14.801389in}{4.854371in}}%
\pgfusepath{stroke}%
\end{pgfscope}%
\begin{pgfscope}%
\definecolor{textcolor}{rgb}{0.333333,0.333333,0.333333}%
\pgfsetstrokecolor{textcolor}%
\pgfsetfillcolor{textcolor}%
\pgftext[x=1.514882in,y=4.800830in,left,base]{\color{textcolor}\sffamily\fontsize{10.000000}{12.000000}\selectfont 1.4.2.3 Write the Server Application}%
\end{pgfscope}%
\begin{pgfscope}%
\pgfpathrectangle{\pgfqpoint{4.009158in}{0.624448in}}{\pgfqpoint{10.792231in}{7.176941in}} %
\pgfusepath{clip}%
\pgfsetrectcap%
\pgfsetroundjoin%
\pgfsetlinewidth{0.803000pt}%
\definecolor{currentstroke}{rgb}{1.000000,1.000000,1.000000}%
\pgfsetstrokecolor{currentstroke}%
\pgfsetdash{}{0pt}%
\pgfpathmoveto{\pgfqpoint{4.009158in}{4.671099in}}%
\pgfpathlineto{\pgfqpoint{14.801389in}{4.671099in}}%
\pgfusepath{stroke}%
\end{pgfscope}%
\begin{pgfscope}%
\definecolor{textcolor}{rgb}{0.333333,0.333333,0.333333}%
\pgfsetstrokecolor{textcolor}%
\pgfsetfillcolor{textcolor}%
\pgftext[x=1.168134in,y=4.617558in,left,base]{\color{textcolor}\sffamily\fontsize{10.000000}{12.000000}\selectfont 1.4.2.4 Document the Server Application}%
\end{pgfscope}%
\begin{pgfscope}%
\pgfpathrectangle{\pgfqpoint{4.009158in}{0.624448in}}{\pgfqpoint{10.792231in}{7.176941in}} %
\pgfusepath{clip}%
\pgfsetrectcap%
\pgfsetroundjoin%
\pgfsetlinewidth{0.803000pt}%
\definecolor{currentstroke}{rgb}{1.000000,1.000000,1.000000}%
\pgfsetstrokecolor{currentstroke}%
\pgfsetdash{}{0pt}%
\pgfpathmoveto{\pgfqpoint{4.009158in}{4.487827in}}%
\pgfpathlineto{\pgfqpoint{14.801389in}{4.487827in}}%
\pgfusepath{stroke}%
\end{pgfscope}%
\begin{pgfscope}%
\definecolor{textcolor}{rgb}{0.333333,0.333333,0.333333}%
\pgfsetstrokecolor{textcolor}%
\pgfsetfillcolor{textcolor}%
\pgftext[x=0.749501in,y=4.433438in,left,base]{\color{textcolor}\sffamily\fontsize{10.000000}{12.000000}\selectfont 1.4.3.1 Minimal Working Example for the FPGA}%
\end{pgfscope}%
\begin{pgfscope}%
\pgfpathrectangle{\pgfqpoint{4.009158in}{0.624448in}}{\pgfqpoint{10.792231in}{7.176941in}} %
\pgfusepath{clip}%
\pgfsetrectcap%
\pgfsetroundjoin%
\pgfsetlinewidth{0.803000pt}%
\definecolor{currentstroke}{rgb}{1.000000,1.000000,1.000000}%
\pgfsetstrokecolor{currentstroke}%
\pgfsetdash{}{0pt}%
\pgfpathmoveto{\pgfqpoint{4.009158in}{4.304555in}}%
\pgfpathlineto{\pgfqpoint{14.801389in}{4.304555in}}%
\pgfusepath{stroke}%
\end{pgfscope}%
\begin{pgfscope}%
\definecolor{textcolor}{rgb}{0.333333,0.333333,0.333333}%
\pgfsetstrokecolor{textcolor}%
\pgfsetfillcolor{textcolor}%
\pgftext[x=0.892594in,y=4.250165in,left,base]{\color{textcolor}\sffamily\fontsize{10.000000}{12.000000}\selectfont 1.4.3.2 Set up a Sane Toolchain for the FPGA}%
\end{pgfscope}%
\begin{pgfscope}%
\pgfpathrectangle{\pgfqpoint{4.009158in}{0.624448in}}{\pgfqpoint{10.792231in}{7.176941in}} %
\pgfusepath{clip}%
\pgfsetrectcap%
\pgfsetroundjoin%
\pgfsetlinewidth{0.803000pt}%
\definecolor{currentstroke}{rgb}{1.000000,1.000000,1.000000}%
\pgfsetstrokecolor{currentstroke}%
\pgfsetdash{}{0pt}%
\pgfpathmoveto{\pgfqpoint{4.009158in}{4.121282in}}%
\pgfpathlineto{\pgfqpoint{14.801389in}{4.121282in}}%
\pgfusepath{stroke}%
\end{pgfscope}%
\begin{pgfscope}%
\definecolor{textcolor}{rgb}{0.333333,0.333333,0.333333}%
\pgfsetstrokecolor{textcolor}%
\pgfsetfillcolor{textcolor}%
\pgftext[x=1.918867in,y=4.067741in,left,base]{\color{textcolor}\sffamily\fontsize{10.000000}{12.000000}\selectfont 1.4.3.2 Port the ZYNQ Logger}%
\end{pgfscope}%
\begin{pgfscope}%
\pgfpathrectangle{\pgfqpoint{4.009158in}{0.624448in}}{\pgfqpoint{10.792231in}{7.176941in}} %
\pgfusepath{clip}%
\pgfsetrectcap%
\pgfsetroundjoin%
\pgfsetlinewidth{0.803000pt}%
\definecolor{currentstroke}{rgb}{1.000000,1.000000,1.000000}%
\pgfsetstrokecolor{currentstroke}%
\pgfsetdash{}{0pt}%
\pgfpathmoveto{\pgfqpoint{4.009158in}{3.938010in}}%
\pgfpathlineto{\pgfqpoint{14.801389in}{3.938010in}}%
\pgfusepath{stroke}%
\end{pgfscope}%
\begin{pgfscope}%
\definecolor{textcolor}{rgb}{0.333333,0.333333,0.333333}%
\pgfsetstrokecolor{textcolor}%
\pgfsetfillcolor{textcolor}%
\pgftext[x=1.520646in,y=3.884469in,left,base]{\color{textcolor}\sffamily\fontsize{10.000000}{12.000000}\selectfont 1.4.3.3 Implement the Filter Chains}%
\end{pgfscope}%
\begin{pgfscope}%
\pgfpathrectangle{\pgfqpoint{4.009158in}{0.624448in}}{\pgfqpoint{10.792231in}{7.176941in}} %
\pgfusepath{clip}%
\pgfsetrectcap%
\pgfsetroundjoin%
\pgfsetlinewidth{0.803000pt}%
\definecolor{currentstroke}{rgb}{1.000000,1.000000,1.000000}%
\pgfsetstrokecolor{currentstroke}%
\pgfsetdash{}{0pt}%
\pgfpathmoveto{\pgfqpoint{4.009158in}{3.754738in}}%
\pgfpathlineto{\pgfqpoint{14.801389in}{3.754738in}}%
\pgfusepath{stroke}%
\end{pgfscope}%
\begin{pgfscope}%
\definecolor{textcolor}{rgb}{0.333333,0.333333,0.333333}%
\pgfsetstrokecolor{textcolor}%
\pgfsetfillcolor{textcolor}%
\pgftext[x=1.659603in,y=3.700349in,left,base]{\color{textcolor}\sffamily\fontsize{10.000000}{12.000000}\selectfont 1.4.3.4 Verification by Simulation}%
\end{pgfscope}%
\begin{pgfscope}%
\pgfpathrectangle{\pgfqpoint{4.009158in}{0.624448in}}{\pgfqpoint{10.792231in}{7.176941in}} %
\pgfusepath{clip}%
\pgfsetrectcap%
\pgfsetroundjoin%
\pgfsetlinewidth{0.803000pt}%
\definecolor{currentstroke}{rgb}{1.000000,1.000000,1.000000}%
\pgfsetstrokecolor{currentstroke}%
\pgfsetdash{}{0pt}%
\pgfpathmoveto{\pgfqpoint{4.009158in}{3.571466in}}%
\pgfpathlineto{\pgfqpoint{14.801389in}{3.571466in}}%
\pgfusepath{stroke}%
\end{pgfscope}%
\begin{pgfscope}%
\definecolor{textcolor}{rgb}{0.333333,0.333333,0.333333}%
\pgfsetstrokecolor{textcolor}%
\pgfsetfillcolor{textcolor}%
\pgftext[x=0.262373in,y=3.517924in,left,base]{\color{textcolor}\sffamily\fontsize{10.000000}{12.000000}\selectfont 1.4.3.3 Document FPGA Toolchain and Filter Structure}%
\end{pgfscope}%
\begin{pgfscope}%
\pgfpathrectangle{\pgfqpoint{4.009158in}{0.624448in}}{\pgfqpoint{10.792231in}{7.176941in}} %
\pgfusepath{clip}%
\pgfsetrectcap%
\pgfsetroundjoin%
\pgfsetlinewidth{0.803000pt}%
\definecolor{currentstroke}{rgb}{1.000000,1.000000,1.000000}%
\pgfsetstrokecolor{currentstroke}%
\pgfsetdash{}{0pt}%
\pgfpathmoveto{\pgfqpoint{4.009158in}{3.388193in}}%
\pgfpathlineto{\pgfqpoint{14.801389in}{3.388193in}}%
\pgfusepath{stroke}%
\end{pgfscope}%
\begin{pgfscope}%
\definecolor{textcolor}{rgb}{0.333333,0.333333,0.333333}%
\pgfsetstrokecolor{textcolor}%
\pgfsetfillcolor{textcolor}%
\pgftext[x=1.653432in,y=3.334652in,left,base]{\color{textcolor}\sffamily\fontsize{10.000000}{12.000000}\selectfont 1.5.0.1 Validate the Build Process}%
\end{pgfscope}%
\begin{pgfscope}%
\pgfpathrectangle{\pgfqpoint{4.009158in}{0.624448in}}{\pgfqpoint{10.792231in}{7.176941in}} %
\pgfusepath{clip}%
\pgfsetrectcap%
\pgfsetroundjoin%
\pgfsetlinewidth{0.803000pt}%
\definecolor{currentstroke}{rgb}{1.000000,1.000000,1.000000}%
\pgfsetstrokecolor{currentstroke}%
\pgfsetdash{}{0pt}%
\pgfpathmoveto{\pgfqpoint{4.009158in}{3.204921in}}%
\pgfpathlineto{\pgfqpoint{14.801389in}{3.204921in}}%
\pgfusepath{stroke}%
\end{pgfscope}%
\begin{pgfscope}%
\definecolor{textcolor}{rgb}{0.333333,0.333333,0.333333}%
\pgfsetstrokecolor{textcolor}%
\pgfsetfillcolor{textcolor}%
\pgftext[x=1.095231in,y=3.151380in,left,base]{\color{textcolor}\sffamily\fontsize{10.000000}{12.000000}\selectfont 1.5.0.2 Validate the Filter Implementation}%
\end{pgfscope}%
\begin{pgfscope}%
\pgfpathrectangle{\pgfqpoint{4.009158in}{0.624448in}}{\pgfqpoint{10.792231in}{7.176941in}} %
\pgfusepath{clip}%
\pgfsetrectcap%
\pgfsetroundjoin%
\pgfsetlinewidth{0.803000pt}%
\definecolor{currentstroke}{rgb}{1.000000,1.000000,1.000000}%
\pgfsetstrokecolor{currentstroke}%
\pgfsetdash{}{0pt}%
\pgfpathmoveto{\pgfqpoint{4.009158in}{3.021649in}}%
\pgfpathlineto{\pgfqpoint{14.801389in}{3.021649in}}%
\pgfusepath{stroke}%
\end{pgfscope}%
\begin{pgfscope}%
\definecolor{textcolor}{rgb}{0.333333,0.333333,0.333333}%
\pgfsetstrokecolor{textcolor}%
\pgfsetfillcolor{textcolor}%
\pgftext[x=2.089765in,y=2.968107in,left,base]{\color{textcolor}\sffamily\fontsize{10.000000}{12.000000}\selectfont 1.6.1 Technology Research}%
\end{pgfscope}%
\begin{pgfscope}%
\pgfpathrectangle{\pgfqpoint{4.009158in}{0.624448in}}{\pgfqpoint{10.792231in}{7.176941in}} %
\pgfusepath{clip}%
\pgfsetrectcap%
\pgfsetroundjoin%
\pgfsetlinewidth{0.803000pt}%
\definecolor{currentstroke}{rgb}{1.000000,1.000000,1.000000}%
\pgfsetstrokecolor{currentstroke}%
\pgfsetdash{}{0pt}%
\pgfpathmoveto{\pgfqpoint{4.009158in}{2.838377in}}%
\pgfpathlineto{\pgfqpoint{14.801389in}{2.838377in}}%
\pgfusepath{stroke}%
\end{pgfscope}%
\begin{pgfscope}%
\definecolor{textcolor}{rgb}{0.333333,0.333333,0.333333}%
\pgfsetstrokecolor{textcolor}%
\pgfsetfillcolor{textcolor}%
\pgftext[x=2.153988in,y=2.783988in,left,base]{\color{textcolor}\sffamily\fontsize{10.000000}{12.000000}\selectfont 1.6.2 Scope Specifications}%
\end{pgfscope}%
\begin{pgfscope}%
\pgfpathrectangle{\pgfqpoint{4.009158in}{0.624448in}}{\pgfqpoint{10.792231in}{7.176941in}} %
\pgfusepath{clip}%
\pgfsetrectcap%
\pgfsetroundjoin%
\pgfsetlinewidth{0.803000pt}%
\definecolor{currentstroke}{rgb}{1.000000,1.000000,1.000000}%
\pgfsetstrokecolor{currentstroke}%
\pgfsetdash{}{0pt}%
\pgfpathmoveto{\pgfqpoint{4.009158in}{2.655104in}}%
\pgfpathlineto{\pgfqpoint{14.801389in}{2.655104in}}%
\pgfusepath{stroke}%
\end{pgfscope}%
\begin{pgfscope}%
\definecolor{textcolor}{rgb}{0.333333,0.333333,0.333333}%
\pgfsetstrokecolor{textcolor}%
\pgfsetfillcolor{textcolor}%
\pgftext[x=1.288645in,y=2.600715in,left,base]{\color{textcolor}\sffamily\fontsize{10.000000}{12.000000}\selectfont 1.6.3 Implementation of a Basic Plotter}%
\end{pgfscope}%
\begin{pgfscope}%
\pgfpathrectangle{\pgfqpoint{4.009158in}{0.624448in}}{\pgfqpoint{10.792231in}{7.176941in}} %
\pgfusepath{clip}%
\pgfsetrectcap%
\pgfsetroundjoin%
\pgfsetlinewidth{0.803000pt}%
\definecolor{currentstroke}{rgb}{1.000000,1.000000,1.000000}%
\pgfsetstrokecolor{currentstroke}%
\pgfsetdash{}{0pt}%
\pgfpathmoveto{\pgfqpoint{4.009158in}{2.471832in}}%
\pgfpathlineto{\pgfqpoint{14.801389in}{2.471832in}}%
\pgfusepath{stroke}%
\end{pgfscope}%
\begin{pgfscope}%
\definecolor{textcolor}{rgb}{0.333333,0.333333,0.333333}%
\pgfsetstrokecolor{textcolor}%
\pgfsetfillcolor{textcolor}%
\pgftext[x=1.169016in,y=2.417443in,left,base]{\color{textcolor}\sffamily\fontsize{10.000000}{12.000000}\selectfont 1.6.4 Implementation of a Basic Receiver}%
\end{pgfscope}%
\begin{pgfscope}%
\pgfpathrectangle{\pgfqpoint{4.009158in}{0.624448in}}{\pgfqpoint{10.792231in}{7.176941in}} %
\pgfusepath{clip}%
\pgfsetrectcap%
\pgfsetroundjoin%
\pgfsetlinewidth{0.803000pt}%
\definecolor{currentstroke}{rgb}{1.000000,1.000000,1.000000}%
\pgfsetstrokecolor{currentstroke}%
\pgfsetdash{}{0pt}%
\pgfpathmoveto{\pgfqpoint{4.009158in}{2.288560in}}%
\pgfpathlineto{\pgfqpoint{14.801389in}{2.288560in}}%
\pgfusepath{stroke}%
\end{pgfscope}%
\begin{pgfscope}%
\definecolor{textcolor}{rgb}{0.333333,0.333333,0.333333}%
\pgfsetstrokecolor{textcolor}%
\pgfsetfillcolor{textcolor}%
\pgftext[x=0.877471in,y=2.235019in,left,base]{\color{textcolor}\sffamily\fontsize{10.000000}{12.000000}\selectfont 1.6.5 Combining the Plotter and the Receiver}%
\end{pgfscope}%
\begin{pgfscope}%
\pgfpathrectangle{\pgfqpoint{4.009158in}{0.624448in}}{\pgfqpoint{10.792231in}{7.176941in}} %
\pgfusepath{clip}%
\pgfsetrectcap%
\pgfsetroundjoin%
\pgfsetlinewidth{0.803000pt}%
\definecolor{currentstroke}{rgb}{1.000000,1.000000,1.000000}%
\pgfsetstrokecolor{currentstroke}%
\pgfsetdash{}{0pt}%
\pgfpathmoveto{\pgfqpoint{4.009158in}{2.105288in}}%
\pgfpathlineto{\pgfqpoint{14.801389in}{2.105288in}}%
\pgfusepath{stroke}%
\end{pgfscope}%
\begin{pgfscope}%
\definecolor{textcolor}{rgb}{0.333333,0.333333,0.333333}%
\pgfsetstrokecolor{textcolor}%
\pgfsetfillcolor{textcolor}%
\pgftext[x=1.045589in,y=2.051746in,left,base]{\color{textcolor}\sffamily\fontsize{10.000000}{12.000000}\selectfont 1.6.6 Setting a Trigger on the Logger Core}%
\end{pgfscope}%
\begin{pgfscope}%
\pgfpathrectangle{\pgfqpoint{4.009158in}{0.624448in}}{\pgfqpoint{10.792231in}{7.176941in}} %
\pgfusepath{clip}%
\pgfsetrectcap%
\pgfsetroundjoin%
\pgfsetlinewidth{0.803000pt}%
\definecolor{currentstroke}{rgb}{1.000000,1.000000,1.000000}%
\pgfsetstrokecolor{currentstroke}%
\pgfsetdash{}{0pt}%
\pgfpathmoveto{\pgfqpoint{4.009158in}{1.922015in}}%
\pgfpathlineto{\pgfqpoint{14.801389in}{1.922015in}}%
\pgfusepath{stroke}%
\end{pgfscope}%
\begin{pgfscope}%
\definecolor{textcolor}{rgb}{0.333333,0.333333,0.333333}%
\pgfsetstrokecolor{textcolor}%
\pgfsetfillcolor{textcolor}%
\pgftext[x=1.792660in,y=1.868474in,left,base]{\color{textcolor}\sffamily\fontsize{10.000000}{12.000000}\selectfont 1.6.7 Implement Proper Scaling}%
\end{pgfscope}%
\begin{pgfscope}%
\pgfpathrectangle{\pgfqpoint{4.009158in}{0.624448in}}{\pgfqpoint{10.792231in}{7.176941in}} %
\pgfusepath{clip}%
\pgfsetrectcap%
\pgfsetroundjoin%
\pgfsetlinewidth{0.803000pt}%
\definecolor{currentstroke}{rgb}{1.000000,1.000000,1.000000}%
\pgfsetstrokecolor{currentstroke}%
\pgfsetdash{}{0pt}%
\pgfpathmoveto{\pgfqpoint{4.009158in}{1.738743in}}%
\pgfpathlineto{\pgfqpoint{14.801389in}{1.738743in}}%
\pgfusepath{stroke}%
\end{pgfscope}%
\begin{pgfscope}%
\definecolor{textcolor}{rgb}{0.333333,0.333333,0.333333}%
\pgfsetstrokecolor{textcolor}%
\pgfsetfillcolor{textcolor}%
\pgftext[x=1.776180in,y=1.685202in,left,base]{\color{textcolor}\sffamily\fontsize{10.000000}{12.000000}\selectfont 1.6.8 Implement Trigger Modes}%
\end{pgfscope}%
\begin{pgfscope}%
\pgfpathrectangle{\pgfqpoint{4.009158in}{0.624448in}}{\pgfqpoint{10.792231in}{7.176941in}} %
\pgfusepath{clip}%
\pgfsetrectcap%
\pgfsetroundjoin%
\pgfsetlinewidth{0.803000pt}%
\definecolor{currentstroke}{rgb}{1.000000,1.000000,1.000000}%
\pgfsetstrokecolor{currentstroke}%
\pgfsetdash{}{0pt}%
\pgfpathmoveto{\pgfqpoint{4.009158in}{1.555471in}}%
\pgfpathlineto{\pgfqpoint{14.801389in}{1.555471in}}%
\pgfusepath{stroke}%
\end{pgfscope}%
\begin{pgfscope}%
\definecolor{textcolor}{rgb}{0.333333,0.333333,0.333333}%
\pgfsetstrokecolor{textcolor}%
\pgfsetfillcolor{textcolor}%
\pgftext[x=2.309288in,y=1.501930in,left,base]{\color{textcolor}\sffamily\fontsize{10.000000}{12.000000}\selectfont 1.6.9 Implement an FFT}%
\end{pgfscope}%
\begin{pgfscope}%
\pgfpathrectangle{\pgfqpoint{4.009158in}{0.624448in}}{\pgfqpoint{10.792231in}{7.176941in}} %
\pgfusepath{clip}%
\pgfsetrectcap%
\pgfsetroundjoin%
\pgfsetlinewidth{0.803000pt}%
\definecolor{currentstroke}{rgb}{1.000000,1.000000,1.000000}%
\pgfsetstrokecolor{currentstroke}%
\pgfsetdash{}{0pt}%
\pgfpathmoveto{\pgfqpoint{4.009158in}{1.372199in}}%
\pgfpathlineto{\pgfqpoint{14.801389in}{1.372199in}}%
\pgfusepath{stroke}%
\end{pgfscope}%
\begin{pgfscope}%
\definecolor{textcolor}{rgb}{0.333333,0.333333,0.333333}%
\pgfsetstrokecolor{textcolor}%
\pgfsetfillcolor{textcolor}%
\pgftext[x=1.236900in,y=1.318657in,left,base]{\color{textcolor}\sffamily\fontsize{10.000000}{12.000000}\selectfont 1.6.10 Implement Power Measurmenets}%
\end{pgfscope}%
\begin{pgfscope}%
\pgfpathrectangle{\pgfqpoint{4.009158in}{0.624448in}}{\pgfqpoint{10.792231in}{7.176941in}} %
\pgfusepath{clip}%
\pgfsetrectcap%
\pgfsetroundjoin%
\pgfsetlinewidth{0.803000pt}%
\definecolor{currentstroke}{rgb}{1.000000,1.000000,1.000000}%
\pgfsetstrokecolor{currentstroke}%
\pgfsetdash{}{0pt}%
\pgfpathmoveto{\pgfqpoint{4.009158in}{1.188926in}}%
\pgfpathlineto{\pgfqpoint{14.801389in}{1.188926in}}%
\pgfusepath{stroke}%
\end{pgfscope}%
\begin{pgfscope}%
\definecolor{textcolor}{rgb}{0.333333,0.333333,0.333333}%
\pgfsetstrokecolor{textcolor}%
\pgfsetfillcolor{textcolor}%
\pgftext[x=2.037478in,y=1.135385in,left,base]{\color{textcolor}\sffamily\fontsize{10.000000}{12.000000}\selectfont 1.6.14 Document the Scope}%
\end{pgfscope}%
\begin{pgfscope}%
\pgfpathrectangle{\pgfqpoint{4.009158in}{0.624448in}}{\pgfqpoint{10.792231in}{7.176941in}} %
\pgfusepath{clip}%
\pgfsetrectcap%
\pgfsetroundjoin%
\pgfsetlinewidth{0.803000pt}%
\definecolor{currentstroke}{rgb}{1.000000,1.000000,1.000000}%
\pgfsetstrokecolor{currentstroke}%
\pgfsetdash{}{0pt}%
\pgfpathmoveto{\pgfqpoint{4.009158in}{1.005654in}}%
\pgfpathlineto{\pgfqpoint{14.801389in}{1.005654in}}%
\pgfusepath{stroke}%
\end{pgfscope}%
\begin{pgfscope}%
\definecolor{textcolor}{rgb}{0.333333,0.333333,0.333333}%
\pgfsetstrokecolor{textcolor}%
\pgfsetfillcolor{textcolor}%
\pgftext[x=1.979020in,y=0.952113in,left,base]{\color{textcolor}\sffamily\fontsize{10.000000}{12.000000}\selectfont General Project Management}%
\end{pgfscope}%
\begin{pgfscope}%
\pgfpathrectangle{\pgfqpoint{4.009158in}{0.624448in}}{\pgfqpoint{10.792231in}{7.176941in}} %
\pgfusepath{clip}%
\pgfsetbuttcap%
\pgfsetmiterjoin%
\definecolor{currentfill}{rgb}{0.796078,0.294118,0.086275}%
\pgfsetfillcolor{currentfill}%
\pgfsetfillopacity{0.800000}%
\pgfsetlinewidth{0.000000pt}%
\definecolor{currentstroke}{rgb}{0.000000,0.000000,0.000000}%
\pgfsetstrokecolor{currentstroke}%
\pgfsetstrokeopacity{0.800000}%
\pgfsetdash{}{0pt}%
\pgfpathmoveto{\pgfqpoint{4.009158in}{7.429346in}}%
\pgfpathlineto{\pgfqpoint{4.453544in}{7.429346in}}%
\pgfpathlineto{\pgfqpoint{4.453544in}{7.475164in}}%
\pgfpathlineto{\pgfqpoint{4.009158in}{7.475164in}}%
\pgfpathclose%
\pgfusepath{fill}%
\end{pgfscope}%
\begin{pgfscope}%
\pgfpathrectangle{\pgfqpoint{4.009158in}{0.624448in}}{\pgfqpoint{10.792231in}{7.176941in}} %
\pgfusepath{clip}%
\pgfsetbuttcap%
\pgfsetmiterjoin%
\definecolor{currentfill}{rgb}{0.796078,0.294118,0.086275}%
\pgfsetfillcolor{currentfill}%
\pgfsetfillopacity{0.800000}%
\pgfsetlinewidth{0.000000pt}%
\definecolor{currentstroke}{rgb}{0.000000,0.000000,0.000000}%
\pgfsetstrokecolor{currentstroke}%
\pgfsetstrokeopacity{0.800000}%
\pgfsetdash{}{0pt}%
\pgfpathmoveto{\pgfqpoint{4.009158in}{7.246074in}}%
\pgfpathlineto{\pgfqpoint{4.453544in}{7.246074in}}%
\pgfpathlineto{\pgfqpoint{4.453544in}{7.291892in}}%
\pgfpathlineto{\pgfqpoint{4.009158in}{7.291892in}}%
\pgfpathclose%
\pgfusepath{fill}%
\end{pgfscope}%
\begin{pgfscope}%
\pgfpathrectangle{\pgfqpoint{4.009158in}{0.624448in}}{\pgfqpoint{10.792231in}{7.176941in}} %
\pgfusepath{clip}%
\pgfsetbuttcap%
\pgfsetmiterjoin%
\definecolor{currentfill}{rgb}{0.796078,0.294118,0.086275}%
\pgfsetfillcolor{currentfill}%
\pgfsetfillopacity{0.800000}%
\pgfsetlinewidth{0.000000pt}%
\definecolor{currentstroke}{rgb}{0.000000,0.000000,0.000000}%
\pgfsetstrokecolor{currentstroke}%
\pgfsetstrokeopacity{0.800000}%
\pgfsetdash{}{0pt}%
\pgfpathmoveto{\pgfqpoint{4.326577in}{7.062802in}}%
\pgfpathlineto{\pgfqpoint{4.897930in}{7.062802in}}%
\pgfpathlineto{\pgfqpoint{4.897930in}{7.108620in}}%
\pgfpathlineto{\pgfqpoint{4.326577in}{7.108620in}}%
\pgfpathclose%
\pgfusepath{fill}%
\end{pgfscope}%
\begin{pgfscope}%
\pgfpathrectangle{\pgfqpoint{4.009158in}{0.624448in}}{\pgfqpoint{10.792231in}{7.176941in}} %
\pgfusepath{clip}%
\pgfsetbuttcap%
\pgfsetmiterjoin%
\definecolor{currentfill}{rgb}{0.796078,0.294118,0.086275}%
\pgfsetfillcolor{currentfill}%
\pgfsetfillopacity{0.800000}%
\pgfsetlinewidth{0.000000pt}%
\definecolor{currentstroke}{rgb}{0.000000,0.000000,0.000000}%
\pgfsetstrokecolor{currentstroke}%
\pgfsetstrokeopacity{0.800000}%
\pgfsetdash{}{0pt}%
\pgfpathmoveto{\pgfqpoint{4.897930in}{6.879530in}}%
\pgfpathlineto{\pgfqpoint{5.278832in}{6.879530in}}%
\pgfpathlineto{\pgfqpoint{5.278832in}{6.925348in}}%
\pgfpathlineto{\pgfqpoint{4.897930in}{6.925348in}}%
\pgfpathclose%
\pgfusepath{fill}%
\end{pgfscope}%
\begin{pgfscope}%
\pgfpathrectangle{\pgfqpoint{4.009158in}{0.624448in}}{\pgfqpoint{10.792231in}{7.176941in}} %
\pgfusepath{clip}%
\pgfsetbuttcap%
\pgfsetmiterjoin%
\definecolor{currentfill}{rgb}{0.796078,0.294118,0.086275}%
\pgfsetfillcolor{currentfill}%
\pgfsetfillopacity{0.800000}%
\pgfsetlinewidth{0.000000pt}%
\definecolor{currentstroke}{rgb}{0.000000,0.000000,0.000000}%
\pgfsetstrokecolor{currentstroke}%
\pgfsetstrokeopacity{0.800000}%
\pgfsetdash{}{0pt}%
\pgfpathmoveto{\pgfqpoint{4.897930in}{6.696257in}}%
\pgfpathlineto{\pgfqpoint{6.802441in}{6.696257in}}%
\pgfpathlineto{\pgfqpoint{6.802441in}{6.742075in}}%
\pgfpathlineto{\pgfqpoint{4.897930in}{6.742075in}}%
\pgfpathclose%
\pgfusepath{fill}%
\end{pgfscope}%
\begin{pgfscope}%
\pgfpathrectangle{\pgfqpoint{4.009158in}{0.624448in}}{\pgfqpoint{10.792231in}{7.176941in}} %
\pgfusepath{clip}%
\pgfsetbuttcap%
\pgfsetmiterjoin%
\definecolor{currentfill}{rgb}{0.796078,0.294118,0.086275}%
\pgfsetfillcolor{currentfill}%
\pgfsetfillopacity{0.800000}%
\pgfsetlinewidth{0.000000pt}%
\definecolor{currentstroke}{rgb}{0.000000,0.000000,0.000000}%
\pgfsetstrokecolor{currentstroke}%
\pgfsetstrokeopacity{0.800000}%
\pgfsetdash{}{0pt}%
\pgfpathmoveto{\pgfqpoint{6.802441in}{6.512985in}}%
\pgfpathlineto{\pgfqpoint{6.992892in}{6.512985in}}%
\pgfpathlineto{\pgfqpoint{6.992892in}{6.558803in}}%
\pgfpathlineto{\pgfqpoint{6.802441in}{6.558803in}}%
\pgfpathclose%
\pgfusepath{fill}%
\end{pgfscope}%
\begin{pgfscope}%
\pgfpathrectangle{\pgfqpoint{4.009158in}{0.624448in}}{\pgfqpoint{10.792231in}{7.176941in}} %
\pgfusepath{clip}%
\pgfsetbuttcap%
\pgfsetmiterjoin%
\definecolor{currentfill}{rgb}{0.796078,0.294118,0.086275}%
\pgfsetfillcolor{currentfill}%
\pgfsetfillopacity{0.800000}%
\pgfsetlinewidth{0.000000pt}%
\definecolor{currentstroke}{rgb}{0.000000,0.000000,0.000000}%
\pgfsetstrokecolor{currentstroke}%
\pgfsetstrokeopacity{0.800000}%
\pgfsetdash{}{0pt}%
\pgfpathmoveto{\pgfqpoint{6.992892in}{6.329713in}}%
\pgfpathlineto{\pgfqpoint{10.674948in}{6.329713in}}%
\pgfpathlineto{\pgfqpoint{10.674948in}{6.375531in}}%
\pgfpathlineto{\pgfqpoint{6.992892in}{6.375531in}}%
\pgfpathclose%
\pgfusepath{fill}%
\end{pgfscope}%
\begin{pgfscope}%
\pgfpathrectangle{\pgfqpoint{4.009158in}{0.624448in}}{\pgfqpoint{10.792231in}{7.176941in}} %
\pgfusepath{clip}%
\pgfsetbuttcap%
\pgfsetmiterjoin%
\definecolor{currentfill}{rgb}{0.796078,0.294118,0.086275}%
\pgfsetfillcolor{currentfill}%
\pgfsetfillopacity{0.800000}%
\pgfsetlinewidth{0.000000pt}%
\definecolor{currentstroke}{rgb}{0.000000,0.000000,0.000000}%
\pgfsetstrokecolor{currentstroke}%
\pgfsetstrokeopacity{0.800000}%
\pgfsetdash{}{0pt}%
\pgfpathmoveto{\pgfqpoint{10.674948in}{6.146441in}}%
\pgfpathlineto{\pgfqpoint{11.563720in}{6.146441in}}%
\pgfpathlineto{\pgfqpoint{11.563720in}{6.192259in}}%
\pgfpathlineto{\pgfqpoint{10.674948in}{6.192259in}}%
\pgfpathclose%
\pgfusepath{fill}%
\end{pgfscope}%
\begin{pgfscope}%
\pgfpathrectangle{\pgfqpoint{4.009158in}{0.624448in}}{\pgfqpoint{10.792231in}{7.176941in}} %
\pgfusepath{clip}%
\pgfsetbuttcap%
\pgfsetmiterjoin%
\definecolor{currentfill}{rgb}{0.796078,0.294118,0.086275}%
\pgfsetfillcolor{currentfill}%
\pgfsetfillopacity{0.800000}%
\pgfsetlinewidth{0.000000pt}%
\definecolor{currentstroke}{rgb}{0.000000,0.000000,0.000000}%
\pgfsetstrokecolor{currentstroke}%
\pgfsetstrokeopacity{0.800000}%
\pgfsetdash{}{0pt}%
\pgfpathmoveto{\pgfqpoint{13.214296in}{5.963168in}}%
\pgfpathlineto{\pgfqpoint{14.801389in}{5.963168in}}%
\pgfpathlineto{\pgfqpoint{14.801389in}{6.008986in}}%
\pgfpathlineto{\pgfqpoint{13.214296in}{6.008986in}}%
\pgfpathclose%
\pgfusepath{fill}%
\end{pgfscope}%
\begin{pgfscope}%
\pgfpathrectangle{\pgfqpoint{4.009158in}{0.624448in}}{\pgfqpoint{10.792231in}{7.176941in}} %
\pgfusepath{clip}%
\pgfsetbuttcap%
\pgfsetmiterjoin%
\definecolor{currentfill}{rgb}{0.796078,0.294118,0.086275}%
\pgfsetfillcolor{currentfill}%
\pgfsetfillopacity{0.800000}%
\pgfsetlinewidth{0.000000pt}%
\definecolor{currentstroke}{rgb}{0.000000,0.000000,0.000000}%
\pgfsetstrokecolor{currentstroke}%
\pgfsetstrokeopacity{0.800000}%
\pgfsetdash{}{0pt}%
\pgfpathmoveto{\pgfqpoint{7.437278in}{5.779896in}}%
\pgfpathlineto{\pgfqpoint{8.008632in}{5.779896in}}%
\pgfpathlineto{\pgfqpoint{8.008632in}{5.825714in}}%
\pgfpathlineto{\pgfqpoint{7.437278in}{5.825714in}}%
\pgfpathclose%
\pgfusepath{fill}%
\end{pgfscope}%
\begin{pgfscope}%
\pgfpathrectangle{\pgfqpoint{4.009158in}{0.624448in}}{\pgfqpoint{10.792231in}{7.176941in}} %
\pgfusepath{clip}%
\pgfsetbuttcap%
\pgfsetmiterjoin%
\definecolor{currentfill}{rgb}{0.796078,0.294118,0.086275}%
\pgfsetfillcolor{currentfill}%
\pgfsetfillopacity{0.800000}%
\pgfsetlinewidth{0.000000pt}%
\definecolor{currentstroke}{rgb}{0.000000,0.000000,0.000000}%
\pgfsetstrokecolor{currentstroke}%
\pgfsetstrokeopacity{0.800000}%
\pgfsetdash{}{0pt}%
\pgfpathmoveto{\pgfqpoint{9.341790in}{5.596624in}}%
\pgfpathlineto{\pgfqpoint{14.801389in}{5.596624in}}%
\pgfpathlineto{\pgfqpoint{14.801389in}{5.642442in}}%
\pgfpathlineto{\pgfqpoint{9.341790in}{5.642442in}}%
\pgfpathclose%
\pgfusepath{fill}%
\end{pgfscope}%
\begin{pgfscope}%
\pgfpathrectangle{\pgfqpoint{4.009158in}{0.624448in}}{\pgfqpoint{10.792231in}{7.176941in}} %
\pgfusepath{clip}%
\pgfsetbuttcap%
\pgfsetmiterjoin%
\definecolor{currentfill}{rgb}{0.796078,0.294118,0.086275}%
\pgfsetfillcolor{currentfill}%
\pgfsetfillopacity{0.800000}%
\pgfsetlinewidth{0.000000pt}%
\definecolor{currentstroke}{rgb}{0.000000,0.000000,0.000000}%
\pgfsetstrokecolor{currentstroke}%
\pgfsetstrokeopacity{0.800000}%
\pgfsetdash{}{0pt}%
\pgfpathmoveto{\pgfqpoint{5.786702in}{5.413352in}}%
\pgfpathlineto{\pgfqpoint{6.675474in}{5.413352in}}%
\pgfpathlineto{\pgfqpoint{6.675474in}{5.459170in}}%
\pgfpathlineto{\pgfqpoint{5.786702in}{5.459170in}}%
\pgfpathclose%
\pgfusepath{fill}%
\end{pgfscope}%
\begin{pgfscope}%
\pgfpathrectangle{\pgfqpoint{4.009158in}{0.624448in}}{\pgfqpoint{10.792231in}{7.176941in}} %
\pgfusepath{clip}%
\pgfsetbuttcap%
\pgfsetmiterjoin%
\definecolor{currentfill}{rgb}{0.796078,0.294118,0.086275}%
\pgfsetfillcolor{currentfill}%
\pgfsetfillopacity{0.800000}%
\pgfsetlinewidth{0.000000pt}%
\definecolor{currentstroke}{rgb}{0.000000,0.000000,0.000000}%
\pgfsetstrokecolor{currentstroke}%
\pgfsetstrokeopacity{0.800000}%
\pgfsetdash{}{0pt}%
\pgfpathmoveto{\pgfqpoint{6.675474in}{5.230079in}}%
\pgfpathlineto{\pgfqpoint{7.119860in}{5.230079in}}%
\pgfpathlineto{\pgfqpoint{7.119860in}{5.275897in}}%
\pgfpathlineto{\pgfqpoint{6.675474in}{5.275897in}}%
\pgfpathclose%
\pgfusepath{fill}%
\end{pgfscope}%
\begin{pgfscope}%
\pgfpathrectangle{\pgfqpoint{4.009158in}{0.624448in}}{\pgfqpoint{10.792231in}{7.176941in}} %
\pgfusepath{clip}%
\pgfsetbuttcap%
\pgfsetmiterjoin%
\definecolor{currentfill}{rgb}{0.796078,0.294118,0.086275}%
\pgfsetfillcolor{currentfill}%
\pgfsetfillopacity{0.800000}%
\pgfsetlinewidth{0.000000pt}%
\definecolor{currentstroke}{rgb}{0.000000,0.000000,0.000000}%
\pgfsetstrokecolor{currentstroke}%
\pgfsetstrokeopacity{0.800000}%
\pgfsetdash{}{0pt}%
\pgfpathmoveto{\pgfqpoint{7.119860in}{5.046807in}}%
\pgfpathlineto{\pgfqpoint{7.945148in}{5.046807in}}%
\pgfpathlineto{\pgfqpoint{7.945148in}{5.092625in}}%
\pgfpathlineto{\pgfqpoint{7.119860in}{5.092625in}}%
\pgfpathclose%
\pgfusepath{fill}%
\end{pgfscope}%
\begin{pgfscope}%
\pgfpathrectangle{\pgfqpoint{4.009158in}{0.624448in}}{\pgfqpoint{10.792231in}{7.176941in}} %
\pgfusepath{clip}%
\pgfsetbuttcap%
\pgfsetmiterjoin%
\definecolor{currentfill}{rgb}{0.796078,0.294118,0.086275}%
\pgfsetfillcolor{currentfill}%
\pgfsetfillopacity{0.800000}%
\pgfsetlinewidth{0.000000pt}%
\definecolor{currentstroke}{rgb}{0.000000,0.000000,0.000000}%
\pgfsetstrokecolor{currentstroke}%
\pgfsetstrokeopacity{0.800000}%
\pgfsetdash{}{0pt}%
\pgfpathmoveto{\pgfqpoint{7.945148in}{4.863535in}}%
\pgfpathlineto{\pgfqpoint{9.913143in}{4.863535in}}%
\pgfpathlineto{\pgfqpoint{9.913143in}{4.909353in}}%
\pgfpathlineto{\pgfqpoint{7.945148in}{4.909353in}}%
\pgfpathclose%
\pgfusepath{fill}%
\end{pgfscope}%
\begin{pgfscope}%
\pgfpathrectangle{\pgfqpoint{4.009158in}{0.624448in}}{\pgfqpoint{10.792231in}{7.176941in}} %
\pgfusepath{clip}%
\pgfsetbuttcap%
\pgfsetmiterjoin%
\definecolor{currentfill}{rgb}{0.796078,0.294118,0.086275}%
\pgfsetfillcolor{currentfill}%
\pgfsetfillopacity{0.800000}%
\pgfsetlinewidth{0.000000pt}%
\definecolor{currentstroke}{rgb}{0.000000,0.000000,0.000000}%
\pgfsetstrokecolor{currentstroke}%
\pgfsetstrokeopacity{0.800000}%
\pgfsetdash{}{0pt}%
\pgfpathmoveto{\pgfqpoint{13.214296in}{4.680263in}}%
\pgfpathlineto{\pgfqpoint{14.801389in}{4.680263in}}%
\pgfpathlineto{\pgfqpoint{14.801389in}{4.726081in}}%
\pgfpathlineto{\pgfqpoint{13.214296in}{4.726081in}}%
\pgfpathclose%
\pgfusepath{fill}%
\end{pgfscope}%
\begin{pgfscope}%
\pgfpathrectangle{\pgfqpoint{4.009158in}{0.624448in}}{\pgfqpoint{10.792231in}{7.176941in}} %
\pgfusepath{clip}%
\pgfsetbuttcap%
\pgfsetmiterjoin%
\definecolor{currentfill}{rgb}{0.796078,0.294118,0.086275}%
\pgfsetfillcolor{currentfill}%
\pgfsetfillopacity{0.800000}%
\pgfsetlinewidth{0.000000pt}%
\definecolor{currentstroke}{rgb}{0.000000,0.000000,0.000000}%
\pgfsetstrokecolor{currentstroke}%
\pgfsetstrokeopacity{0.800000}%
\pgfsetdash{}{0pt}%
\pgfpathmoveto{\pgfqpoint{5.342316in}{4.496990in}}%
\pgfpathlineto{\pgfqpoint{5.786702in}{4.496990in}}%
\pgfpathlineto{\pgfqpoint{5.786702in}{4.542808in}}%
\pgfpathlineto{\pgfqpoint{5.342316in}{4.542808in}}%
\pgfpathclose%
\pgfusepath{fill}%
\end{pgfscope}%
\begin{pgfscope}%
\pgfpathrectangle{\pgfqpoint{4.009158in}{0.624448in}}{\pgfqpoint{10.792231in}{7.176941in}} %
\pgfusepath{clip}%
\pgfsetbuttcap%
\pgfsetmiterjoin%
\definecolor{currentfill}{rgb}{0.796078,0.294118,0.086275}%
\pgfsetfillcolor{currentfill}%
\pgfsetfillopacity{0.800000}%
\pgfsetlinewidth{0.000000pt}%
\definecolor{currentstroke}{rgb}{0.000000,0.000000,0.000000}%
\pgfsetstrokecolor{currentstroke}%
\pgfsetstrokeopacity{0.800000}%
\pgfsetdash{}{0pt}%
\pgfpathmoveto{\pgfqpoint{5.786702in}{4.313718in}}%
\pgfpathlineto{\pgfqpoint{6.738958in}{4.313718in}}%
\pgfpathlineto{\pgfqpoint{6.738958in}{4.359536in}}%
\pgfpathlineto{\pgfqpoint{5.786702in}{4.359536in}}%
\pgfpathclose%
\pgfusepath{fill}%
\end{pgfscope}%
\begin{pgfscope}%
\pgfpathrectangle{\pgfqpoint{4.009158in}{0.624448in}}{\pgfqpoint{10.792231in}{7.176941in}} %
\pgfusepath{clip}%
\pgfsetbuttcap%
\pgfsetmiterjoin%
\definecolor{currentfill}{rgb}{0.796078,0.294118,0.086275}%
\pgfsetfillcolor{currentfill}%
\pgfsetfillopacity{0.800000}%
\pgfsetlinewidth{0.000000pt}%
\definecolor{currentstroke}{rgb}{0.000000,0.000000,0.000000}%
\pgfsetstrokecolor{currentstroke}%
\pgfsetstrokeopacity{0.800000}%
\pgfsetdash{}{0pt}%
\pgfpathmoveto{\pgfqpoint{6.738958in}{4.130446in}}%
\pgfpathlineto{\pgfqpoint{7.627730in}{4.130446in}}%
\pgfpathlineto{\pgfqpoint{7.627730in}{4.176264in}}%
\pgfpathlineto{\pgfqpoint{6.738958in}{4.176264in}}%
\pgfpathclose%
\pgfusepath{fill}%
\end{pgfscope}%
\begin{pgfscope}%
\pgfpathrectangle{\pgfqpoint{4.009158in}{0.624448in}}{\pgfqpoint{10.792231in}{7.176941in}} %
\pgfusepath{clip}%
\pgfsetbuttcap%
\pgfsetmiterjoin%
\definecolor{currentfill}{rgb}{0.796078,0.294118,0.086275}%
\pgfsetfillcolor{currentfill}%
\pgfsetfillopacity{0.800000}%
\pgfsetlinewidth{0.000000pt}%
\definecolor{currentstroke}{rgb}{0.000000,0.000000,0.000000}%
\pgfsetstrokecolor{currentstroke}%
\pgfsetstrokeopacity{0.800000}%
\pgfsetdash{}{0pt}%
\pgfpathmoveto{\pgfqpoint{7.627730in}{3.947174in}}%
\pgfpathlineto{\pgfqpoint{10.103594in}{3.947174in}}%
\pgfpathlineto{\pgfqpoint{10.103594in}{3.992992in}}%
\pgfpathlineto{\pgfqpoint{7.627730in}{3.992992in}}%
\pgfpathclose%
\pgfusepath{fill}%
\end{pgfscope}%
\begin{pgfscope}%
\pgfpathrectangle{\pgfqpoint{4.009158in}{0.624448in}}{\pgfqpoint{10.792231in}{7.176941in}} %
\pgfusepath{clip}%
\pgfsetbuttcap%
\pgfsetmiterjoin%
\definecolor{currentfill}{rgb}{0.796078,0.294118,0.086275}%
\pgfsetfillcolor{currentfill}%
\pgfsetfillopacity{0.800000}%
\pgfsetlinewidth{0.000000pt}%
\definecolor{currentstroke}{rgb}{0.000000,0.000000,0.000000}%
\pgfsetstrokecolor{currentstroke}%
\pgfsetstrokeopacity{0.800000}%
\pgfsetdash{}{0pt}%
\pgfpathmoveto{\pgfqpoint{10.103594in}{3.763901in}}%
\pgfpathlineto{\pgfqpoint{10.547980in}{3.763901in}}%
\pgfpathlineto{\pgfqpoint{10.547980in}{3.809720in}}%
\pgfpathlineto{\pgfqpoint{10.103594in}{3.809720in}}%
\pgfpathclose%
\pgfusepath{fill}%
\end{pgfscope}%
\begin{pgfscope}%
\pgfpathrectangle{\pgfqpoint{4.009158in}{0.624448in}}{\pgfqpoint{10.792231in}{7.176941in}} %
\pgfusepath{clip}%
\pgfsetbuttcap%
\pgfsetmiterjoin%
\definecolor{currentfill}{rgb}{0.796078,0.294118,0.086275}%
\pgfsetfillcolor{currentfill}%
\pgfsetfillopacity{0.800000}%
\pgfsetlinewidth{0.000000pt}%
\definecolor{currentstroke}{rgb}{0.000000,0.000000,0.000000}%
\pgfsetstrokecolor{currentstroke}%
\pgfsetstrokeopacity{0.800000}%
\pgfsetdash{}{0pt}%
\pgfpathmoveto{\pgfqpoint{13.214296in}{3.580629in}}%
\pgfpathlineto{\pgfqpoint{14.801389in}{3.580629in}}%
\pgfpathlineto{\pgfqpoint{14.801389in}{3.626447in}}%
\pgfpathlineto{\pgfqpoint{13.214296in}{3.626447in}}%
\pgfpathclose%
\pgfusepath{fill}%
\end{pgfscope}%
\begin{pgfscope}%
\pgfpathrectangle{\pgfqpoint{4.009158in}{0.624448in}}{\pgfqpoint{10.792231in}{7.176941in}} %
\pgfusepath{clip}%
\pgfsetbuttcap%
\pgfsetmiterjoin%
\definecolor{currentfill}{rgb}{0.796078,0.294118,0.086275}%
\pgfsetfillcolor{currentfill}%
\pgfsetfillopacity{0.800000}%
\pgfsetlinewidth{0.000000pt}%
\definecolor{currentstroke}{rgb}{0.000000,0.000000,0.000000}%
\pgfsetstrokecolor{currentstroke}%
\pgfsetstrokeopacity{0.800000}%
\pgfsetdash{}{0pt}%
\pgfpathmoveto{\pgfqpoint{13.912617in}{3.397357in}}%
\pgfpathlineto{\pgfqpoint{14.357003in}{3.397357in}}%
\pgfpathlineto{\pgfqpoint{14.357003in}{3.443175in}}%
\pgfpathlineto{\pgfqpoint{13.912617in}{3.443175in}}%
\pgfpathclose%
\pgfusepath{fill}%
\end{pgfscope}%
\begin{pgfscope}%
\pgfpathrectangle{\pgfqpoint{4.009158in}{0.624448in}}{\pgfqpoint{10.792231in}{7.176941in}} %
\pgfusepath{clip}%
\pgfsetbuttcap%
\pgfsetmiterjoin%
\definecolor{currentfill}{rgb}{0.796078,0.294118,0.086275}%
\pgfsetfillcolor{currentfill}%
\pgfsetfillopacity{0.800000}%
\pgfsetlinewidth{0.000000pt}%
\definecolor{currentstroke}{rgb}{0.000000,0.000000,0.000000}%
\pgfsetstrokecolor{currentstroke}%
\pgfsetstrokeopacity{0.800000}%
\pgfsetdash{}{0pt}%
\pgfpathmoveto{\pgfqpoint{13.912617in}{3.214085in}}%
\pgfpathlineto{\pgfqpoint{14.801389in}{3.214085in}}%
\pgfpathlineto{\pgfqpoint{14.801389in}{3.259903in}}%
\pgfpathlineto{\pgfqpoint{13.912617in}{3.259903in}}%
\pgfpathclose%
\pgfusepath{fill}%
\end{pgfscope}%
\begin{pgfscope}%
\pgfpathrectangle{\pgfqpoint{4.009158in}{0.624448in}}{\pgfqpoint{10.792231in}{7.176941in}} %
\pgfusepath{clip}%
\pgfsetbuttcap%
\pgfsetmiterjoin%
\definecolor{currentfill}{rgb}{0.796078,0.294118,0.086275}%
\pgfsetfillcolor{currentfill}%
\pgfsetfillopacity{0.800000}%
\pgfsetlinewidth{0.000000pt}%
\definecolor{currentstroke}{rgb}{0.000000,0.000000,0.000000}%
\pgfsetstrokecolor{currentstroke}%
\pgfsetstrokeopacity{0.800000}%
\pgfsetdash{}{0pt}%
\pgfpathmoveto{\pgfqpoint{5.278832in}{3.030812in}}%
\pgfpathlineto{\pgfqpoint{6.231088in}{3.030812in}}%
\pgfpathlineto{\pgfqpoint{6.231088in}{3.076631in}}%
\pgfpathlineto{\pgfqpoint{5.278832in}{3.076631in}}%
\pgfpathclose%
\pgfusepath{fill}%
\end{pgfscope}%
\begin{pgfscope}%
\pgfpathrectangle{\pgfqpoint{4.009158in}{0.624448in}}{\pgfqpoint{10.792231in}{7.176941in}} %
\pgfusepath{clip}%
\pgfsetbuttcap%
\pgfsetmiterjoin%
\definecolor{currentfill}{rgb}{0.796078,0.294118,0.086275}%
\pgfsetfillcolor{currentfill}%
\pgfsetfillopacity{0.800000}%
\pgfsetlinewidth{0.000000pt}%
\definecolor{currentstroke}{rgb}{0.000000,0.000000,0.000000}%
\pgfsetstrokecolor{currentstroke}%
\pgfsetstrokeopacity{0.800000}%
\pgfsetdash{}{0pt}%
\pgfpathmoveto{\pgfqpoint{6.675474in}{2.847540in}}%
\pgfpathlineto{\pgfqpoint{7.119860in}{2.847540in}}%
\pgfpathlineto{\pgfqpoint{7.119860in}{2.893358in}}%
\pgfpathlineto{\pgfqpoint{6.675474in}{2.893358in}}%
\pgfpathclose%
\pgfusepath{fill}%
\end{pgfscope}%
\begin{pgfscope}%
\pgfpathrectangle{\pgfqpoint{4.009158in}{0.624448in}}{\pgfqpoint{10.792231in}{7.176941in}} %
\pgfusepath{clip}%
\pgfsetbuttcap%
\pgfsetmiterjoin%
\definecolor{currentfill}{rgb}{0.796078,0.294118,0.086275}%
\pgfsetfillcolor{currentfill}%
\pgfsetfillopacity{0.800000}%
\pgfsetlinewidth{0.000000pt}%
\definecolor{currentstroke}{rgb}{0.000000,0.000000,0.000000}%
\pgfsetstrokecolor{currentstroke}%
\pgfsetstrokeopacity{0.800000}%
\pgfsetdash{}{0pt}%
\pgfpathmoveto{\pgfqpoint{7.119860in}{2.664268in}}%
\pgfpathlineto{\pgfqpoint{7.564246in}{2.664268in}}%
\pgfpathlineto{\pgfqpoint{7.564246in}{2.710086in}}%
\pgfpathlineto{\pgfqpoint{7.119860in}{2.710086in}}%
\pgfpathclose%
\pgfusepath{fill}%
\end{pgfscope}%
\begin{pgfscope}%
\pgfpathrectangle{\pgfqpoint{4.009158in}{0.624448in}}{\pgfqpoint{10.792231in}{7.176941in}} %
\pgfusepath{clip}%
\pgfsetbuttcap%
\pgfsetmiterjoin%
\definecolor{currentfill}{rgb}{0.796078,0.294118,0.086275}%
\pgfsetfillcolor{currentfill}%
\pgfsetfillopacity{0.800000}%
\pgfsetlinewidth{0.000000pt}%
\definecolor{currentstroke}{rgb}{0.000000,0.000000,0.000000}%
\pgfsetstrokecolor{currentstroke}%
\pgfsetstrokeopacity{0.800000}%
\pgfsetdash{}{0pt}%
\pgfpathmoveto{\pgfqpoint{7.564246in}{2.480996in}}%
\pgfpathlineto{\pgfqpoint{8.008632in}{2.480996in}}%
\pgfpathlineto{\pgfqpoint{8.008632in}{2.526814in}}%
\pgfpathlineto{\pgfqpoint{7.564246in}{2.526814in}}%
\pgfpathclose%
\pgfusepath{fill}%
\end{pgfscope}%
\begin{pgfscope}%
\pgfpathrectangle{\pgfqpoint{4.009158in}{0.624448in}}{\pgfqpoint{10.792231in}{7.176941in}} %
\pgfusepath{clip}%
\pgfsetbuttcap%
\pgfsetmiterjoin%
\definecolor{currentfill}{rgb}{0.796078,0.294118,0.086275}%
\pgfsetfillcolor{currentfill}%
\pgfsetfillopacity{0.800000}%
\pgfsetlinewidth{0.000000pt}%
\definecolor{currentstroke}{rgb}{0.000000,0.000000,0.000000}%
\pgfsetstrokecolor{currentstroke}%
\pgfsetstrokeopacity{0.800000}%
\pgfsetdash{}{0pt}%
\pgfpathmoveto{\pgfqpoint{8.008632in}{2.297724in}}%
\pgfpathlineto{\pgfqpoint{8.453018in}{2.297724in}}%
\pgfpathlineto{\pgfqpoint{8.453018in}{2.343542in}}%
\pgfpathlineto{\pgfqpoint{8.008632in}{2.343542in}}%
\pgfpathclose%
\pgfusepath{fill}%
\end{pgfscope}%
\begin{pgfscope}%
\pgfpathrectangle{\pgfqpoint{4.009158in}{0.624448in}}{\pgfqpoint{10.792231in}{7.176941in}} %
\pgfusepath{clip}%
\pgfsetbuttcap%
\pgfsetmiterjoin%
\definecolor{currentfill}{rgb}{0.796078,0.294118,0.086275}%
\pgfsetfillcolor{currentfill}%
\pgfsetfillopacity{0.800000}%
\pgfsetlinewidth{0.000000pt}%
\definecolor{currentstroke}{rgb}{0.000000,0.000000,0.000000}%
\pgfsetstrokecolor{currentstroke}%
\pgfsetstrokeopacity{0.800000}%
\pgfsetdash{}{0pt}%
\pgfpathmoveto{\pgfqpoint{8.453018in}{2.114451in}}%
\pgfpathlineto{\pgfqpoint{8.897404in}{2.114451in}}%
\pgfpathlineto{\pgfqpoint{8.897404in}{2.160269in}}%
\pgfpathlineto{\pgfqpoint{8.453018in}{2.160269in}}%
\pgfpathclose%
\pgfusepath{fill}%
\end{pgfscope}%
\begin{pgfscope}%
\pgfpathrectangle{\pgfqpoint{4.009158in}{0.624448in}}{\pgfqpoint{10.792231in}{7.176941in}} %
\pgfusepath{clip}%
\pgfsetbuttcap%
\pgfsetmiterjoin%
\definecolor{currentfill}{rgb}{0.796078,0.294118,0.086275}%
\pgfsetfillcolor{currentfill}%
\pgfsetfillopacity{0.800000}%
\pgfsetlinewidth{0.000000pt}%
\definecolor{currentstroke}{rgb}{0.000000,0.000000,0.000000}%
\pgfsetstrokecolor{currentstroke}%
\pgfsetstrokeopacity{0.800000}%
\pgfsetdash{}{0pt}%
\pgfpathmoveto{\pgfqpoint{8.897404in}{1.931179in}}%
\pgfpathlineto{\pgfqpoint{9.341790in}{1.931179in}}%
\pgfpathlineto{\pgfqpoint{9.341790in}{1.976997in}}%
\pgfpathlineto{\pgfqpoint{8.897404in}{1.976997in}}%
\pgfpathclose%
\pgfusepath{fill}%
\end{pgfscope}%
\begin{pgfscope}%
\pgfpathrectangle{\pgfqpoint{4.009158in}{0.624448in}}{\pgfqpoint{10.792231in}{7.176941in}} %
\pgfusepath{clip}%
\pgfsetbuttcap%
\pgfsetmiterjoin%
\definecolor{currentfill}{rgb}{0.796078,0.294118,0.086275}%
\pgfsetfillcolor{currentfill}%
\pgfsetfillopacity{0.800000}%
\pgfsetlinewidth{0.000000pt}%
\definecolor{currentstroke}{rgb}{0.000000,0.000000,0.000000}%
\pgfsetstrokecolor{currentstroke}%
\pgfsetstrokeopacity{0.800000}%
\pgfsetdash{}{0pt}%
\pgfpathmoveto{\pgfqpoint{9.341790in}{1.747907in}}%
\pgfpathlineto{\pgfqpoint{9.849659in}{1.747907in}}%
\pgfpathlineto{\pgfqpoint{9.849659in}{1.793725in}}%
\pgfpathlineto{\pgfqpoint{9.341790in}{1.793725in}}%
\pgfpathclose%
\pgfusepath{fill}%
\end{pgfscope}%
\begin{pgfscope}%
\pgfpathrectangle{\pgfqpoint{4.009158in}{0.624448in}}{\pgfqpoint{10.792231in}{7.176941in}} %
\pgfusepath{clip}%
\pgfsetbuttcap%
\pgfsetmiterjoin%
\definecolor{currentfill}{rgb}{0.796078,0.294118,0.086275}%
\pgfsetfillcolor{currentfill}%
\pgfsetfillopacity{0.800000}%
\pgfsetlinewidth{0.000000pt}%
\definecolor{currentstroke}{rgb}{0.000000,0.000000,0.000000}%
\pgfsetstrokecolor{currentstroke}%
\pgfsetstrokeopacity{0.800000}%
\pgfsetdash{}{0pt}%
\pgfpathmoveto{\pgfqpoint{9.849659in}{1.564635in}}%
\pgfpathlineto{\pgfqpoint{10.294045in}{1.564635in}}%
\pgfpathlineto{\pgfqpoint{10.294045in}{1.610453in}}%
\pgfpathlineto{\pgfqpoint{9.849659in}{1.610453in}}%
\pgfpathclose%
\pgfusepath{fill}%
\end{pgfscope}%
\begin{pgfscope}%
\pgfpathrectangle{\pgfqpoint{4.009158in}{0.624448in}}{\pgfqpoint{10.792231in}{7.176941in}} %
\pgfusepath{clip}%
\pgfsetbuttcap%
\pgfsetmiterjoin%
\definecolor{currentfill}{rgb}{0.796078,0.294118,0.086275}%
\pgfsetfillcolor{currentfill}%
\pgfsetfillopacity{0.800000}%
\pgfsetlinewidth{0.000000pt}%
\definecolor{currentstroke}{rgb}{0.000000,0.000000,0.000000}%
\pgfsetstrokecolor{currentstroke}%
\pgfsetstrokeopacity{0.800000}%
\pgfsetdash{}{0pt}%
\pgfpathmoveto{\pgfqpoint{10.294045in}{1.381362in}}%
\pgfpathlineto{\pgfqpoint{11.373268in}{1.381362in}}%
\pgfpathlineto{\pgfqpoint{11.373268in}{1.427180in}}%
\pgfpathlineto{\pgfqpoint{10.294045in}{1.427180in}}%
\pgfpathclose%
\pgfusepath{fill}%
\end{pgfscope}%
\begin{pgfscope}%
\pgfpathrectangle{\pgfqpoint{4.009158in}{0.624448in}}{\pgfqpoint{10.792231in}{7.176941in}} %
\pgfusepath{clip}%
\pgfsetbuttcap%
\pgfsetmiterjoin%
\definecolor{currentfill}{rgb}{0.796078,0.294118,0.086275}%
\pgfsetfillcolor{currentfill}%
\pgfsetfillopacity{0.800000}%
\pgfsetlinewidth{0.000000pt}%
\definecolor{currentstroke}{rgb}{0.000000,0.000000,0.000000}%
\pgfsetstrokecolor{currentstroke}%
\pgfsetstrokeopacity{0.800000}%
\pgfsetdash{}{0pt}%
\pgfpathmoveto{\pgfqpoint{13.214296in}{1.198090in}}%
\pgfpathlineto{\pgfqpoint{14.801389in}{1.198090in}}%
\pgfpathlineto{\pgfqpoint{14.801389in}{1.243908in}}%
\pgfpathlineto{\pgfqpoint{13.214296in}{1.243908in}}%
\pgfpathclose%
\pgfusepath{fill}%
\end{pgfscope}%
\begin{pgfscope}%
\pgfpathrectangle{\pgfqpoint{4.009158in}{0.624448in}}{\pgfqpoint{10.792231in}{7.176941in}} %
\pgfusepath{clip}%
\pgfsetbuttcap%
\pgfsetmiterjoin%
\definecolor{currentfill}{rgb}{0.796078,0.294118,0.086275}%
\pgfsetfillcolor{currentfill}%
\pgfsetfillopacity{0.800000}%
\pgfsetlinewidth{0.000000pt}%
\definecolor{currentstroke}{rgb}{0.000000,0.000000,0.000000}%
\pgfsetstrokecolor{currentstroke}%
\pgfsetstrokeopacity{0.800000}%
\pgfsetdash{}{0pt}%
\pgfpathmoveto{\pgfqpoint{4.009158in}{1.014818in}}%
\pgfpathlineto{\pgfqpoint{14.801389in}{1.014818in}}%
\pgfpathlineto{\pgfqpoint{14.801389in}{1.060636in}}%
\pgfpathlineto{\pgfqpoint{4.009158in}{1.060636in}}%
\pgfpathclose%
\pgfusepath{fill}%
\end{pgfscope}%
\begin{pgfscope}%
\pgfpathrectangle{\pgfqpoint{4.009158in}{0.624448in}}{\pgfqpoint{10.792231in}{7.176941in}} %
\pgfusepath{clip}%
\pgfsetbuttcap%
\pgfsetmiterjoin%
\definecolor{currentfill}{rgb}{0.149020,0.545098,0.823529}%
\pgfsetfillcolor{currentfill}%
\pgfsetfillopacity{0.800000}%
\pgfsetlinewidth{0.000000pt}%
\definecolor{currentstroke}{rgb}{0.000000,0.000000,0.000000}%
\pgfsetstrokecolor{currentstroke}%
\pgfsetstrokeopacity{0.800000}%
\pgfsetdash{}{0pt}%
\pgfpathmoveto{\pgfqpoint{4.009158in}{7.365201in}}%
\pgfpathlineto{\pgfqpoint{4.580511in}{7.365201in}}%
\pgfpathlineto{\pgfqpoint{4.580511in}{7.411019in}}%
\pgfpathlineto{\pgfqpoint{4.009158in}{7.411019in}}%
\pgfpathclose%
\pgfusepath{fill}%
\end{pgfscope}%
\begin{pgfscope}%
\pgfpathrectangle{\pgfqpoint{4.009158in}{0.624448in}}{\pgfqpoint{10.792231in}{7.176941in}} %
\pgfusepath{clip}%
\pgfsetbuttcap%
\pgfsetmiterjoin%
\definecolor{currentfill}{rgb}{0.149020,0.545098,0.823529}%
\pgfsetfillcolor{currentfill}%
\pgfsetfillopacity{0.800000}%
\pgfsetlinewidth{0.000000pt}%
\definecolor{currentstroke}{rgb}{0.000000,0.000000,0.000000}%
\pgfsetstrokecolor{currentstroke}%
\pgfsetstrokeopacity{0.800000}%
\pgfsetdash{}{0pt}%
\pgfpathmoveto{\pgfqpoint{4.009158in}{7.181929in}}%
\pgfpathlineto{\pgfqpoint{4.072642in}{7.181929in}}%
\pgfpathlineto{\pgfqpoint{4.072642in}{7.227747in}}%
\pgfpathlineto{\pgfqpoint{4.009158in}{7.227747in}}%
\pgfpathclose%
\pgfusepath{fill}%
\end{pgfscope}%
\begin{pgfscope}%
\pgfpathrectangle{\pgfqpoint{4.009158in}{0.624448in}}{\pgfqpoint{10.792231in}{7.176941in}} %
\pgfusepath{clip}%
\pgfsetbuttcap%
\pgfsetmiterjoin%
\definecolor{currentfill}{rgb}{0.149020,0.545098,0.823529}%
\pgfsetfillcolor{currentfill}%
\pgfsetfillopacity{0.800000}%
\pgfsetlinewidth{0.000000pt}%
\definecolor{currentstroke}{rgb}{0.000000,0.000000,0.000000}%
\pgfsetstrokecolor{currentstroke}%
\pgfsetstrokeopacity{0.800000}%
\pgfsetdash{}{0pt}%
\pgfpathmoveto{\pgfqpoint{4.453544in}{6.998656in}}%
\pgfpathlineto{\pgfqpoint{4.897930in}{6.998656in}}%
\pgfpathlineto{\pgfqpoint{4.897930in}{7.044475in}}%
\pgfpathlineto{\pgfqpoint{4.453544in}{7.044475in}}%
\pgfpathclose%
\pgfusepath{fill}%
\end{pgfscope}%
\begin{pgfscope}%
\pgfpathrectangle{\pgfqpoint{4.009158in}{0.624448in}}{\pgfqpoint{10.792231in}{7.176941in}} %
\pgfusepath{clip}%
\pgfsetbuttcap%
\pgfsetmiterjoin%
\definecolor{currentfill}{rgb}{0.149020,0.545098,0.823529}%
\pgfsetfillcolor{currentfill}%
\pgfsetfillopacity{0.800000}%
\pgfsetlinewidth{0.000000pt}%
\definecolor{currentstroke}{rgb}{0.000000,0.000000,0.000000}%
\pgfsetstrokecolor{currentstroke}%
\pgfsetstrokeopacity{0.800000}%
\pgfsetdash{}{0pt}%
\pgfpathmoveto{\pgfqpoint{4.453544in}{6.815384in}}%
\pgfpathlineto{\pgfqpoint{5.786702in}{6.815384in}}%
\pgfpathlineto{\pgfqpoint{5.786702in}{6.861202in}}%
\pgfpathlineto{\pgfqpoint{4.453544in}{6.861202in}}%
\pgfpathclose%
\pgfusepath{fill}%
\end{pgfscope}%
\begin{pgfscope}%
\pgfpathrectangle{\pgfqpoint{4.009158in}{0.624448in}}{\pgfqpoint{10.792231in}{7.176941in}} %
\pgfusepath{clip}%
\pgfsetbuttcap%
\pgfsetmiterjoin%
\definecolor{currentfill}{rgb}{0.149020,0.545098,0.823529}%
\pgfsetfillcolor{currentfill}%
\pgfsetfillopacity{0.800000}%
\pgfsetlinewidth{0.000000pt}%
\definecolor{currentstroke}{rgb}{0.000000,0.000000,0.000000}%
\pgfsetstrokecolor{currentstroke}%
\pgfsetstrokeopacity{0.800000}%
\pgfsetdash{}{0pt}%
\pgfpathmoveto{\pgfqpoint{4.897930in}{6.632112in}}%
\pgfpathlineto{\pgfqpoint{7.119860in}{6.632112in}}%
\pgfpathlineto{\pgfqpoint{7.119860in}{6.677930in}}%
\pgfpathlineto{\pgfqpoint{4.897930in}{6.677930in}}%
\pgfpathclose%
\pgfusepath{fill}%
\end{pgfscope}%
\begin{pgfscope}%
\pgfpathrectangle{\pgfqpoint{4.009158in}{0.624448in}}{\pgfqpoint{10.792231in}{7.176941in}} %
\pgfusepath{clip}%
\pgfsetbuttcap%
\pgfsetmiterjoin%
\definecolor{currentfill}{rgb}{0.149020,0.545098,0.823529}%
\pgfsetfillcolor{currentfill}%
\pgfsetfillopacity{0.800000}%
\pgfsetlinewidth{0.000000pt}%
\definecolor{currentstroke}{rgb}{0.000000,0.000000,0.000000}%
\pgfsetstrokecolor{currentstroke}%
\pgfsetstrokeopacity{0.800000}%
\pgfsetdash{}{0pt}%
\pgfpathmoveto{\pgfqpoint{7.119860in}{6.448840in}}%
\pgfpathlineto{\pgfqpoint{7.310311in}{6.448840in}}%
\pgfpathlineto{\pgfqpoint{7.310311in}{6.494658in}}%
\pgfpathlineto{\pgfqpoint{7.119860in}{6.494658in}}%
\pgfpathclose%
\pgfusepath{fill}%
\end{pgfscope}%
\begin{pgfscope}%
\pgfpathrectangle{\pgfqpoint{4.009158in}{0.624448in}}{\pgfqpoint{10.792231in}{7.176941in}} %
\pgfusepath{clip}%
\pgfsetbuttcap%
\pgfsetmiterjoin%
\definecolor{currentfill}{rgb}{0.149020,0.545098,0.823529}%
\pgfsetfillcolor{currentfill}%
\pgfsetfillopacity{0.800000}%
\pgfsetlinewidth{0.000000pt}%
\definecolor{currentstroke}{rgb}{0.000000,0.000000,0.000000}%
\pgfsetstrokecolor{currentstroke}%
\pgfsetstrokeopacity{0.800000}%
\pgfsetdash{}{0pt}%
\pgfpathmoveto{\pgfqpoint{8.008632in}{6.265568in}}%
\pgfpathlineto{\pgfqpoint{12.960361in}{6.265568in}}%
\pgfpathlineto{\pgfqpoint{12.960361in}{6.311386in}}%
\pgfpathlineto{\pgfqpoint{8.008632in}{6.311386in}}%
\pgfpathclose%
\pgfusepath{fill}%
\end{pgfscope}%
\begin{pgfscope}%
\pgfpathrectangle{\pgfqpoint{4.009158in}{0.624448in}}{\pgfqpoint{10.792231in}{7.176941in}} %
\pgfusepath{clip}%
\pgfsetbuttcap%
\pgfsetmiterjoin%
\definecolor{currentfill}{rgb}{0.149020,0.545098,0.823529}%
\pgfsetfillcolor{currentfill}%
\pgfsetfillopacity{0.800000}%
\pgfsetlinewidth{0.000000pt}%
\definecolor{currentstroke}{rgb}{0.000000,0.000000,0.000000}%
\pgfsetstrokecolor{currentstroke}%
\pgfsetstrokeopacity{0.800000}%
\pgfsetdash{}{0pt}%
\pgfpathmoveto{\pgfqpoint{12.960361in}{6.082295in}}%
\pgfpathlineto{\pgfqpoint{13.976101in}{6.082295in}}%
\pgfpathlineto{\pgfqpoint{13.976101in}{6.128113in}}%
\pgfpathlineto{\pgfqpoint{12.960361in}{6.128113in}}%
\pgfpathclose%
\pgfusepath{fill}%
\end{pgfscope}%
\begin{pgfscope}%
\pgfpathrectangle{\pgfqpoint{4.009158in}{0.624448in}}{\pgfqpoint{10.792231in}{7.176941in}} %
\pgfusepath{clip}%
\pgfsetbuttcap%
\pgfsetmiterjoin%
\definecolor{currentfill}{rgb}{0.149020,0.545098,0.823529}%
\pgfsetfillcolor{currentfill}%
\pgfsetfillopacity{0.800000}%
\pgfsetlinewidth{0.000000pt}%
\definecolor{currentstroke}{rgb}{0.000000,0.000000,0.000000}%
\pgfsetstrokecolor{currentstroke}%
\pgfsetstrokeopacity{0.800000}%
\pgfsetdash{}{0pt}%
\pgfpathmoveto{\pgfqpoint{12.579459in}{5.899023in}}%
\pgfpathlineto{\pgfqpoint{14.483970in}{5.899023in}}%
\pgfpathlineto{\pgfqpoint{14.483970in}{5.944841in}}%
\pgfpathlineto{\pgfqpoint{12.579459in}{5.944841in}}%
\pgfpathclose%
\pgfusepath{fill}%
\end{pgfscope}%
\begin{pgfscope}%
\pgfpathrectangle{\pgfqpoint{4.009158in}{0.624448in}}{\pgfqpoint{10.792231in}{7.176941in}} %
\pgfusepath{clip}%
\pgfsetbuttcap%
\pgfsetmiterjoin%
\definecolor{currentfill}{rgb}{0.149020,0.545098,0.823529}%
\pgfsetfillcolor{currentfill}%
\pgfsetfillopacity{0.800000}%
\pgfsetlinewidth{0.000000pt}%
\definecolor{currentstroke}{rgb}{0.000000,0.000000,0.000000}%
\pgfsetstrokecolor{currentstroke}%
\pgfsetstrokeopacity{0.800000}%
\pgfsetdash{}{0pt}%
\pgfpathmoveto{\pgfqpoint{5.151865in}{5.715751in}}%
\pgfpathlineto{\pgfqpoint{5.215348in}{5.715751in}}%
\pgfpathlineto{\pgfqpoint{5.215348in}{5.761569in}}%
\pgfpathlineto{\pgfqpoint{5.151865in}{5.761569in}}%
\pgfpathclose%
\pgfusepath{fill}%
\end{pgfscope}%
\begin{pgfscope}%
\pgfpathrectangle{\pgfqpoint{4.009158in}{0.624448in}}{\pgfqpoint{10.792231in}{7.176941in}} %
\pgfusepath{clip}%
\pgfsetbuttcap%
\pgfsetmiterjoin%
\definecolor{currentfill}{rgb}{0.149020,0.545098,0.823529}%
\pgfsetfillcolor{currentfill}%
\pgfsetfillopacity{0.800000}%
\pgfsetlinewidth{0.000000pt}%
\definecolor{currentstroke}{rgb}{0.000000,0.000000,0.000000}%
\pgfsetstrokecolor{currentstroke}%
\pgfsetstrokeopacity{0.800000}%
\pgfsetdash{}{0pt}%
\pgfpathmoveto{\pgfqpoint{6.358055in}{5.532479in}}%
\pgfpathlineto{\pgfqpoint{14.737905in}{5.532479in}}%
\pgfpathlineto{\pgfqpoint{14.737905in}{5.578297in}}%
\pgfpathlineto{\pgfqpoint{6.358055in}{5.578297in}}%
\pgfpathclose%
\pgfusepath{fill}%
\end{pgfscope}%
\begin{pgfscope}%
\pgfpathrectangle{\pgfqpoint{4.009158in}{0.624448in}}{\pgfqpoint{10.792231in}{7.176941in}} %
\pgfusepath{clip}%
\pgfsetbuttcap%
\pgfsetmiterjoin%
\definecolor{currentfill}{rgb}{0.149020,0.545098,0.823529}%
\pgfsetfillcolor{currentfill}%
\pgfsetfillopacity{0.800000}%
\pgfsetlinewidth{0.000000pt}%
\definecolor{currentstroke}{rgb}{0.000000,0.000000,0.000000}%
\pgfsetstrokecolor{currentstroke}%
\pgfsetstrokeopacity{0.800000}%
\pgfsetdash{}{0pt}%
\pgfpathmoveto{\pgfqpoint{6.802441in}{5.349206in}}%
\pgfpathlineto{\pgfqpoint{8.135599in}{5.349206in}}%
\pgfpathlineto{\pgfqpoint{8.135599in}{5.395024in}}%
\pgfpathlineto{\pgfqpoint{6.802441in}{5.395024in}}%
\pgfpathclose%
\pgfusepath{fill}%
\end{pgfscope}%
\begin{pgfscope}%
\pgfpathrectangle{\pgfqpoint{4.009158in}{0.624448in}}{\pgfqpoint{10.792231in}{7.176941in}} %
\pgfusepath{clip}%
\pgfsetbuttcap%
\pgfsetmiterjoin%
\definecolor{currentfill}{rgb}{0.149020,0.545098,0.823529}%
\pgfsetfillcolor{currentfill}%
\pgfsetfillopacity{0.800000}%
\pgfsetlinewidth{0.000000pt}%
\definecolor{currentstroke}{rgb}{0.000000,0.000000,0.000000}%
\pgfsetstrokecolor{currentstroke}%
\pgfsetstrokeopacity{0.800000}%
\pgfsetdash{}{0pt}%
\pgfpathmoveto{\pgfqpoint{7.500762in}{5.165934in}}%
\pgfpathlineto{\pgfqpoint{7.945148in}{5.165934in}}%
\pgfpathlineto{\pgfqpoint{7.945148in}{5.211752in}}%
\pgfpathlineto{\pgfqpoint{7.500762in}{5.211752in}}%
\pgfpathclose%
\pgfusepath{fill}%
\end{pgfscope}%
\begin{pgfscope}%
\pgfpathrectangle{\pgfqpoint{4.009158in}{0.624448in}}{\pgfqpoint{10.792231in}{7.176941in}} %
\pgfusepath{clip}%
\pgfsetbuttcap%
\pgfsetmiterjoin%
\definecolor{currentfill}{rgb}{0.149020,0.545098,0.823529}%
\pgfsetfillcolor{currentfill}%
\pgfsetfillopacity{0.800000}%
\pgfsetlinewidth{0.000000pt}%
\definecolor{currentstroke}{rgb}{0.000000,0.000000,0.000000}%
\pgfsetstrokecolor{currentstroke}%
\pgfsetstrokeopacity{0.800000}%
\pgfsetdash{}{0pt}%
\pgfpathmoveto{\pgfqpoint{7.945148in}{4.982662in}}%
\pgfpathlineto{\pgfqpoint{8.579985in}{4.982662in}}%
\pgfpathlineto{\pgfqpoint{8.579985in}{5.028480in}}%
\pgfpathlineto{\pgfqpoint{7.945148in}{5.028480in}}%
\pgfpathclose%
\pgfusepath{fill}%
\end{pgfscope}%
\begin{pgfscope}%
\pgfpathrectangle{\pgfqpoint{4.009158in}{0.624448in}}{\pgfqpoint{10.792231in}{7.176941in}} %
\pgfusepath{clip}%
\pgfsetbuttcap%
\pgfsetmiterjoin%
\definecolor{currentfill}{rgb}{0.149020,0.545098,0.823529}%
\pgfsetfillcolor{currentfill}%
\pgfsetfillopacity{0.800000}%
\pgfsetlinewidth{0.000000pt}%
\definecolor{currentstroke}{rgb}{0.000000,0.000000,0.000000}%
\pgfsetstrokecolor{currentstroke}%
\pgfsetstrokeopacity{0.800000}%
\pgfsetdash{}{0pt}%
\pgfpathmoveto{\pgfqpoint{8.579985in}{4.799390in}}%
\pgfpathlineto{\pgfqpoint{9.659208in}{4.799390in}}%
\pgfpathlineto{\pgfqpoint{9.659208in}{4.845208in}}%
\pgfpathlineto{\pgfqpoint{8.579985in}{4.845208in}}%
\pgfpathclose%
\pgfusepath{fill}%
\end{pgfscope}%
\begin{pgfscope}%
\pgfpathrectangle{\pgfqpoint{4.009158in}{0.624448in}}{\pgfqpoint{10.792231in}{7.176941in}} %
\pgfusepath{clip}%
\pgfsetbuttcap%
\pgfsetmiterjoin%
\definecolor{currentfill}{rgb}{0.149020,0.545098,0.823529}%
\pgfsetfillcolor{currentfill}%
\pgfsetfillopacity{0.800000}%
\pgfsetlinewidth{0.000000pt}%
\definecolor{currentstroke}{rgb}{0.000000,0.000000,0.000000}%
\pgfsetstrokecolor{currentstroke}%
\pgfsetstrokeopacity{0.800000}%
\pgfsetdash{}{0pt}%
\pgfpathmoveto{\pgfqpoint{14.103068in}{4.616117in}}%
\pgfpathlineto{\pgfqpoint{14.420487in}{4.616117in}}%
\pgfpathlineto{\pgfqpoint{14.420487in}{4.661935in}}%
\pgfpathlineto{\pgfqpoint{14.103068in}{4.661935in}}%
\pgfpathclose%
\pgfusepath{fill}%
\end{pgfscope}%
\begin{pgfscope}%
\pgfpathrectangle{\pgfqpoint{4.009158in}{0.624448in}}{\pgfqpoint{10.792231in}{7.176941in}} %
\pgfusepath{clip}%
\pgfsetbuttcap%
\pgfsetmiterjoin%
\definecolor{currentfill}{rgb}{0.149020,0.545098,0.823529}%
\pgfsetfillcolor{currentfill}%
\pgfsetfillopacity{0.800000}%
\pgfsetlinewidth{0.000000pt}%
\definecolor{currentstroke}{rgb}{0.000000,0.000000,0.000000}%
\pgfsetstrokecolor{currentstroke}%
\pgfsetstrokeopacity{0.800000}%
\pgfsetdash{}{0pt}%
\pgfpathmoveto{\pgfqpoint{5.342316in}{4.432845in}}%
\pgfpathlineto{\pgfqpoint{5.405800in}{4.432845in}}%
\pgfpathlineto{\pgfqpoint{5.405800in}{4.478663in}}%
\pgfpathlineto{\pgfqpoint{5.342316in}{4.478663in}}%
\pgfpathclose%
\pgfusepath{fill}%
\end{pgfscope}%
\begin{pgfscope}%
\pgfpathrectangle{\pgfqpoint{4.009158in}{0.624448in}}{\pgfqpoint{10.792231in}{7.176941in}} %
\pgfusepath{clip}%
\pgfsetbuttcap%
\pgfsetmiterjoin%
\definecolor{currentfill}{rgb}{0.149020,0.545098,0.823529}%
\pgfsetfillcolor{currentfill}%
\pgfsetfillopacity{0.800000}%
\pgfsetlinewidth{0.000000pt}%
\definecolor{currentstroke}{rgb}{0.000000,0.000000,0.000000}%
\pgfsetstrokecolor{currentstroke}%
\pgfsetstrokeopacity{0.800000}%
\pgfsetdash{}{0pt}%
\pgfpathmoveto{\pgfqpoint{5.405800in}{4.249573in}}%
\pgfpathlineto{\pgfqpoint{5.723218in}{4.249573in}}%
\pgfpathlineto{\pgfqpoint{5.723218in}{4.295391in}}%
\pgfpathlineto{\pgfqpoint{5.405800in}{4.295391in}}%
\pgfpathclose%
\pgfusepath{fill}%
\end{pgfscope}%
\begin{pgfscope}%
\pgfpathrectangle{\pgfqpoint{4.009158in}{0.624448in}}{\pgfqpoint{10.792231in}{7.176941in}} %
\pgfusepath{clip}%
\pgfsetbuttcap%
\pgfsetmiterjoin%
\definecolor{currentfill}{rgb}{0.149020,0.545098,0.823529}%
\pgfsetfillcolor{currentfill}%
\pgfsetfillopacity{0.800000}%
\pgfsetlinewidth{0.000000pt}%
\definecolor{currentstroke}{rgb}{0.000000,0.000000,0.000000}%
\pgfsetstrokecolor{currentstroke}%
\pgfsetstrokeopacity{0.800000}%
\pgfsetdash{}{0pt}%
\pgfpathmoveto{\pgfqpoint{6.738958in}{4.066301in}}%
\pgfpathlineto{\pgfqpoint{8.960887in}{4.066301in}}%
\pgfpathlineto{\pgfqpoint{8.960887in}{4.112119in}}%
\pgfpathlineto{\pgfqpoint{6.738958in}{4.112119in}}%
\pgfpathclose%
\pgfusepath{fill}%
\end{pgfscope}%
\begin{pgfscope}%
\pgfpathrectangle{\pgfqpoint{4.009158in}{0.624448in}}{\pgfqpoint{10.792231in}{7.176941in}} %
\pgfusepath{clip}%
\pgfsetbuttcap%
\pgfsetmiterjoin%
\definecolor{currentfill}{rgb}{0.149020,0.545098,0.823529}%
\pgfsetfillcolor{currentfill}%
\pgfsetfillopacity{0.800000}%
\pgfsetlinewidth{0.000000pt}%
\definecolor{currentstroke}{rgb}{0.000000,0.000000,0.000000}%
\pgfsetstrokecolor{currentstroke}%
\pgfsetstrokeopacity{0.800000}%
\pgfsetdash{}{0pt}%
\pgfpathmoveto{\pgfqpoint{8.579985in}{3.883028in}}%
\pgfpathlineto{\pgfqpoint{13.849133in}{3.883028in}}%
\pgfpathlineto{\pgfqpoint{13.849133in}{3.928846in}}%
\pgfpathlineto{\pgfqpoint{8.579985in}{3.928846in}}%
\pgfpathclose%
\pgfusepath{fill}%
\end{pgfscope}%
\begin{pgfscope}%
\pgfpathrectangle{\pgfqpoint{4.009158in}{0.624448in}}{\pgfqpoint{10.792231in}{7.176941in}} %
\pgfusepath{clip}%
\pgfsetbuttcap%
\pgfsetmiterjoin%
\definecolor{currentfill}{rgb}{0.149020,0.545098,0.823529}%
\pgfsetfillcolor{currentfill}%
\pgfsetfillopacity{0.800000}%
\pgfsetlinewidth{0.000000pt}%
\definecolor{currentstroke}{rgb}{0.000000,0.000000,0.000000}%
\pgfsetstrokecolor{currentstroke}%
\pgfsetstrokeopacity{0.800000}%
\pgfsetdash{}{0pt}%
\pgfpathmoveto{\pgfqpoint{13.658682in}{3.699756in}}%
\pgfpathlineto{\pgfqpoint{13.912617in}{3.699756in}}%
\pgfpathlineto{\pgfqpoint{13.912617in}{3.745574in}}%
\pgfpathlineto{\pgfqpoint{13.658682in}{3.745574in}}%
\pgfpathclose%
\pgfusepath{fill}%
\end{pgfscope}%
\begin{pgfscope}%
\pgfpathrectangle{\pgfqpoint{4.009158in}{0.624448in}}{\pgfqpoint{10.792231in}{7.176941in}} %
\pgfusepath{clip}%
\pgfsetbuttcap%
\pgfsetmiterjoin%
\definecolor{currentfill}{rgb}{0.149020,0.545098,0.823529}%
\pgfsetfillcolor{currentfill}%
\pgfsetfillopacity{0.800000}%
\pgfsetlinewidth{0.000000pt}%
\definecolor{currentstroke}{rgb}{0.000000,0.000000,0.000000}%
\pgfsetstrokecolor{currentstroke}%
\pgfsetstrokeopacity{0.800000}%
\pgfsetdash{}{0pt}%
\pgfpathmoveto{\pgfqpoint{14.483970in}{3.516484in}}%
\pgfpathlineto{\pgfqpoint{14.547454in}{3.516484in}}%
\pgfpathlineto{\pgfqpoint{14.547454in}{3.562302in}}%
\pgfpathlineto{\pgfqpoint{14.483970in}{3.562302in}}%
\pgfpathclose%
\pgfusepath{fill}%
\end{pgfscope}%
\begin{pgfscope}%
\pgfpathrectangle{\pgfqpoint{4.009158in}{0.624448in}}{\pgfqpoint{10.792231in}{7.176941in}} %
\pgfusepath{clip}%
\pgfsetbuttcap%
\pgfsetmiterjoin%
\definecolor{currentfill}{rgb}{0.149020,0.545098,0.823529}%
\pgfsetfillcolor{currentfill}%
\pgfsetfillopacity{0.800000}%
\pgfsetlinewidth{0.000000pt}%
\definecolor{currentstroke}{rgb}{0.000000,0.000000,0.000000}%
\pgfsetstrokecolor{currentstroke}%
\pgfsetstrokeopacity{0.800000}%
\pgfsetdash{}{0pt}%
\pgfpathmoveto{\pgfqpoint{14.166552in}{3.333212in}}%
\pgfpathlineto{\pgfqpoint{14.737905in}{3.333212in}}%
\pgfpathlineto{\pgfqpoint{14.737905in}{3.379030in}}%
\pgfpathlineto{\pgfqpoint{14.166552in}{3.379030in}}%
\pgfpathclose%
\pgfusepath{fill}%
\end{pgfscope}%
\begin{pgfscope}%
\pgfpathrectangle{\pgfqpoint{4.009158in}{0.624448in}}{\pgfqpoint{10.792231in}{7.176941in}} %
\pgfusepath{clip}%
\pgfsetbuttcap%
\pgfsetmiterjoin%
\definecolor{currentfill}{rgb}{0.149020,0.545098,0.823529}%
\pgfsetfillcolor{currentfill}%
\pgfsetfillopacity{0.800000}%
\pgfsetlinewidth{0.000000pt}%
\definecolor{currentstroke}{rgb}{0.000000,0.000000,0.000000}%
\pgfsetstrokecolor{currentstroke}%
\pgfsetstrokeopacity{0.800000}%
\pgfsetdash{}{0pt}%
\pgfpathmoveto{\pgfqpoint{14.420487in}{3.149939in}}%
\pgfpathlineto{\pgfqpoint{14.737905in}{3.149939in}}%
\pgfpathlineto{\pgfqpoint{14.737905in}{3.195757in}}%
\pgfpathlineto{\pgfqpoint{14.420487in}{3.195757in}}%
\pgfpathclose%
\pgfusepath{fill}%
\end{pgfscope}%
\begin{pgfscope}%
\pgfpathrectangle{\pgfqpoint{4.009158in}{0.624448in}}{\pgfqpoint{10.792231in}{7.176941in}} %
\pgfusepath{clip}%
\pgfsetbuttcap%
\pgfsetmiterjoin%
\definecolor{currentfill}{rgb}{0.149020,0.545098,0.823529}%
\pgfsetfillcolor{currentfill}%
\pgfsetfillopacity{0.800000}%
\pgfsetlinewidth{0.000000pt}%
\definecolor{currentstroke}{rgb}{0.000000,0.000000,0.000000}%
\pgfsetstrokecolor{currentstroke}%
\pgfsetstrokeopacity{0.800000}%
\pgfsetdash{}{0pt}%
\pgfpathmoveto{\pgfqpoint{4.009158in}{2.966667in}}%
\pgfpathlineto{\pgfqpoint{4.897930in}{2.966667in}}%
\pgfpathlineto{\pgfqpoint{4.897930in}{3.012485in}}%
\pgfpathlineto{\pgfqpoint{4.009158in}{3.012485in}}%
\pgfpathclose%
\pgfusepath{fill}%
\end{pgfscope}%
\begin{pgfscope}%
\pgfpathrectangle{\pgfqpoint{4.009158in}{0.624448in}}{\pgfqpoint{10.792231in}{7.176941in}} %
\pgfusepath{clip}%
\pgfsetbuttcap%
\pgfsetmiterjoin%
\definecolor{currentfill}{rgb}{0.149020,0.545098,0.823529}%
\pgfsetfillcolor{currentfill}%
\pgfsetfillopacity{0.800000}%
\pgfsetlinewidth{0.000000pt}%
\definecolor{currentstroke}{rgb}{0.000000,0.000000,0.000000}%
\pgfsetstrokecolor{currentstroke}%
\pgfsetstrokeopacity{0.800000}%
\pgfsetdash{}{0pt}%
\pgfpathmoveto{\pgfqpoint{7.500762in}{2.783395in}}%
\pgfpathlineto{\pgfqpoint{7.945148in}{2.783395in}}%
\pgfpathlineto{\pgfqpoint{7.945148in}{2.829213in}}%
\pgfpathlineto{\pgfqpoint{7.500762in}{2.829213in}}%
\pgfpathclose%
\pgfusepath{fill}%
\end{pgfscope}%
\begin{pgfscope}%
\pgfpathrectangle{\pgfqpoint{4.009158in}{0.624448in}}{\pgfqpoint{10.792231in}{7.176941in}} %
\pgfusepath{clip}%
\pgfsetbuttcap%
\pgfsetmiterjoin%
\definecolor{currentfill}{rgb}{0.149020,0.545098,0.823529}%
\pgfsetfillcolor{currentfill}%
\pgfsetfillopacity{0.800000}%
\pgfsetlinewidth{0.000000pt}%
\definecolor{currentstroke}{rgb}{0.000000,0.000000,0.000000}%
\pgfsetstrokecolor{currentstroke}%
\pgfsetstrokeopacity{0.800000}%
\pgfsetdash{}{0pt}%
\pgfpathmoveto{\pgfqpoint{7.500762in}{2.600123in}}%
\pgfpathlineto{\pgfqpoint{7.564246in}{2.600123in}}%
\pgfpathlineto{\pgfqpoint{7.564246in}{2.645941in}}%
\pgfpathlineto{\pgfqpoint{7.500762in}{2.645941in}}%
\pgfpathclose%
\pgfusepath{fill}%
\end{pgfscope}%
\begin{pgfscope}%
\pgfpathrectangle{\pgfqpoint{4.009158in}{0.624448in}}{\pgfqpoint{10.792231in}{7.176941in}} %
\pgfusepath{clip}%
\pgfsetbuttcap%
\pgfsetmiterjoin%
\definecolor{currentfill}{rgb}{0.149020,0.545098,0.823529}%
\pgfsetfillcolor{currentfill}%
\pgfsetfillopacity{0.800000}%
\pgfsetlinewidth{0.000000pt}%
\definecolor{currentstroke}{rgb}{0.000000,0.000000,0.000000}%
\pgfsetstrokecolor{currentstroke}%
\pgfsetstrokeopacity{0.800000}%
\pgfsetdash{}{0pt}%
\pgfpathmoveto{\pgfqpoint{7.500762in}{2.416850in}}%
\pgfpathlineto{\pgfqpoint{7.564246in}{2.416850in}}%
\pgfpathlineto{\pgfqpoint{7.564246in}{2.462669in}}%
\pgfpathlineto{\pgfqpoint{7.500762in}{2.462669in}}%
\pgfpathclose%
\pgfusepath{fill}%
\end{pgfscope}%
\begin{pgfscope}%
\pgfpathrectangle{\pgfqpoint{4.009158in}{0.624448in}}{\pgfqpoint{10.792231in}{7.176941in}} %
\pgfusepath{clip}%
\pgfsetbuttcap%
\pgfsetmiterjoin%
\definecolor{currentfill}{rgb}{0.149020,0.545098,0.823529}%
\pgfsetfillcolor{currentfill}%
\pgfsetfillopacity{0.800000}%
\pgfsetlinewidth{0.000000pt}%
\definecolor{currentstroke}{rgb}{0.000000,0.000000,0.000000}%
\pgfsetstrokecolor{currentstroke}%
\pgfsetstrokeopacity{0.800000}%
\pgfsetdash{}{0pt}%
\pgfpathmoveto{\pgfqpoint{7.500762in}{2.233578in}}%
\pgfpathlineto{\pgfqpoint{7.564246in}{2.233578in}}%
\pgfpathlineto{\pgfqpoint{7.564246in}{2.279396in}}%
\pgfpathlineto{\pgfqpoint{7.500762in}{2.279396in}}%
\pgfpathclose%
\pgfusepath{fill}%
\end{pgfscope}%
\begin{pgfscope}%
\pgfpathrectangle{\pgfqpoint{4.009158in}{0.624448in}}{\pgfqpoint{10.792231in}{7.176941in}} %
\pgfusepath{clip}%
\pgfsetbuttcap%
\pgfsetmiterjoin%
\definecolor{currentfill}{rgb}{0.149020,0.545098,0.823529}%
\pgfsetfillcolor{currentfill}%
\pgfsetfillopacity{0.800000}%
\pgfsetlinewidth{0.000000pt}%
\definecolor{currentstroke}{rgb}{0.000000,0.000000,0.000000}%
\pgfsetstrokecolor{currentstroke}%
\pgfsetstrokeopacity{0.800000}%
\pgfsetdash{}{0pt}%
\pgfpathmoveto{\pgfqpoint{7.691213in}{2.050306in}}%
\pgfpathlineto{\pgfqpoint{8.262567in}{2.050306in}}%
\pgfpathlineto{\pgfqpoint{8.262567in}{2.096124in}}%
\pgfpathlineto{\pgfqpoint{7.691213in}{2.096124in}}%
\pgfpathclose%
\pgfusepath{fill}%
\end{pgfscope}%
\begin{pgfscope}%
\pgfpathrectangle{\pgfqpoint{4.009158in}{0.624448in}}{\pgfqpoint{10.792231in}{7.176941in}} %
\pgfusepath{clip}%
\pgfsetbuttcap%
\pgfsetmiterjoin%
\definecolor{currentfill}{rgb}{0.149020,0.545098,0.823529}%
\pgfsetfillcolor{currentfill}%
\pgfsetfillopacity{0.800000}%
\pgfsetlinewidth{0.000000pt}%
\definecolor{currentstroke}{rgb}{0.000000,0.000000,0.000000}%
\pgfsetstrokecolor{currentstroke}%
\pgfsetstrokeopacity{0.800000}%
\pgfsetdash{}{0pt}%
\pgfpathmoveto{\pgfqpoint{8.262567in}{1.867034in}}%
\pgfpathlineto{\pgfqpoint{9.151339in}{1.867034in}}%
\pgfpathlineto{\pgfqpoint{9.151339in}{1.912852in}}%
\pgfpathlineto{\pgfqpoint{8.262567in}{1.912852in}}%
\pgfpathclose%
\pgfusepath{fill}%
\end{pgfscope}%
\begin{pgfscope}%
\pgfpathrectangle{\pgfqpoint{4.009158in}{0.624448in}}{\pgfqpoint{10.792231in}{7.176941in}} %
\pgfusepath{clip}%
\pgfsetbuttcap%
\pgfsetmiterjoin%
\definecolor{currentfill}{rgb}{0.149020,0.545098,0.823529}%
\pgfsetfillcolor{currentfill}%
\pgfsetfillopacity{0.800000}%
\pgfsetlinewidth{0.000000pt}%
\definecolor{currentstroke}{rgb}{0.000000,0.000000,0.000000}%
\pgfsetstrokecolor{currentstroke}%
\pgfsetstrokeopacity{0.800000}%
\pgfsetdash{}{0pt}%
\pgfpathmoveto{\pgfqpoint{8.262567in}{1.683762in}}%
\pgfpathlineto{\pgfqpoint{8.326050in}{1.683762in}}%
\pgfpathlineto{\pgfqpoint{8.326050in}{1.729580in}}%
\pgfpathlineto{\pgfqpoint{8.262567in}{1.729580in}}%
\pgfpathclose%
\pgfusepath{fill}%
\end{pgfscope}%
\begin{pgfscope}%
\pgfpathrectangle{\pgfqpoint{4.009158in}{0.624448in}}{\pgfqpoint{10.792231in}{7.176941in}} %
\pgfusepath{clip}%
\pgfsetbuttcap%
\pgfsetmiterjoin%
\definecolor{currentfill}{rgb}{0.149020,0.545098,0.823529}%
\pgfsetfillcolor{currentfill}%
\pgfsetfillopacity{0.800000}%
\pgfsetlinewidth{0.000000pt}%
\definecolor{currentstroke}{rgb}{0.000000,0.000000,0.000000}%
\pgfsetstrokecolor{currentstroke}%
\pgfsetstrokeopacity{0.800000}%
\pgfsetdash{}{0pt}%
\pgfpathmoveto{\pgfqpoint{8.706953in}{1.500489in}}%
\pgfpathlineto{\pgfqpoint{8.770436in}{1.500489in}}%
\pgfpathlineto{\pgfqpoint{8.770436in}{1.546307in}}%
\pgfpathlineto{\pgfqpoint{8.706953in}{1.546307in}}%
\pgfpathclose%
\pgfusepath{fill}%
\end{pgfscope}%
\begin{pgfscope}%
\pgfpathrectangle{\pgfqpoint{4.009158in}{0.624448in}}{\pgfqpoint{10.792231in}{7.176941in}} %
\pgfusepath{clip}%
\pgfsetbuttcap%
\pgfsetmiterjoin%
\definecolor{currentfill}{rgb}{0.149020,0.545098,0.823529}%
\pgfsetfillcolor{currentfill}%
\pgfsetfillopacity{0.800000}%
\pgfsetlinewidth{0.000000pt}%
\definecolor{currentstroke}{rgb}{0.000000,0.000000,0.000000}%
\pgfsetstrokecolor{currentstroke}%
\pgfsetstrokeopacity{0.800000}%
\pgfsetdash{}{0pt}%
\pgfpathmoveto{\pgfqpoint{9.151339in}{1.317217in}}%
\pgfpathlineto{\pgfqpoint{14.674421in}{1.317217in}}%
\pgfpathlineto{\pgfqpoint{14.674421in}{1.363035in}}%
\pgfpathlineto{\pgfqpoint{9.151339in}{1.363035in}}%
\pgfpathclose%
\pgfusepath{fill}%
\end{pgfscope}%
\begin{pgfscope}%
\pgfpathrectangle{\pgfqpoint{4.009158in}{0.624448in}}{\pgfqpoint{10.792231in}{7.176941in}} %
\pgfusepath{clip}%
\pgfsetbuttcap%
\pgfsetmiterjoin%
\definecolor{currentfill}{rgb}{0.149020,0.545098,0.823529}%
\pgfsetfillcolor{currentfill}%
\pgfsetfillopacity{0.800000}%
\pgfsetlinewidth{0.000000pt}%
\definecolor{currentstroke}{rgb}{0.000000,0.000000,0.000000}%
\pgfsetstrokecolor{currentstroke}%
\pgfsetstrokeopacity{0.800000}%
\pgfsetdash{}{0pt}%
\pgfpathmoveto{\pgfqpoint{13.023845in}{1.133945in}}%
\pgfpathlineto{\pgfqpoint{14.230035in}{1.133945in}}%
\pgfpathlineto{\pgfqpoint{14.230035in}{1.179763in}}%
\pgfpathlineto{\pgfqpoint{13.023845in}{1.179763in}}%
\pgfpathclose%
\pgfusepath{fill}%
\end{pgfscope}%
\begin{pgfscope}%
\pgfpathrectangle{\pgfqpoint{4.009158in}{0.624448in}}{\pgfqpoint{10.792231in}{7.176941in}} %
\pgfusepath{clip}%
\pgfsetbuttcap%
\pgfsetmiterjoin%
\definecolor{currentfill}{rgb}{0.149020,0.545098,0.823529}%
\pgfsetfillcolor{currentfill}%
\pgfsetfillopacity{0.800000}%
\pgfsetlinewidth{0.000000pt}%
\definecolor{currentstroke}{rgb}{0.000000,0.000000,0.000000}%
\pgfsetstrokecolor{currentstroke}%
\pgfsetstrokeopacity{0.800000}%
\pgfsetdash{}{0pt}%
\pgfpathmoveto{\pgfqpoint{4.009158in}{0.950673in}}%
\pgfpathlineto{\pgfqpoint{14.737905in}{0.950673in}}%
\pgfpathlineto{\pgfqpoint{14.737905in}{0.996491in}}%
\pgfpathlineto{\pgfqpoint{4.009158in}{0.996491in}}%
\pgfpathclose%
\pgfusepath{fill}%
\end{pgfscope}%
\begin{pgfscope}%
\pgfsetrectcap%
\pgfsetmiterjoin%
\pgfsetlinewidth{1.003750pt}%
\definecolor{currentstroke}{rgb}{1.000000,1.000000,1.000000}%
\pgfsetstrokecolor{currentstroke}%
\pgfsetdash{}{0pt}%
\pgfpathmoveto{\pgfqpoint{4.009158in}{0.624448in}}%
\pgfpathlineto{\pgfqpoint{4.009158in}{7.801389in}}%
\pgfusepath{stroke}%
\end{pgfscope}%
\begin{pgfscope}%
\pgfsetrectcap%
\pgfsetmiterjoin%
\pgfsetlinewidth{1.003750pt}%
\definecolor{currentstroke}{rgb}{1.000000,1.000000,1.000000}%
\pgfsetstrokecolor{currentstroke}%
\pgfsetdash{}{0pt}%
\pgfpathmoveto{\pgfqpoint{14.801389in}{0.624448in}}%
\pgfpathlineto{\pgfqpoint{14.801389in}{7.801389in}}%
\pgfusepath{stroke}%
\end{pgfscope}%
\begin{pgfscope}%
\pgfsetrectcap%
\pgfsetmiterjoin%
\pgfsetlinewidth{1.003750pt}%
\definecolor{currentstroke}{rgb}{1.000000,1.000000,1.000000}%
\pgfsetstrokecolor{currentstroke}%
\pgfsetdash{}{0pt}%
\pgfpathmoveto{\pgfqpoint{4.009158in}{0.624448in}}%
\pgfpathlineto{\pgfqpoint{14.801389in}{0.624448in}}%
\pgfusepath{stroke}%
\end{pgfscope}%
\begin{pgfscope}%
\pgfsetrectcap%
\pgfsetmiterjoin%
\pgfsetlinewidth{1.003750pt}%
\definecolor{currentstroke}{rgb}{1.000000,1.000000,1.000000}%
\pgfsetstrokecolor{currentstroke}%
\pgfsetdash{}{0pt}%
\pgfpathmoveto{\pgfqpoint{4.009158in}{7.801389in}}%
\pgfpathlineto{\pgfqpoint{14.801389in}{7.801389in}}%
\pgfusepath{stroke}%
\end{pgfscope}%
\begin{pgfscope}%
\pgfsetbuttcap%
\pgfsetmiterjoin%
\definecolor{currentfill}{rgb}{0.898039,0.898039,0.898039}%
\pgfsetfillcolor{currentfill}%
\pgfsetfillopacity{0.800000}%
\pgfsetlinewidth{0.501875pt}%
\definecolor{currentstroke}{rgb}{0.800000,0.800000,0.800000}%
\pgfsetstrokecolor{currentstroke}%
\pgfsetstrokeopacity{0.800000}%
\pgfsetdash{}{0pt}%
\pgfpathmoveto{\pgfqpoint{13.696707in}{7.283648in}}%
\pgfpathlineto{\pgfqpoint{14.704167in}{7.283648in}}%
\pgfpathquadraticcurveto{\pgfqpoint{14.731944in}{7.283648in}}{\pgfqpoint{14.731944in}{7.311426in}}%
\pgfpathlineto{\pgfqpoint{14.731944in}{7.704167in}}%
\pgfpathquadraticcurveto{\pgfqpoint{14.731944in}{7.731944in}}{\pgfqpoint{14.704167in}{7.731944in}}%
\pgfpathlineto{\pgfqpoint{13.696707in}{7.731944in}}%
\pgfpathquadraticcurveto{\pgfqpoint{13.668929in}{7.731944in}}{\pgfqpoint{13.668929in}{7.704167in}}%
\pgfpathlineto{\pgfqpoint{13.668929in}{7.311426in}}%
\pgfpathquadraticcurveto{\pgfqpoint{13.668929in}{7.283648in}}{\pgfqpoint{13.696707in}{7.283648in}}%
\pgfpathclose%
\pgfusepath{stroke,fill}%
\end{pgfscope}%
\begin{pgfscope}%
\pgfsetbuttcap%
\pgfsetmiterjoin%
\definecolor{currentfill}{rgb}{0.796078,0.294118,0.086275}%
\pgfsetfillcolor{currentfill}%
\pgfsetfillopacity{0.800000}%
\pgfsetlinewidth{0.000000pt}%
\definecolor{currentstroke}{rgb}{0.000000,0.000000,0.000000}%
\pgfsetstrokecolor{currentstroke}%
\pgfsetstrokeopacity{0.800000}%
\pgfsetdash{}{0pt}%
\pgfpathmoveto{\pgfqpoint{13.724485in}{7.569306in}}%
\pgfpathlineto{\pgfqpoint{14.002262in}{7.569306in}}%
\pgfpathlineto{\pgfqpoint{14.002262in}{7.666528in}}%
\pgfpathlineto{\pgfqpoint{13.724485in}{7.666528in}}%
\pgfpathclose%
\pgfusepath{fill}%
\end{pgfscope}%
\begin{pgfscope}%
\pgftext[x=14.113373in,y=7.569306in,left,base]{\sffamily\fontsize{10.000000}{12.000000}\selectfont Planned}%
\end{pgfscope}%
\begin{pgfscope}%
\pgfsetbuttcap%
\pgfsetmiterjoin%
\definecolor{currentfill}{rgb}{0.149020,0.545098,0.823529}%
\pgfsetfillcolor{currentfill}%
\pgfsetfillopacity{0.800000}%
\pgfsetlinewidth{0.000000pt}%
\definecolor{currentstroke}{rgb}{0.000000,0.000000,0.000000}%
\pgfsetstrokecolor{currentstroke}%
\pgfsetstrokeopacity{0.800000}%
\pgfsetdash{}{0pt}%
\pgfpathmoveto{\pgfqpoint{13.724485in}{7.365991in}}%
\pgfpathlineto{\pgfqpoint{14.002262in}{7.365991in}}%
\pgfpathlineto{\pgfqpoint{14.002262in}{7.463214in}}%
\pgfpathlineto{\pgfqpoint{13.724485in}{7.463214in}}%
\pgfpathclose%
\pgfusepath{fill}%
\end{pgfscope}%
\begin{pgfscope}%
\pgftext[x=14.113373in,y=7.365991in,left,base]{\sffamily\fontsize{10.000000}{12.000000}\selectfont Effective}%
\end{pgfscope}%
\end{pgfpicture}%
\makeatother%
\endgroup%

        \caption{The timetables of the project. The planned and effective time spans are visible.}
    \end{figure}
\end{a3pages}

\chapter{Work Packages (revised)} % <<< ------------------------------------------------ %
\label{ch:wpac}
% ---------------------------------------------------------------------------- %

        %wpac:
        %{Start Date}
        %{End Date}
        %{Hours}
        %{Ancestors}
        %{Descendants}
        %{Assignee}
        %{Description}

\section{General Prerequisite Tasks}
\label{sec:general}

Contains tasks to get an overview  of the previous project's current state and
to define  further steps and  requirements to  this project which  follow from
that.

\subsection{Read Literature about the Existing Project}
\label{subsec:general:read}

\wpac
    {}
    {}
    {}
    {None}
    {}
    {Raphael Frey, Noah Hüsser}
    {%
        Read  the preceding  project's documentation  and general  information
        about  the STEMlab  board to  get familiar  with the  project and  its
        general state. Based  on this,  the specifications will  be determined
        and  the  project requirements  (this  document)  will be  revised  to
        reflect the new findings.%
    }

\subsection{Tutorial to IIR Filters}
\label{subsec:general:iir}

\wpac
    {}
    {}
    {}
    {None}
    {}
    {Raphael Frey, Noah Hüsser}
    {%
        Complete  the  lab course  about  IIR  filters  and their  bit  growth
        problems and fixes provided by Mr. Pichler.
    }

\subsection{Assess Results of the Predecessor Project}
\label{subsec:general:assess}

\wpac
    {}
    {}
    {}
    {Read Literature about the Existing Project}
    {}
    {Raphael Frey, Noah Hüsser}
    {%
        Define what the current project status  is and where we can connect to
        the predecessor project.
    }

\subsection{Outline the Project}
\label{subsec:general:outline}

\wpac
    {}
    {}
    {}
    {\ref{subsec:general:assess}}
    {}
    {Raphael Frey, Noah Hüsser}
    {%
        Define what  the requirements for the  project are and which  of those
        are a must and which a nice to have.
    }

\section{Filter Design}
\label{sec:filters}

Contains the  work packages  relating to  filter design  in Matlab  and filter
implementation in the FPGA tool chain.

\subsection{Filter Research}
\label{subsec:filter:research}

\wpac
    {}
    {}
    {}
    {\ref{subsec:general:outline}}
    {}
    {Raphael Frey}
    {%
        Research the digital filter technologies which will be needed for this
        project. In particular, this includes FIR and CIC filters, since Xilinx
        provide predefined blocks for these filters in their toolchain.

        Besides the general knowledge on digital filters, the Xilinx toolchain
        must be researched in order to understand the FIR and CIC filter blocks
        and allow their usage.
    }

\subsection{Determine Filter Specifications}
\label{subsec:filter:specification}

\wpac
    {}
    {}
    {}
    {\ref{subsec:filter:research}, \ref{subsec:general:outline}}
    {}
    {Raphael Frey}
    {%
        Specify the minimum and nice to have requirements for filters.
        Document this well.
    }

\subsection{Matlab Scripts for Filter Design}
\label{subsec:filter:matlab}

\wpac
    {}
    {}
    {}
    {\ref{subsec:filter:specification}}
    {}
    {Raphael Frey}
    {%
        Matlab scripts  which design  various filter chains. These  consist of
        one  or  several  dispatcher  scripts  where  the  filter  chains  are
        specified in  terms of  their frequency  band properties,  and scripts
        which actually design the filters according to these specifications.

        Specifically,   scripts  for   generating   CIC   filters  and   their
        compensation  filters, as  well  as FIR  filters  are needed. For  FIR
        filters, the resulting  coefficients are to be saved in  files so that
        they can be loaded by the Xilinx tool chain.

        If possible, Matlab's parallelism should be exploited to reduce filter
        design times.
    }

\subsection{TCL Scripts for Filter Evaluation}
\label{subsec:filter:tcl}

\wpac
    {}
    {}
    {}
    {\ref{subsec:filter:matlab}}
    {}
    {Raphael Frey}
    {%
    In  order to  assess the  resource usage  of the  filter chains  which are
    designed by  Matlab, the respective  filter blocks  (CIC, FIR) need  to be
    implemented in Vivado and a bitstream must be compiled.

    Since       the      number       of       filters      generated       by
    package~\ref{subsec:filter:matlab}  can  be  very large  (dozens  or  even
    hundreds of filters), this process must be automated to be of any use. For
    this, Tcl scripts are used.

    For FIR  filters, the scripts load  the coefficient files which  have been
    generated by Matlab.  Vivado will generate usage reports which can then be
    automatically post-processed.
    }


\subsection{Document Filter Design and Implementation}
\label{subsec:filter:doc}

\wpac
    {}
    {}
    {}
    {\ref{subsec:filter:tcl}}
    {}
    {Raphael Frey}
    {%
        Documentation  for the  filter  design  tool chain. The  documentation
        should be sufficiently detailed so that  the tool chain can be used by
        a person who wishes to create new filters for a new bitstream.
    }


\section{Documentation}
\label{sec:docs}

\subsection{Disposition}
\wpac
    {}
    {}
    {}
    {None}
    {}
    {Noah H\"usser}
    {%
    Create a general outline of the thesis document.%
    }

\subsection{Report}
\wpac
    {}
    {}
    {}
    {None}
    {}
    {Raphael Frey, Noah H\"usser}
    {%
    Continuously  check, restructure  and  update the  report  to current  new
    insights.  Compile and round out the  actual thesis document after all the
    different parts have been written.%
    }


\section{Firmware}
\label{sec:firmware}

\subsection{Linux}
\label{subsec:fw:linux}

\wpac
    {}
    {}
    {}
    {None}
    {}
    {Noah H\"usser}
    {%
    Boot an Ubuntu Linux onto the Red Pitaya.%
    The target here  is to have a  ARM Linux distribution which  boots fine on
    the STEMlab.  The official sources  from the STEMlab project should be
    used as they contain a lot of patches for the STEMlab hardware.
    A standalone toolchain has to be developed as the compiled versions do not
    contain the right device tree (DT) and board support package (BSP).
    Compiling in the right DT and BSP is not part of this task.
    }

\subsection{Server Application}
\label{subsec:fw:server}

\subsubsection{Server Design Choices}
\label{subsubsec:fw:server:design-decisions}
\wpac
    {}
    {}
    {}
    {None}
    {}
    {Noah H\"usser}
    {%
        It has to be determined what  technology should be used to provide the
        server side of the data transmission.   The server needs to be able to
        perform IOCTL  calls and to serve  data with the chosen  protocol over
        the network. It has to run on the arm GNU/Linux and shouldn't have too
        many  dependencies.  A  C-like language  will most  likely serve  best
        since it's close to the kernel.
    }

\subsubsection{Build External Libraries and Test Them}
\label{subsubsec:fw:server:external}
\wpac
    {}
    {}
    {}
    {\ref{subsubsec:fw:server:design-decisions}}
    {}
    {Noah H\"usser}
    {%
        A  library  for  networking  and  another  for  JSON  are  needed. The
        requirements  will  be decided  by  the  specifications of  the  scope
        application.
    }

\subsubsection{Write the Server Application}
\label{subsubsec:fw:server:server}
\wpac
    {}
    {}
    {}
    {\ref{subsubsec:fw:server:external}}
    {}
    {Noah H\"usser}
    {%
        A server application has to be written that reads data from the Logger
        and transmits it  over the network.  It needs to  be able to configure
        the logger  via RPCs  that are mapped  to \texttt{IOCTL}s.   It should
        read data from  \texttt{/dev/logger} which is presented  by the kernel
        module.
    }

\subsubsection{Document the Server Application}
\label{subsubsec:fw:server:docs}
\wpac
    {}
    {}
    {}
    {\ref{subsubsec:fw:server:external}}
    {}
    {Raphael Frey, Noah H\"usser}
    {%
        Write documentation for the server application build process.
    }

\subsection{FPGA}
\label{subsec:fw:fpga}

\subsubsection{Minimal Working Example}
\label{subsubsec:fw:fpga:mwe}
\wpac
    {}
    {}
    {}
    {None}
    {}
    {Noah H\"usser}
    {%
        Get a minimal working example project running on the Red Pitaya.  This
        requires importing  the cores  from Pavel Demin's  project as  well as
        compiling the Linux for the STEMlab.  Correct functionality of the Red
        Pitaya hardware,  particularly the  ADC, is to  be verified  via Anton
        Potochnik's frequency counter.
    }

\subsubsection{Set up a Sane Toolchain for the FPGA}
\label{subsubsec:fw:fpga:sane}

\wpac
    {}
    {}
    {}
    {\ref{subsec:general:assess},\ref{subsubsec:fw:fpga:mwe}}
    {}
    {Raphael Frey}
    {%
    Depending on the previous project's  results, an entirely new Toolchain is
    required or  not. Either way, Vivado  can be very  finicky and we  need to
    have a  sane toolchain for  the entire FPGA build  process so we  can work
    cleanly.
    }

\subsubsection{Port the ZYNQ Logger}
\label{subsubsec:fw:fpga:logger}
\wpac
    {}
    {}
    {}
    {\ref{subsubsec:fw:fpga:sane}}
    {}
    {Noah H\"usser}
    {%
        Port the Zynq Logger to the Red Pitaya. The main challenges will be
        the porting of its interface (for which a block \code{axis2datalanes}
        is developed), as well as understanding and properly implementing the
        device tree and the kernel module which is needed for the Linux OS
        to interface with the logger hardware.
    }

\subsubsection{Implement the Filter Chains}
\label{subsubsec:fw:fpga:filters}
\wpac
    {}
    {}
    {}
    {\ref{subsubsec:fw:fpga:logger}}
    {}
    {Raphael Frey}
    {%
        Implement the filter  chains as designed in Matlab on  the FPGA.  This
        requires correct functionality of the FIR and CIC filter blocks.
    }

\subsubsection{Verification by Simulation}
\label{subsubsec:fw:fpga:sim}
\wpac
    {}
    {}
    {}
    {\ref{subsubsec:fw:fpga:filters}}
    {}
    {Raphael Frey}
    {%
        Simulate all  the FPGA blocks.   Simulate the  filters with a  sine to
        make sure the filter work as expected.
    }

\subsubsection{Document FPGA Toolchain and Filter Structure}
\label{subsubsec:fw:fpga:doc}
\wpac
    {}
    {}
    {}
    {\ref{subsubsec:fw:fpga:sim}}
    {}
    {Raphael Frey}
    {%
        Document the implemented filter chains and surrounding FPGA structure.
    }

\section{Validation}
\label{sec:validation}

\subsection{Validate the Build Process}
\label{subsec:validation:build}
\wpac
    {}
    {}
    {}
    {\ref{subsubsec:fw:fpga:sim}, \ref{subsubsec:fw:server:server}, \ref{subsubsec:subsec:fw:linux}}
    {}
    {Raphael Frey}
    {%
        Build the entire project (Linux, bitstream), flash result onto the Red
        Pitaya and verify correct functionality.
    }

\subsection{Validate the Filter Implementation}
\label{subsec:validation:impl}
\wpac
    {}
    {}
    {}
    {\ref{subsubsec:fw:fpga:sim}}
    {}
    {Noah Hüsser}
    {%
        Validate  the specified  frequency  response of  the  filters and  the
        attenuation in the passband.
    }

\section{Front-end}
\label{sec:frontend}

\subsection{Technology Research}
\label{subsec:frontend:research}
\wpac
    {}
    {}
    {}
    {None}
    {}
    {Noah Hüsser}
    {%
        Research the possibilities to design and implement a nice front-end for
        the STEMlab recording system.
        The  STEMlab Project  has a  built-in web  interface which  allows the
        display of signals and rudimentary calculations.
        Furthermore  Prof. Gut has  implemented  a spectrum  analyzer in  Java
        which interfaces with the STEMlab.
        Both approaches are feasible and should be evaluated thoroughly.
        Furthermore it has to be evaluated  how the data should be transferred
        to the front-end.
        This decision  correlates a  lot with the  selection of  a programming
        language since e.g. web technologies do not support UDP.
        Finally, existing libraries have to examined very carefully. There are
        a lot  of libraries around data  transmission and GUI design,  but not
        all of them are viable solutions.
    }

\subsection{Scope Specifications}
\label{subsec:frontend:specify}
\wpac
    {}
    {}
    {}
    {\ref{subsec:frontend:research}}
    {}
    {Noah Hüsser}
    {%
        Specify the scope's requirements and nice-to-haves.
    }

\subsection{Implementation of a Basic Plotter}
\label{subsec:frontend:plotter}
\wpac
    {}
    {}
    {}
    {\ref{subsec:frontend:research}}
    {}
    {Noah Hüsser}
    {%
        At first a  simple plotter is implemented. It should  not do more
        than plot a given vector of data points on a canvas.
        The plotting  should be done  with high  performance at 60  frames per
        second independently of the size of the data vector.
        Of course more data will slow down the plotting, but it should perform
        well enough that vectors up to 64 kilosamples in size plot nicely.
        For this task  it is recommended to use existing  libraries which make
        use of the GPU such as OpenGL.
    }

\subsection{Implementation of a Basic Receiver}
\label{subsec:frontend:receiver}
\wpac
    {}
    {}
    {}
    {\ref{subsec:frontend:research}\ref{subsubsec:fw:server:server}}
    {}
    {Noah Hüsser}
    {%
        A simple receiver which receives recorded samples over the network should be created.
        For  this task  a simple  emitter which  emits test  samples without
        actually recording anything should be implemented.
        In further  tasks it can  be used conveniently  without the need  of a
        running STEMlab board.
        The receiver should simply receive  samples as packets with lengths of
        powers of two over the network.
        Again, a  good, high-performance and publicly  accepted library should
        be used to reduce workload and ensure reliability.
    }

\subsection{Combining the Plotter and the Receiver}
\label{subsec:frontend:combining}
\wpac
    {}
    {}
    {}
    {\ref{subsec:frontend:plotter}, \ref{subsec:frontend:receiver}}
    {}
    {Noah Hüsser}
    {%
        The plotter  and the  receiver should be  combined to  display vectors
        received over  the network.  This  should happen fluently  without any
        delays.
    }

\subsection{Setting a Trigger on the Logger Core}
\label{subsec:frontend:trig}
\wpac
    {}
    {}
    {}
    {\ref{subsec:frontend:combining}}
    {}
    {Noah Hüsser}
    {%
        The application  should have a ability  to set a trigger  type for
        any channel with the possibility to do this using a GUI.
    }

\subsection{Implement Proper Scaling}
\label{subsec:frontend:scaling}
\wpac
    {}
    {}
    {}
    {\ref{subsec:frontend:trig}}
    {}
    {Noah Hüsser}
    {%
        The GUI should  not only plot a qualitatively correct  signal but also
        show its  quantitative properties.  A  grid should be  programmed into
        the GUI which can be used to  read metrics of a signal. A panel on the
        GUI should  display the current grid  size.  It should be  possible to
        scale signals using the mouse or the GUI. Grids should adjust with the
        current scale.
    }

\subsection{Implement Trigger Modes}
\label{subsec:frontend:modes}
\wpac
    {}
    {}
    {}
    {\ref{subsec:frontend:trig}}
    {}
    {Noah Hüsser}
    {%
        It should be possible to use the scope application in \textit{normal},
        \textit{auto}, and \textit{single} mode. Those modes should behave as follows:
        \begin{tabular}{l l}
        \textit{normal} & send a packet when a trigger was fired and repeats this procedure \\
        \textit{auto} & send a packet when a trigger was fired or a timeout occurred and repeat \\
        \textit{single} & send a packet after a first trigger and does not repeat this procedure
        \end{tabular}
    }

\subsection{Implement an FFT}
\label{subsec:frontend:fft}
\wpac
    {}
    {}
    {}
    {\ref{subsec:frontend:specify}}
    {}
    {Noah Hüsser}
    {%
        An FFT should be calculated and displayed live.
    }

\subsection{Implement Power Measurements}
\label{subsec:frontend:}
\wpac
    {}
    {}
    {}
    {\ref{subsec:frontend:specify}}
    {}
    {Noah Hüsser}
    {%
        Calculate the power in the signal.
        It has  to be possible to  calculate the power between  two markers to
        limit the calculation on a certain part of the signal.
    }

\subsection{Implement SNR Calculation}
\label{subsec:frontend:}
\wpac
    {}
    {}
    {}
    {\ref{subsec:frontend:specify}}
    {}
    {Noah Hüsser}
    {%
        The SNR of the signal between two markers has to be calculated automatically.
        As a reference use the Java application of Prof. Gut.
    }

\subsection{Implement THD Calculation}
\label{subsec:frontend:}
\wpac
    {}
    {}
    {}
    {\ref{subsec:frontend:specify}}
    {}
    {Noah Hüsser}
    {%
        The THD of the signal spectrum has to be calculated automatically.
        As a reference, use the Java application of Prof. Gut.
    }

\subsection{Implement Dual Channel Possibilities}
\label{subsec:frontend:}
\wpac
    {}
    {}
    {}
    {\ref{subsec:frontend:specify}}
    {}
    {Noah Hüsser}
    {%
        The  entire transmission  and plotting  of all  the signals  should be
        implemented for dual channel capabilities.
        The FPGA code supports up to  eight channels whilst only supplying two
        to  the server  since the  STEMlab  board only  supports two  physical
        channels.
    }

\subsection{Document the Scope}
\label{subsec:frontend:}
\wpac
    {}
    {}
    {}
    {\ref{subsec:frontend:specify}}
    {}
    {Noah Hüsser}
    {%
        Document the entire oscilloscope application and how to use it.
    }

% >>>

\backmatter
\end{document}
%^^A vim: foldenable foldcolumn=4 foldmethod=marker foldmarker=<<<,>>>
